% Motion of Particles in two Dimensions — Further Mathematics Upper Sixth — Cours signé OnBuch+
% Official MINESEC syllabus. Compiler avec : tectonic FMA11-motion-of-particles-in-two-dimensions.tex
\newcommand{\DOCMATIERE}{Further Mathematics}
\newcommand{\DOCNIVEAU}{Upper Sixth}
\newcommand{\DOCMODULE}{Upper Sixth — Further mathematics}
\newcommand{\DOCLECON}{11}
\newcommand{\DOCTITRE}{Motion of Particles in two Dimensions}
% OnBuch+ — préambule « cours signés » ENRICHI (compilé avec tectonic / XeLaTeX)
% Système visuel « Pop » d'OnBuch+ : fond crème (jamais blanc), bordures encre
% épaisses, ombres « dures » décalées, titres Archivo Black en capitales.
% Version MATHÉMATIQUES : graphes (pgfplots/tikz), boîtes pédagogiques
% (définition, propriété, méthode, attention, le savais-tu, exercice résolu,
% exercice, TP, bilan), page de garde et signature OnBuch+ ancrée en bas.
%
% Macros attendues dans main.tex AVANT \input{preamble} :
%   \DOCMATIERE  \DOCNIVEAU  \DOCMODULE  \DOCLECON (numéro)  \DOCTITRE
\documentclass[11pt]{article}

\usepackage{fontspec}
\usepackage[english]{babel}
\usepackage[a4paper,margin=2cm,top=3.4cm,bottom=2.8cm,headheight=1.4cm,footskip=1.3cm]{geometry}
\usepackage{amsmath,amssymb,amsfonts}
\usepackage{mathtools} % \xmapsto, \coloneqq... souvent utilisés par les modèles
\usepackage{cancel} % simplifications barrées (bilans de forces)
% \abs{z} (module, valeur absolue) et \norm{u} (norme), utilisés par les modèles
\providecommand{\abs}{}\let\abs\relax\DeclarePairedDelimiter{\abs}{\lvert}{\rvert}
\providecommand{\norm}{}\let\norm\relax\DeclarePairedDelimiter{\norm}{\lVert}{\rVert}
\usepackage[table]{xcolor} % [table] : fournit aussi \rowcolors (alternance de lignes)
\usepackage{pagecolor}
\usepackage{tikz}
\usetikzlibrary{arrows.meta,positioning,calc,shapes.geometric,shapes.misc,shapes.arrows,decorations.pathmorphing,decorations.pathreplacing,decorations.markings,patterns,shadows,mindmap,trees,fit,backgrounds,matrix,intersections,angles,quotes}
\usepackage{pgfplots}
\usepackage[version=4]{mhchem} % équations nucléaires \ce{^{235}_{92}U}
\usepackage[european,siunitx]{circuitikz} % schémas électriques
\pgfplotsset{compat=1.18}
\usepgfplotslibrary{fillbetween,groupplots} % aires entre courbes (calcul intégral)
\usepackage{enumitem}
\usepackage{fancyhdr}
% [most] charge toutes les bibliothèques tcolorbox (listings, minted, raster,
% documentation...) et gonfle inutilement la mémoire TeX sur un document long
% (~300 boîtes) : on ne charge que ce qui est réellement utilisé.
\usepackage[skins,breakable]{tcolorbox}
\usepackage{graphicx}
\usepackage{colortbl}
\usepackage{booktabs}
\usepackage{tabularx}
\usepackage{diagbox} % case d'en-tête diagonale des tableaux à double entrée
% Colonnes extensibles centrée / gauche / droite (C, L, R), souvent utilisées par les modèles
\newcolumntype{C}{>{\centering\arraybackslash}X}
\newcolumntype{L}{>{\raggedright\arraybackslash}X}
\newcolumntype{R}{>{\raggedleft\arraybackslash}X}
\usepackage{array}
\usepackage{multicol}
\usepackage{siunitx}
% parse-numbers=false : accepte aussi \SI{2,5 \times 10^{-3}}{...}
\sisetup{locale=UK,output-decimal-marker={.},group-minimum-digits=4,parse-numbers=false}
\DeclareSIUnit\ohm{\ensuremath{\symup{\Omega}}} % U+2126 absent de la police
\DeclareSIUnit\division{div} % réglages d'oscilloscope (ms/div, V/div)
\DeclareSIUnit\tr{tr} % tours (vitesse de rotation, ex. \SI{1500}{\tr\per\minute})
\usepackage{qrcode}
% \degree : commande fréquemment utilisée par les modèles pour un angle ou une
% température en dehors de \SI{}{} (non fournie par les paquets chargés).
\providecommand{\degree}{^{\circ}}
% \qed (fin de démonstration), utilisé par les modèles sans amsthm
\providecommand{\qed}{\hfill$\square$}
% \barre : double barre des valeurs interdites dans un tableau de variations (tkz-tab)
\providecommand{\barre}{\ensuremath{\|}}
% environnement proof (amsthm non chargé) utilisé par les modèles
\ifdefined\proof\else\newenvironment{proof}[1][Proof]{\par\noindent\textit{#1.}\ }{\qed\par}\fi
\usepackage{lastpage}
\usepackage{hyperref}
\hypersetup{hidelinks}

% ============================================================
% Palette « Pop » OnBuch+ — mêmes hex que la classe P (plus_ui.dart)
% ============================================================
\definecolor{poporange}{HTML}{F59321}
\definecolor{popdark}{HTML}{E07A0C}
\definecolor{popcream}{HTML}{FFF6E8}
\definecolor{popink}{HTML}{1C1714}
\definecolor{poppurple}{HTML}{7A5AE0}
\definecolor{popgreen}{HTML}{2E9E6A}
\definecolor{poppink}{HTML}{E0457B}
\definecolor{popblue}{HTML}{2F7DE1}
\definecolor{popgold}{HTML}{C98A12}
\definecolor{popmuted}{HTML}{8A7F76}
\definecolor{popline}{HTML}{EADFD2}
\definecolor{popgray}{HTML}{8A7F76}
\colorlet{popgrey}{popgray}
% Alias tolérants (le modèle nomme parfois les couleurs différemment)
\colorlet{popwhite}{white}
\colorlet{popblack}{popink}
\colorlet{popred}{poppink}
\colorlet{popyellow}{popgold}
% Teintes claires pour fonds de tableaux / zones
\colorlet{popblueL}{popblue!10!white}
\colorlet{poporangeL}{poporange!14!white}
\colorlet{popgreenL}{popgreen!12!white}
\colorlet{poppurpleL}{poppurple!12!white}
\colorlet{poppinkL}{poppink!12!white}
\colorlet{popgoldL}{popgold!14!white}

\pagecolor{popcream}

% ============================================================
% Typographie
% ============================================================
\defaultfontfeatures{Ligatures=TeX}
\setmainfont{PlusJakartaSans-Medium.ttf}[Path=../../../fonts/,BoldFont=PlusJakartaSans-Bold.ttf]
% Maths en sans-serif (Fira Math) avec chiffres et lettres droites de Plus
% Jakarta, pour que les formules se fondent dans le texte.
\usepackage[math-style=ISO,bold-style=ISO]{unicode-math}
\setmathfont{FiraMath-Regular.otf}[Path=../../../fonts/]
\setmathfont{PlusJakartaSans-Medium.ttf}[Path=../../../fonts/,range={up/num,up/{Latin,latin},\mathup}]
% (icomma retiré : décimales avec point en anglais)
% Caractères mathématiques Unicode absents des polices de texte (Plus Jakarta,
% Archivo Black) : rendus par leur équivalent mathématique, en texte comme en
% formule — sinon ils s'affichent en carré vide (ex. « Arithmétique dans ℤ »).
\usepackage{newunicodechar}
\newunicodechar{ℝ}{\ensuremath{\mathbb{R}}}
\newunicodechar{ℤ}{\ensuremath{\mathbb{Z}}}
\newunicodechar{ℂ}{\ensuremath{\mathbb{C}}}
\newunicodechar{ℕ}{\ensuremath{\mathbb{N}}}
\newunicodechar{ℚ}{\ensuremath{\mathbb{Q}}}
\newunicodechar{𝔻}{\ensuremath{\mathbb{D}}}
\newunicodechar{α}{\ensuremath{\alpha}}
\newunicodechar{β}{\ensuremath{\beta}}
\newunicodechar{γ}{\ensuremath{\gamma}}
\newunicodechar{δ}{\ensuremath{\delta}}
\newunicodechar{ε}{\ensuremath{\varepsilon}}
\newunicodechar{θ}{\ensuremath{\theta}}
\newunicodechar{λ}{\ensuremath{\lambda}}
\newunicodechar{μ}{\ensuremath{\mu}}
\newunicodechar{σ}{\ensuremath{\sigma}}
\newunicodechar{τ}{\ensuremath{\tau}}
\newunicodechar{φ}{\ensuremath{\varphi}}
\newunicodechar{ψ}{\ensuremath{\psi}}
\newunicodechar{ω}{\ensuremath{\omega}}
\newunicodechar{Σ}{\ensuremath{\Sigma}}
% majuscules produites par \MakeUppercase dans les titres (α-aminés -> Α)
\newunicodechar{Α}{\ensuremath{\alpha}}
\newunicodechar{Β}{\ensuremath{\beta}}
\newunicodechar{Γ}{\ensuremath{\Gamma}}
\newunicodechar{Δ}{\ensuremath{\Delta}}
\newunicodechar{Ε}{\ensuremath{\varepsilon}}
\newunicodechar{Θ}{\ensuremath{\Theta}}
\newunicodechar{Λ}{\ensuremath{\Lambda}}
\newunicodechar{Μ}{\ensuremath{\mu}}
\newunicodechar{Τ}{\ensuremath{\tau}}
\newunicodechar{Φ}{\ensuremath{\Phi}}
\newunicodechar{Ψ}{\ensuremath{\Psi}}
\newunicodechar{Ω}{\ensuremath{\Omega}}
\newunicodechar{↦}{\ensuremath{\mapsto}}
\newunicodechar{∈}{\ensuremath{\in}}
\newunicodechar{∉}{\ensuremath{\notin}}
\newunicodechar{∧}{\ensuremath{\wedge}}
\newunicodechar{⇔}{\ensuremath{\Leftrightarrow}}
\newunicodechar{⟺}{\ensuremath{\Longleftrightarrow}}
\newunicodechar{⇒}{\ensuremath{\Rightarrow}}
\newunicodechar{⟹}{\ensuremath{\Longrightarrow}}
\newunicodechar{∘}{\ensuremath{\circ}}
\newunicodechar{∪}{\ensuremath{\cup}}
\newunicodechar{∩}{\ensuremath{\cap}}
\newunicodechar{≡}{\ensuremath{\equiv}}
\newunicodechar{⊂}{\ensuremath{\subset}}
\newunicodechar{⊥}{\ensuremath{\perp}}
\newunicodechar{∃}{\ensuremath{\exists}}
\newunicodechar{∀}{\ensuremath{\forall}}
\newunicodechar{∛}{\ensuremath{\sqrt[3]{\,}}}
\newunicodechar{∜}{\ensuremath{\sqrt[4]{\,}}}
\newunicodechar{⃗}{\ensuremath{{}^{\to}}}
\newunicodechar{ⁿ}{\ifmmode^{n}\else\textsuperscript{\ensuremath{n}}\fi}
\newunicodechar{ˣ}{\ifmmode^{x}\else\textsuperscript{\ensuremath{x}}\fi}
\newunicodechar{ᵇ}{\ifmmode^{b}\else\textsuperscript{\ensuremath{b}}\fi}
\newunicodechar{ʳ}{\ifmmode^{r}\else\textsuperscript{\ensuremath{r}}\fi}
\newunicodechar{ᵘ}{\ifmmode^{u}\else\textsuperscript{\ensuremath{u}}\fi}
\newunicodechar{ʸ}{\ifmmode^{y}\else\textsuperscript{\ensuremath{y}}\fi}
\newunicodechar{ᵃ}{\ifmmode^{a}\else\textsuperscript{\ensuremath{a}}\fi}
\newunicodechar{ᵉ}{\ifmmode^{e}\else\textsuperscript{\ensuremath{e}}\fi}
\newunicodechar{ᵏ}{\ifmmode^{k}\else\textsuperscript{\ensuremath{k}}\fi}
\newunicodechar{ᵖ}{\ifmmode^{p}\else\textsuperscript{\ensuremath{p}}\fi}
\newunicodechar{ᴺ}{\ifmmode^{N}\else\textsuperscript{\ensuremath{N}}\fi}
\newunicodechar{ᶜ}{\ifmmode^{c}\else\textsuperscript{\ensuremath{c}}\fi}
\newunicodechar{ʰ}{\ifmmode^{h}\else\textsuperscript{\ensuremath{h}}\fi}
\newunicodechar{ᵠ}{\ifmmode^{q}\else\textsuperscript{\ensuremath{q}}\fi}
\newunicodechar{ₙ}{\ifmmode_{n}\else\textsubscript{\ensuremath{n}}\fi}
\newunicodechar{ₐ}{\ifmmode_{a}\else\textsubscript{\ensuremath{a}}\fi}
\newunicodechar{ᵢ}{\ifmmode_{i}\else\textsubscript{\ensuremath{i}}\fi}
\newunicodechar{ᵣ}{\ifmmode_{r}\else\textsubscript{\ensuremath{r}}\fi}
\newunicodechar{ₖ}{\ifmmode_{k}\else\textsubscript{\ensuremath{k}}\fi}
\newunicodechar{ᵦ}{\ifmmode_{\beta}\else\textsubscript{\ensuremath{\beta}}\fi}
\newunicodechar{ₓ}{\ifmmode_{x}\else\textsubscript{\ensuremath{x}}\fi}
\newfontfamily\popdisplay{ArchivoBlack-Regular.ttf}[Path=../../../fonts/]
\newfontfamily\poplogo{SpaceGrotesk-Bold.ttf}[Path=../../../fonts/]
\newfontfamily\popsemi{PlusJakartaSans-SemiBold.ttf}[Path=../../../fonts/]
\newfontfamily\popxbold{PlusJakartaSans-ExtraBold.ttf}[Path=../../../fonts/]
\newfontfamily\popxboldls{PlusJakartaSans-ExtraBold.ttf}[Path=../../../fonts/,LetterSpace=3.0]

\setlist[enumerate]{leftmargin=1.7em,itemsep=5pt,topsep=4pt}
\setlist[itemize]{leftmargin=1.7em,itemsep=4pt,topsep=4pt,label={\color{poporange}\textbullet}}
\setlength{\parindent}{0pt}
\setlength{\parskip}{5pt}
\renewcommand{\arraystretch}{1.25}

% pgfplots : style Pop par défaut
\pgfplotsset{
  every axis/.append style={axis line style={popink,line width=1pt},tick style={popink},
    grid style={popline,line width=0.5pt},label style={font=\small},tick label style={font=\footnotesize}},
  popcurve/.style={poporange,line width=1.8pt,no markers,smooth},
}
% Même style disponible dans un \draw TikZ simple (hors pgfplots)
\tikzset{popcurve/.style={poporange,line width=1.8pt}}
\tikzset{popgrid/.style={popline,very thin}}
\colorlet{popcurve}{poporange} % le nom du style est parfois utilisé comme couleur

% ============================================================
% Pill / squiggle
% ============================================================
\newcommand{\poppill}[3]{%
  \tikz[baseline=-0.55ex]\node[draw=popink,line width=1.2pt,rounded corners=2.6pt,%
    fill=#2,inner xsep=6.5pt,inner ysep=3pt]{%
    \popxboldls\fontsize{7}{7}\selectfont\color{#3}\MakeUppercase{#1}};%
}
\newcommand{\popsquiggle}[1]{%
  \tikz[baseline=-0.4ex]{
    \draw[#1,line width=2.6pt,line cap=round]
      plot [smooth] coordinates {(0,0.09) (0.34,0.01) (0.68,0.21) (1.02,0.01) (1.36,0.10)};
  }%
}
% Mot-clé surligné (marqueur orange sous le texte)
\newcommand{\cle}[1]{{\popxbold\color{popdark}#1}}

% ============================================================
% En-tête / pied
% ============================================================
\pagestyle{fancy}
\fancyhf{}
\renewcommand{\headrulewidth}{0pt}
\renewcommand{\footrulewidth}{0pt}
\lhead{\raisebox{-0.15ex}{\poplogo\fontsize{13}{13}\selectfont\color{popink}OnBuch\color{poporange}+}}
\rhead{\poppill{\DOCMATIERE\ \textbullet\ \DOCNIVEAU\ \textbullet\ Chapter \DOCLECON}{white}{popink}}
\lfoot{\popxbold\fontsize{7.6}{7.6}\selectfont\color{popmuted}\MakeUppercase{Signed course OnBuch+}\ \textbullet\ {\popsemi\DOCTITRE}}
\rfoot{\tikz[baseline=-0.6ex]\node[rounded corners=8.5pt,fill=popink,minimum height=17pt,minimum width=17pt,inner xsep=5pt,inner sep=0]{\popxbold\fontsize{8}{8}\selectfont\color{popcream}\thepage\,/\,\pageref*{LastPage}};}

% ============================================================
% Titres
% ============================================================
\newcommand{\chaptitre}[2]{%
  \begin{tcolorbox}[enhanced,colback=poporange,colframe=popink,boxrule=2.6pt,arc=11pt,
      drop shadow=popink,left=15pt,right=15pt,top=12pt,bottom=11pt,before skip=3pt,after skip=18pt]
    \poppill{#2}{popink}{popcream}\par\vspace{9pt}
    {\raggedright\hyphenpenalty=10000\exhyphenpenalty=10000\popdisplay\fontsize{22}{24}\selectfont\color{popcream}\MakeUppercase{#1}\par}
  \end{tcolorbox}
}
% Section numérotée : pastille orange avec le numéro + titre Archivo
\newcounter{popsec}
\newcommand{\coursec}[1]{%
  \stepcounter{popsec}\setcounter{popsub}{0}%
  \par\vspace{16pt}\noindent
  \begin{minipage}{\linewidth}
  \tikz[baseline=-0.7ex]\node[draw=popink,line width=1.6pt,fill=poporange,rounded corners=4pt,minimum size=22pt,inner sep=0]{\popdisplay\fontsize{12}{12}\selectfont\color{popcream}\thepopsec};%
  \hspace{8pt}{\popdisplay\fontsize{15}{17}\selectfont\color{popink}\MakeUppercase{#1}}\par
  \vspace{4pt}\noindent{\color{poporange}\rule{36pt}{3.6pt}}
  \end{minipage}\par\nopagebreak\vspace{8pt}
}
\newcounter{popsub}
\newcommand{\courssub}[1]{%
  \stepcounter{popsub}%
  \par\vspace{9pt}\noindent{\popxbold\fontsize{11.5}{13}\selectfont\color{popink}{\color{poporange}\thepopsec.\thepopsub}\hspace{6pt}#1}\par\nopagebreak\vspace{4pt}
}

% ============================================================
% Boîtes pédagogiques — même recette : fond blanc, bordure épaisse colorée,
% ombre dure encre, badge plein. Titre optionnel [#1].
% ============================================================
\tcbset{popbase/.style={breakable,enhanced,colback=white,boxrule=2.3pt,arc=9pt,
  shadow={3pt}{-3pt}{0pt}{popink},left=14pt,right=14pt,top=11pt,bottom=11pt,before skip=11pt,after skip=15pt,
  fontupper=\popsemi\fontsize{10.6}{14.5}\selectfont\color{popink},
  fontlower=\popsemi\fontsize{10.6}{14.5}\selectfont\color{popink}}}
% Coche dessinée (le glyphe ✓ n'existe pas dans les polices du document)
\newcommand{\popcheck}[1]{\tikz[baseline=-0.5ex]\draw[#1,line width=1.6pt,line cap=round,line join=round](-0.13,0.0)--(-0.04,-0.09)--(0.13,0.1);}
\providecommand{\coche}{\popcheck{popgreen}}
\providecommand{\resultat}[1]{\par\smallskip\noindent\fcolorbox{popgreen}{popgreenL}{\parbox{\dimexpr\linewidth-2\fboxsep-2\fboxrule\relax}{\textbf{Result.} #1}}\par\smallskip}
\newcommand{\popboxhead}[3]{\poppill{#1}{#2}{white}\ifx\relax#3\relax\else\hspace{6pt}{\popxbold\fontsize{10.6}{12}\selectfont\color{popink}#3}\fi\par\vspace{7pt}}

\newtcolorbox{aretenir}[1][]{popbase,colframe=popgreen,before upper={\popboxhead{Key points}{popgreen}{#1}}}
\newtcolorbox{exemplebox}[1][]{popbase,colframe=poporange,before upper={\popboxhead{Exemple}{poporange}{#1}}}
\newtcolorbox{definition}[1][]{popbase,colframe=popblue,before upper={\popboxhead{Definition}{popblue}{#1}}}
\newtcolorbox{propriete}[1][]{popbase,colframe=poppurple,before upper={\popboxhead{Property}{poppurple}{#1}}}
\newtcolorbox{methode}[1][]{popbase,colframe=popgold,before upper={\popboxhead{Method}{popgold}{#1}}}
\newtcolorbox{attention}[1][]{popbase,colframe=poppink,before upper={\popboxhead{Warning}{poppink}{#1}}}
\newtcolorbox{experience}[1][]{popbase,colframe=popblue,colback=popblueL,before upper={\popboxhead{Experiment}{popblue}{#1}}}
% « Le savais-tu ? » : carte violette pleine (contexte, culture, Cameroun)
\newtcolorbox{savaistu}[1][]{popbase,colback=poppurple,colframe=popink,
  fontupper=\popsemi\fontsize{10.4}{14.2}\selectfont\color{white},
  before upper={\poppill{Did you know?}{white}{poppurple}\ifx\relax#1\relax\else\hspace{6pt}{\popxbold\color{white}#1}\fi\par\vspace{7pt}}}
% Exercice résolu : énoncé en haut, \tcblower puis la solution sur fond crème
\newcounter{popexo}\newcounter{popexr}
\newtcolorbox{exoresolu}[1][]{popbase,colframe=poporange,colbacklower=popcream,
  segmentation style={popink,line width=1.2pt,dashed},
  before upper={\stepcounter{popexr}\popboxhead{Worked example \thepopexr}{poporange}{#1}},
  before lower={\poppill{Solution}{popink}{popcream}\par\vspace{6pt}}}
% Exercice (non résolu en place) — la correction est dans la partie Corrigés
\newtcolorbox{exercice}[1][]{popbase,colframe=popink,
  before upper={\stepcounter{popexo}\popboxhead{Exercise \thepopexo}{popink}{#1}}}
% Corrigé
\newtcolorbox{corrige}[1][]{popbase,colframe=popgreen,colback=popgreenL,
  before upper={\popboxhead{Solution}{popgreen}{#1}}}
% Objectifs / prérequis / bilan
\newtcolorbox{objectifs}{popbase,colframe=popink,colback=poporangeL,
  before upper={\popboxhead{Chapter objectives}{popink}{}}}
\newtcolorbox{prerequis}{popbase,colframe=popink,colback=popblueL,
  before upper={\popboxhead{Prerequisites}{popblue}{}}}
\newtcolorbox{bilan}{popbase,colframe=popink,colback=white,boxrule=2.8pt,
  before upper={\popboxhead{Fiche bilan}{poporange}{L'essentiel à connaître pour l'examen}}}
% Figure encadrée avec légende
\newtcolorbox{popfigure}[1][]{enhanced,breakable=false,colback=white,colframe=popline,boxrule=1.4pt,arc=8pt,
  left=10pt,right=10pt,top=10pt,bottom=8pt,before skip=10pt,after skip=12pt,halign=center,
  lower separated=false,halign lower=center,
  fontlower=\popsemi\fontsize{9}{11.5}\selectfont\color{popmuted},
  #1}
\newcommand{\legende}[1]{\par\vspace{6pt}{\popsemi\fontsize{9}{11.5}\selectfont\color{popmuted}#1}}

% ============================================================
% Signature OnBuch+
% ============================================================
\newcommand{\poplogoinline}[1]{{\poplogo\fontsize{#1}{#1}\selectfont\color{popink}OnBuch\color{poporange}+}}
\newcommand{\signatureonbuch}{%
  \begin{tcolorbox}[enhanced,colback=popink,colframe=popink,boxrule=2.2pt,arc=10pt,
      drop shadow=poporange,left=14pt,right=14pt,top=10pt,bottom=10pt,before skip=0pt,after skip=0pt]
    \tikz[baseline=-0.7ex]\node[circle,fill=popgreen,draw=popcream,line width=1.3pt,minimum size=22pt,inner sep=0]{\popcheck{white}};%
    \hspace{9pt}\begin{minipage}[c]{0.62\linewidth}
      {\popxbold\fontsize{12}{13}\selectfont\color{popcream}Signed\ \textbullet\ {\poplogo OnBuch\color{poporange}+}}\par\vspace{2pt}
      {\raggedright\popsemi\fontsize{8}{10.5}\selectfont\color{popline}\MakeUppercase{OnBuch+ teaching team}\\\MakeUppercase{Follows the official MINESEC syllabus}\par}
    \end{minipage}\hfill
    {\poplogo\fontsize{22}{22}\selectfont\color{popcream}OnBuch\color{poporange}+}
  \end{tcolorbox}%
}

% ============================================================
% Page de garde
% #1 titre · #2 matière · #3 niveau · #4 module · #5 n° leçon · #6 durée ·
% #7 description / objectifs (texte ou itemize)
% ============================================================
\newcommand{\pagegarde}[7]{%
  \thispagestyle{empty}
  \newgeometry{a4paper,margin=1.7cm,top=1.6cm,bottom=1.4cm}%
  \begin{tikzpicture}[remember picture,overlay]
    \fill[poporange!18] ([xshift=-3.2cm,yshift=-3.6cm]current page.north east) circle (4.6cm);
    \fill[poppurple!14] ([xshift=2.4cm,yshift=7.6cm]current page.south west) circle (3.8cm);
  \end{tikzpicture}%
  {\poplogo\fontsize{30}{30}\selectfont\color{popink}OnBuch\color{poporange}+}\hfill
  \raisebox{0.6ex}{\poppill{Signed course \textbullet\ Premium}{popink}{popcream}}\par
  \vspace{1pt}\popsquiggle{poppurple}\par
  \vspace{8pt}
  \begin{tcolorbox}[enhanced,colback=poporange,colframe=popink,boxrule=2.8pt,arc=13pt,
      drop shadow=popink,left=18pt,right=18pt,top=12pt,bottom=13pt,after skip=10pt]
    \poppill{#2 \textbullet\ #3}{popink}{popcream}\hspace{5pt}\poppill{#4}{white}{popink}\par\vspace{8pt}
    {\popdisplay\fontsize{14}{14}\selectfont\color{popink}CHAPTER #5}\par\vspace{4pt}
    {\raggedright\hyphenpenalty=10000\exhyphenpenalty=10000\popdisplay\fontsize{25}{27}\selectfont\color{popcream}\MakeUppercase{#1}\par}
    \vspace{8pt}
    {\popxbold\fontsize{9.5}{9.5}\selectfont\color{popink}\MakeUppercase{Suggested time: #6}}
  \end{tcolorbox}
  \begin{tcolorbox}[enhanced,colback=white,colframe=popink,boxrule=2.2pt,arc=10pt,
      drop shadow=popink,left=16pt,right=16pt,top=10pt,bottom=10pt,after skip=8pt]
    \poppill{In this chapter}{poppurple}{white}\par\vspace{5pt}
    \popsemi\fontsize{9.3}{12.6}\selectfont\color{popink}#7
  \end{tcolorbox}
  \noindent\begin{minipage}[t]{0.487\linewidth}
    \begin{tcolorbox}[enhanced,colback=popgreen,colframe=popink,boxrule=2.2pt,arc=10pt,
        drop shadow=popink,left=11pt,right=11pt,top=10pt,bottom=10pt,height=3.1cm,valign=center]
      \tcbox[colback=white,colframe=popink,boxrule=1.3pt,arc=5pt,boxsep=0pt,nobeforeafter,
        left=3pt,right=3pt,top=3pt,bottom=3pt,valign=center]{\qrcode[height=1.9cm]{https://play.google.com/store/apps/details?id=com.luvvix.onbuch&hl=en}}%
      \hspace{8pt}\begin{minipage}[c]{0.5\linewidth}
        {\popdisplay\fontsize{11}{12.5}\selectfont\color{white}\MakeUppercase{Download OnBuch}}\par\vspace{4pt}
        {\raggedright\popsemi\fontsize{8.4}{11}\selectfont\color{white}Free \textbullet\ Google Play\\AI tutor, past papers, results\par}
      \end{minipage}
    \end{tcolorbox}
  \end{minipage}\hfill
  \begin{minipage}[t]{0.487\linewidth}
    \begin{tcolorbox}[enhanced,colback=poppurple,colframe=popink,boxrule=2.2pt,arc=10pt,
        drop shadow=popink,left=11pt,right=11pt,top=10pt,bottom=10pt,height=3.1cm,valign=center]
      \tikz[baseline=-0.7ex]\node[circle,fill=white,draw=popink,line width=1.3pt,minimum size=1.6cm,inner sep=0]{\popdisplay\fontsize{24}{24}\selectfont\color{poppurple}+};%
      \hspace{8pt}\begin{minipage}[c]{0.6\linewidth}
        {\popdisplay\fontsize{11}{12.5}\selectfont\color{white}\MakeUppercase{Go OnBuch+}}\par\vspace{4pt}
        {\raggedright\popsemi\fontsize{8.4}{11}\selectfont\color{white}All signed courses, unlimited AI tutor, PDF sheets — OnBuch+ tab in the app\par}
      \end{minipage}
    \end{tcolorbox}
  \end{minipage}
  \vspace{8pt}
  \popcontenu
  \vfill
  \signatureonbuch
  \restoregeometry
  \newpage
}

% Rangée « Ce cours contient » (page de garde)
\newcommand{\poptile}[3]{% #1 couleur #2 gros texte #3 légende
  \begin{tcolorbox}[enhanced,colback=#1,colframe=popink,boxrule=2pt,arc=9pt,shadow={2.5pt}{-2.5pt}{0pt}{popink},
      left=6pt,right=6pt,top=9pt,bottom=9pt,halign=center,valign=center,height=1.75cm,width=\linewidth,nobeforeafter]
    {\popdisplay\fontsize{12.5}{13}\selectfont\color{white}\MakeUppercase{#2}}\par\vspace{3pt}
    {\centering\popsemi\fontsize{7.8}{9.5}\selectfont\color{white}#3\par}
  \end{tcolorbox}}
\newcommand{\popcontenu}{%
  \par\noindent{\popxboldls\fontsize{7.5}{7.5}\selectfont\color{popmuted}THIS SIGNED COURSE CONTAINS}\par\vspace{7pt}
  \noindent\begin{minipage}[t]{0.235\linewidth}\poptile{popblue}{Course}{complete, illustrated, step by step}\end{minipage}\hfill
  \begin{minipage}[t]{0.235\linewidth}\poptile{poporange}{Methods}{and worked examples}\end{minipage}\hfill
  \begin{minipage}[t]{0.235\linewidth}\poptile{poppink}{Exercises}{exam style + solutions}\end{minipage}\hfill
  \begin{minipage}[t]{0.235\linewidth}\poptile{popgreen}{Summary sheet}{the essentials to remember}\end{minipage}\par}

% Sommaire
\newtcolorbox{sommaire}{enhanced,colback=white,colframe=popink,boxrule=2.2pt,arc=10pt,
  drop shadow=popink,left=18pt,right=18pt,top=13pt,bottom=13pt,before skip=6pt,after skip=20pt,
  before upper={\poppill{Contents}{poppurple}{white}\par\vspace{9pt}\popsemi\fontsize{10.6}{14.5}\selectfont\color{popink}}}

% Signature de fin de cours — poussée tout en bas de la dernière page
\newcommand{\findecours}{%
  \par\vfill\nopagebreak
  \begin{center}\popsquiggle{poporange}\end{center}\vspace{4pt}
  \signatureonbuch
}
\AtBeginDocument{\def\llbracket{\mathopen{[\mkern-3.5mu[}}\def\rrbracket{\mathclose{]\mkern-3.5mu]}}} % intervalles d'entiers (glyphe ⟦ absent de la police)
% Glyphes absents de Fira Math : redéfinis à partir de glyphes disponibles
\AtBeginDocument{%
  \def\setminus{\mathbin{\backslash}}\let\smallsetminus\setminus
  \def\vdots{\mathord{\vbox{\baselineskip3.2pt\lineskiplimit0pt\kern4pt\hbox{.}\hbox{.}\hbox{.}}}}%
  \def\ddots{\mathinner{\mkern1mu\raise6.5pt\hbox{.}\mkern2mu\raise3.6pt\hbox{.}\mkern2mu\raise0.7pt\hbox{.}\mkern1mu}}%
  \def\triangle{\mathord{\scalebox{0.9}{$\symup{\Delta}$}}}%
  \def\nabla{\mathord{\raisebox{\depth}{\scalebox{1}[-1]{$\symup{\Delta}$}}}}%
  \def\checkmark{\popcheck{popdark}}%
  \def\star{\mathbin{\ast}}\def\bigstar{\mathord{\ast}}%
  \def\diamond{\mathbin{\scalebox{0.6}{$\Diamond$}}}%
  \def\bigcup{\mathop{\mathchoice{\raisebox{-0.3em}{\scalebox{1.7}{$\cup$}}}{\scalebox{1.25}{$\cup$}}{\cup}{\cup}}\displaylimits}%
  \def\bigcap{\mathop{\mathchoice{\raisebox{-0.3em}{\scalebox{1.7}{$\cap$}}}{\scalebox{1.25}{$\cap$}}{\cap}{\cap}}\displaylimits}%
  \def\lesseqqgtr{\mathrel{\lessgtr}}\def\bigoplus{\mathop{\oplus}\displaylimits}%
}
\providecommand{\ssi}{if and only if} % abréviation fréquente
\AtBeginDocument{\providecommand{\wideparen}{\overparen}} % arc de cercle

% Fonctions usuelles que les modèles utilisent sans que amsmath les fournisse
\providecommand{\arsinh}{\operatorname{arsinh}}\providecommand{\arcosh}{\operatorname{arcosh}}
\providecommand{\artanh}{\operatorname{artanh}}\providecommand{\arsech}{\operatorname{arsech}}
\providecommand{\arcsch}{\operatorname{arcsch}}\providecommand{\arcoth}{\operatorname{arcoth}}
\providecommand{\arcsinh}{\operatorname{arsinh}}\providecommand{\arccosh}{\operatorname{arcosh}}
\providecommand{\arctanh}{\operatorname{artanh}}\providecommand{\sech}{\operatorname{sech}}
\providecommand{\csch}{\operatorname{csch}}\providecommand{\arcsec}{\operatorname{arcsec}}
\providecommand{\arccsc}{\operatorname{arccsc}}\providecommand{\arccot}{\operatorname{arccot}}
\providecommand{\sgn}{\operatorname{sgn}}\providecommand{\lcm}{\operatorname{lcm}}
\providecommand{\Var}{\operatorname{Var}}\providecommand{\Cov}{\operatorname{Cov}}

%% FIGFIX-BEGIN v4 — correctifs globaux des figures (docs/onbuch-generation/tools/figfix)
\usepackage{adjustbox}\usepackage{etoolbox}
% toute figure TikZ/pgfplots plus large que la ligne est réduite à la largeur disponible
\BeforeBeginEnvironment{tikzpicture}{\begin{adjustbox}{max width=\linewidth}}
\AfterEndEnvironment{tikzpicture}{\end{adjustbox}}
% légendes pgfplots lisibles par-dessus les courbes
\pgfplotsset{every axis/.append style={legend style={fill=white,fill opacity=0.92,text opacity=1,draw=black!15,inner sep=3pt}}}
% étiquettes d'arêtes (syntaxe edge["..."]) avec fond blanc
\tikzset{every edge quotes/.append style={fill=white,inner sep=1.2pt}}
%% FIGFIX-END
\begin{document}

\pagegarde{Motion of Particles in two Dimensions}{Further Mathematics}{Upper Sixth}{Upper Sixth — Further mathematics}{11}{4 periods}{This chapter extends kinematics to plane motion: velocity and acceleration are studied as vectors, first in Cartesian coordinates, then in polar coordinates through the radial and transverse components. These tools, essential for the GCE Further Mathematics paper, allow the analysis of circular and general curvilinear motion.
\vspace{4pt}
\begin{multicols}{2}\small
\begin{itemize}
\item By the end of this chapter I can write the position vector of a particle in Cartesian and in polar form.
\item By the end of this chapter I can obtain velocity and acceleration by differentiating position vectors with respect to time.
\item By the end of this chapter I can compute Cartesian components of velocity and acceleration from parametric equations of a path.
\item By the end of this chapter I can differentiate the rotating unit vectors e\_r and e\_theta with respect to time.
\item By the end of this chapter I can derive the radial and transverse components of velocity (r-dot, r theta-dot).
\item By the end of this chapter I can derive the radial and transverse components of acceleration (r-ddot - r theta-dot\^{}2, r theta-ddot + 2 r-dot theta-dot).
\end{itemize}\end{multicols}}

\begin{sommaire}
\begin{multicols}{2}
\begin{enumerate}
\item Cartesian components of velocity and acceleration
\item Polar coordinates and the rotating basis
\item Derivation of radial-transverse components of velocity and acceleration
\item Simple cases of radial-transverse motion
\item Integration activity
\item Methods and worked examples
\item Exercises
\item Solutions to the exercises
\item Summary sheet
\end{enumerate}
\end{multicols}
\end{sommaire}

\begin{objectifs}
\begin{itemize}
\item By the end of this chapter I can write the position vector of a particle in Cartesian and in polar form.
\item By the end of this chapter I can obtain velocity and acceleration by differentiating position vectors with respect to time.
\item By the end of this chapter I can compute Cartesian components of velocity and acceleration from parametric equations of a path.
\item By the end of this chapter I can differentiate the rotating unit vectors e\_r and e\_theta with respect to time.
\item By the end of this chapter I can derive the radial and transverse components of velocity (r-dot, r theta-dot).
\item By the end of this chapter I can derive the radial and transverse components of acceleration (r-ddot - r theta-dot\^{}2, r theta-ddot + 2 r-dot theta-dot).
\item By the end of this chapter I can apply radial-transverse formulae to circular motion, including non-uniform circular motion.
\item By the end of this chapter I can solve GCE-style problems on particles moving in a plane and interpret the physical meaning of each component.
\end{itemize}
\end{objectifs}

\begin{prerequis}
\begin{itemize}
\item Differentiation of functions of one variable, including chain rule and product rule (Pure Mathematics, Lower Sixth)
\item Vectors in the plane: components, magnitude, unit vectors i and j (Lower Sixth)
\item Parametric equations of curves and parametric differentiation
\item One-dimensional kinematics: displacement, velocity, acceleration as rates of change (Mechanics, Lower Sixth)
\item Radian measure and basic trigonometry (sin, cos of a varying angle)
\end{itemize}
\end{prerequis}

\newpage

\coursec{Cartesian components of velocity and acceleration}

At the Douala seaport, a crane swings a container in a circular arc while simultaneously hoisting it upward: how fast is the container really moving, and in which direction? When a cyclist leans into the roundabout at Carrefour Warda in Yaoundé, what acceleration keeps her on the curved path? Describing motion along a straight line with a single coordinate is no longer enough — we need two-dimensional descriptions. Cartesian coordinates suit some paths, but for rotation, polar coordinates and radial-transverse components are far more natural. This chapter gives you the mathematical tools to analyse any planar motion, whatever its shape. We begin with the Cartesian description, which you already met informally in Lower Sixth through projectile motion.

\courssub{Position vector and parametric trajectory}

Consider a particle moving in a plane. Choose an origin $O$ and two perpendicular axes $Ox$ and $Oy$ with unit vectors $\vec{i}$ and $\vec{j}$. At time $t$, the particle is at point $M$ with coordinates $(x(t), y(t))$. Its \cle{position vector} is
\[
\vec{r}(t) = x(t)\,\vec{i} + y(t)\,\vec{j}.
\]
As $t$ varies, the tip of $\vec{r}(t)$ traces a curve $\mathcal{C}$ — the \cle{trajectory} of the particle. The equations $x = x(t)$, $y = y(t)$ give a \cle{parametric representation} of $\mathcal{C}$ with parameter $t$ (time).

\begin{definition}[Position vector of a particle in plane motion]
The position vector of a particle at time $t$ is $\vec{r}(t) = x(t)\vec{i} + y(t)\vec{j}$, where $x(t)$ and $y(t)$ are the Cartesian coordinates of the particle. The trajectory is the set of points $M(x(t), y(t))$ for $t$ in the domain of motion.
\end{definition}

\begin{popfigure}
\begin{tikzpicture}[scale=1.2, >=stealth]
% Axes
\draw[->] (-0.5,0) -- (5,0) node[right] {$x$};
\draw[->] (0,-0.5) -- (0,4) node[above] {$y$};
% Parametric curve (example: x = t, y = t^2/4 for t in [0,4])
\draw[popblue, thick, domain=0:4, smooth, variable=\t] plot ({\t}, {\t*\t/4});
% Points at t=1, t=2, t=3
\foreach \t/\pos in {1/below right, 2/above right, 3/above left} {
  \pgfmathsetmacro{\x}{\t}
  \pgfmathsetmacro{\y}{\t*\t/4}
  \fill[popink] (\x,\y) circle (2pt);
  \node[\pos] at (\x,\y) {$t=\t$};
}
% Position vector at t=2
\draw[->, popdark, thick] (0,0) -- (2,1) node[midway, above left] {$\vec{r}(t)$};
% Velocity vector (tangent) at t=2: derivative (1, t/2) -> (1,1)
\draw[->, poporange, thick] (2,1) -- ++(0.7,0.7) node[midway, above right] {$\vec{v}(t)$};
% Acceleration vector at t=2: derivative of velocity (0, 1/2)
\draw[->, popgreen, thick] (2,1) -- ++(0,0.5) node[midway, left] {$\vec{a}(t)$};
% Origin
\node[below left] at (0,0) {$O$};
\end{tikzpicture}
\legende{A particle's trajectory $\mathcal{C}$ with position vector $\vec{r}(t)$, velocity $\vec{v}(t)$ (tangent) and acceleration $\vec{a}(t)$ at a point.}
\end{popfigure}

\courssub{Velocity vector: definition, geometric meaning, speed and direction}

The \cle{velocity vector} is the time derivative of the position vector. Since $\vec{i}$ and $\vec{j}$ are \emph{fixed} unit vectors (they do not change with time), differentiation acts only on the coordinates:
\[
\vec{v}(t) = \frac{d\vec{r}}{dt} = \frac{dx}{dt}\,\vec{i} + \frac{dy}{dt}\,\vec{j} = \dot{x}(t)\,\vec{i} + \dot{y}(t)\,\vec{j}.
\]

\begin{propriete}[Velocity is tangent to the trajectory]
The velocity vector $\vec{v}(t)$ is always tangent to the trajectory at the point $M(t)$, and points in the direction of motion at that instant. (Proof not examinable; the idea is given below.)
\end{propriete}

\textbf{Justification.} For a small $\Delta t$, the displacement is $\Delta\vec{r} = \vec{r}(t+\Delta t) - \vec{r}(t)$, which lies along the chord $MM'$ joining two nearby points of the path. The vector $\Delta\vec{r}/\Delta t$ has the same direction as this chord. As $\Delta t \to 0$, the chord tends to the tangent at $M$, and $\Delta\vec{r}/\Delta t \to d\vec{r}/dt = \vec{v}(t)$. Hence $\vec{v}(t)$ is directed along the tangent to the trajectory.

The \cle{speed} (magnitude of the velocity) is
\[
v(t) = |\vec{v}(t)| = \sqrt{\dot{x}^2 + \dot{y}^2}.
\]
The \cle{direction of motion} is given by the angle $\alpha$ that $\vec{v}$ makes with the positive $x$-axis:
\[
\tan\alpha = \frac{\dot{y}}{\dot{x}}, \qquad \text{provided } \dot{x} \neq 0.
\]

\begin{definition}[Velocity vector, speed and direction]
For a particle with position vector $\vec{r}(t) = x(t)\vec{i} + y(t)\vec{j}$:
\begin{itemize}
\item Velocity: $\vec{v}(t) = \dot{x}(t)\vec{i} + \dot{y}(t)\vec{j}$.
\item Speed: $v(t) = |\vec{v}(t)| = \sqrt{\dot{x}^2 + \dot{y}^2}$.
\item Direction angle $\alpha$ with $Ox$: $\tan\alpha = \dot{y}/\dot{x}$ (if $\dot{x}\neq0$); if $\dot{x}=0$, then $\alpha = \pi/2$ or $3\pi/2$ according to the sign of $\dot{y}$.
\end{itemize}
\end{definition}

\begin{methode}[Finding velocity, speed and direction in Cartesian coordinates]
\begin{enumerate}
\item Write $x(t)$ and $y(t)$ explicitly.
\item Differentiate: $\dot{x}(t) = dx/dt$, $\dot{y}(t) = dy/dt$.
\item Velocity vector: $\vec{v}(t) = \dot{x}\vec{i} + \dot{y}\vec{j}$.
\item Speed: $v = \sqrt{\dot{x}^2 + \dot{y}^2}$.
\item Direction: compute $\tan\alpha = \dot{y}/\dot{x}$ and determine the quadrant from the signs of $\dot{x}$ and $\dot{y}$.
\end{enumerate}
\end{methode}

\courssub{Acceleration vector}

The \cle{acceleration vector} is the time derivative of the velocity vector:
\[
\vec{a}(t) = \frac{d\vec{v}}{dt} = \frac{d^2\vec{r}}{dt^2} = \ddot{x}(t)\,\vec{i} + \ddot{y}(t)\,\vec{j}.
\]
Its magnitude is $a(t) = |\vec{a}(t)| = \sqrt{\ddot{x}^2 + \ddot{y}^2}$.

\begin{definition}[Acceleration vector]
The acceleration vector of a particle in plane motion is $\vec{a}(t) = \ddot{x}(t)\vec{i} + \ddot{y}(t)\vec{j}$. It measures the rate of change of the velocity \emph{vector} — that is, changes in the speed, in the direction of motion, or in both.
\end{definition}

\begin{attention}[Acceleration is not generally tangent to the path]
Unlike the velocity, the acceleration is \textbf{not} tangent to the trajectory in general. For example, in circular motion the acceleration points towards the centre while the velocity is tangent. Never draw $\vec{a}$ along the path unless you have verified it.
\end{attention}

\begin{attention}[Speed vs. derivative of the distance from the origin]
Do \textbf{not} confuse the speed $v = |d\vec{r}/dt|$ with the rate of change of the distance from the origin $d|\vec{r}|/dt$. In general,
\[
\frac{d}{dt}|\vec{r}| \neq \left|\frac{d\vec{r}}{dt}\right|.
\]
For example, in uniform circular motion $|\vec{r}|$ is constant so $d|\vec{r}|/dt = 0$, but the speed $v > 0$.
\end{attention}

\courssub{Worked examples}

\begin{exoresolu}[Parabolic flight of a thrown stone]
A stone is thrown from ground level with initial velocity $\vec{u} = 10\vec{i} + 20\vec{j}$ (m/s). Take $g = 10$ m/s$^2$ downward. Find:
\begin{enumerate}
\item The position vector $\vec{r}(t)$ at time $t$.
\item The velocity $\vec{v}(t)$ and speed $v(t)$.
\item The time when the stone is at its highest point and the maximum height.
\item The angle of the velocity with the horizontal when $t = 1$ s.
\end{enumerate}
\tcblower
\textbf{Solution.}
Take the origin at the launch point, $Ox$ horizontal, $Oy$ vertical upward. The acceleration is $\vec{a} = -g\vec{j} = -10\vec{j}$.

Integrating with $\vec{v}(0) = \vec{u}$: $\vec{v}(t) = \vec{u} + \vec{a}t = 10\vec{i} + (20 - 10t)\vec{j}$.

Integrating again with $\vec{r}(0) = \vec{0}$: $\vec{r}(t) = 10t\,\vec{i} + (20t - 5t^2)\vec{j}$.

Speed: $v(t) = \sqrt{10^2 + (20-10t)^2} = 10\sqrt{1 + (2-t)^2}$ m/s.

Highest point when the vertical component of velocity is zero: $20 - 10t = 0 \Rightarrow t = 2$ s.

Maximum height: $y(2) = 20(2) - 5(2^2) = 40 - 20 = 20$ m.

At $t=1$: $\vec{v}(1) = 10\vec{i} + 10\vec{j}$, so $\tan\alpha = 10/10 = 1 \Rightarrow \alpha = 45^\circ$ above the horizontal.
\end{exoresolu}

\begin{exoresolu}[Elliptical parametric path]
A particle moves so that its position vector is $\vec{r}(t) = 3\cos t\,\vec{i} + 2\sin t\,\vec{j}$ ($t$ in seconds, distances in metres). Find:
\begin{enumerate}
\item The Cartesian equation of the trajectory.
\item The velocity and acceleration vectors.
\item The speed at $t = \pi/4$.
\item Show that $\vec{a}$ is always directed towards the origin.
\end{enumerate}
\tcblower
\textbf{Solution.}
\begin{enumerate}
\item $x = 3\cos t$, $y = 2\sin t$. Using $\cos^2 t + \sin^2 t = 1$: $\dfrac{x^2}{9} + \dfrac{y^2}{4} = 1$ — an ellipse with semi-axes $3$ and $2$.
\item $\vec{v}(t) = -3\sin t\,\vec{i} + 2\cos t\,\vec{j}$; \quad $\vec{a}(t) = -3\cos t\,\vec{i} - 2\sin t\,\vec{j} = -\vec{r}(t)$.
\item At $t = \pi/4$: $\dot{x} = -3\sin(\pi/4) = -\dfrac{3}{\sqrt{2}}$, $\dot{y} = 2\cos(\pi/4) = \dfrac{2}{\sqrt{2}}$.
\[
v = \sqrt{\frac{9}{2} + \frac{4}{2}} = \sqrt{\frac{13}{2}} \approx 2.55\ \mathrm{m\,s^{-1}}.
\]
\item Since $\vec{a}(t) = -\vec{r}(t)$, the acceleration is opposite to the position vector, i.e. always directed towards the origin $O$.
\end{enumerate}
\end{exoresolu}

\begin{exoresolu}[A particle that is never at rest]
A particle has coordinates $x = 2\cos 3t$, $y = \sin 2t$. Find the velocity and acceleration vectors. Determine whether the particle is ever momentarily at rest.
\tcblower
\textbf{Solution.}
$\dot{x} = -6\sin 3t$, $\dot{y} = 2\cos 2t$, so
\[
\vec{v}(t) = -6\sin 3t\,\vec{i} + 2\cos 2t\,\vec{j}, \qquad
\vec{a}(t) = -18\cos 3t\,\vec{i} - 4\sin 2t\,\vec{j}.
\]
The particle is at rest when $\vec{v} = \vec{0}$, i.e. $\sin 3t = 0$ \emph{and} $\cos 2t = 0$ simultaneously.

$\sin 3t = 0 \Rightarrow 3t = k\pi \Rightarrow t = \dfrac{k\pi}{3}$ ($k \in \mathbb{Z}$).

$\cos 2t = 0 \Rightarrow 2t = \dfrac{\pi}{2} + m\pi \Rightarrow t = \dfrac{\pi}{4} + \dfrac{m\pi}{2}$ ($m \in \mathbb{Z}$).

A common value would require $\dfrac{k\pi}{3} = \dfrac{\pi}{4} + \dfrac{m\pi}{2}$, i.e. $4k = 3 + 6m$, i.e. $4k - 6m = 3$. The left-hand side is even and the right-hand side is odd: no integer solutions exist. Hence the particle is \textbf{never} at rest.
\end{exoresolu}

\courssub{Special cases: velocity perpendicular to acceleration; particle at rest}

\begin{propriete}[Condition for $\vec{v} \perp \vec{a}$]
\[
\vec{v} \perp \vec{a} \iff \vec{v} \cdot \vec{a} = 0 \iff \dot{x}\ddot{x} + \dot{y}\ddot{y} = 0.
\]
Since $\vec{v}\cdot\vec{a} = \dfrac{d}{dt}\!\left(\tfrac{1}{2}v^2\right)$, this condition means the \textbf{speed is stationary} at that instant (in particular, it holds throughout any motion at constant speed).
\end{propriete}

\begin{exoresolu}[When is velocity perpendicular to acceleration?]
For the motion $\vec{r}(t) = t^2\vec{i} + t^3\vec{j}$ ($t \geq 0$), find the times when $\vec{v} \perp \vec{a}$.
\tcblower
\textbf{Solution.}
$\vec{v} = 2t\,\vec{i} + 3t^2\,\vec{j}$, \quad $\vec{a} = 2\,\vec{i} + 6t\,\vec{j}$.
\[
\vec{v}\cdot\vec{a} = (2t)(2) + (3t^2)(6t) = 4t + 18t^3 = 2t(2 + 9t^2).
\]
Set to zero: $2t(2+9t^2)=0 \Rightarrow t=0$ (since $2+9t^2>0$ for all $t$).

At $t=0$, $\vec{v}=\vec{0}$: the particle is at rest, and the zero vector is conventionally perpendicular to every vector. For $t>0$, $\vec{v}\cdot\vec{a} > 0$, so the angle between $\vec{v}$ and $\vec{a}$ is acute: the speed is increasing.
\end{exoresolu}

\begin{aretenir}[Key points for Cartesian components]
\begin{itemize}
\item Position: $\vec{r}(t) = x(t)\vec{i} + y(t)\vec{j}$.
\item Velocity: $\vec{v}(t) = \dot{x}\vec{i} + \dot{y}\vec{j}$; tangent to the path; speed $v = \sqrt{\dot{x}^2+\dot{y}^2}$.
\item Acceleration: $\vec{a}(t) = \ddot{x}\vec{i} + \ddot{y}\vec{j}$; not tangent to the path in general.
\item Direction of $\vec{v}$: $\tan\alpha = \dot{y}/\dot{x}$, quadrant from the signs of $\dot{x}$, $\dot{y}$.
\item $\vec{v}\perp\vec{a} \iff \dot{x}\ddot{x}+\dot{y}\ddot{y}=0$ (speed stationary).
\item Particle at rest $\iff \dot{x}=\dot{y}=0$ simultaneously.
\end{itemize}
\end{aretenir}

\courssub{Link with projectiles (revision of Lower Sixth material)}

Projectile motion under constant gravity is the fundamental example of planar motion in Cartesian coordinates. With $Ox$ horizontal and $Oy$ vertical upward, $\vec{a} = -g\vec{j}$. If the initial velocity is $\vec{u} = u\cos\theta\,\vec{i} + u\sin\theta\,\vec{j}$, then
\begin{align*}
\vec{v}(t) &= u\cos\theta\,\vec{i} + (u\sin\theta - gt)\vec{j}, \\
\vec{r}(t) &= (u\cos\theta)\,t\,\vec{i} + \left(u\sin\theta\,t - \tfrac{1}{2}gt^2\right)\vec{j}.
\end{align*}
Eliminating $t$ gives the trajectory, a parabola:
\[
y = x\tan\theta - \frac{gx^2}{2u^2\cos^2\theta}.
\]
Time of flight $T = \dfrac{2u\sin\theta}{g}$, range $R = \dfrac{u^2\sin2\theta}{g}$, maximum height $H = \dfrac{u^2\sin^2\theta}{2g}$. These results are \textbf{assumed known} from Lower Sixth; they will be used without re-derivation in GCE questions unless a derivation is explicitly requested.

\begin{methode}[GCE exam technique: speed and direction at a given instant]
Questions often ask: ``Find the speed and the angle the velocity makes with the horizontal at $t = 2$ s.''
\begin{enumerate}
\item Compute $\dot{x}(2)$ and $\dot{y}(2)$.
\item Speed $v = \sqrt{\dot{x}^2 + \dot{y}^2}$, with units (m/s).
\item $\tan\alpha = \dot{y}/\dot{x}$; give $\alpha$ in degrees to one decimal place, or as an exact value such as $\alpha = \arctan(3/4)$.
\item State the quadrant clearly if $\dot{x}<0$ (angle measured anticlockwise from the positive $x$-axis), and say ``above'' or ``below'' the horizontal.
\end{enumerate}
\end{methode}

\begin{popfigure}
\begin{tikzpicture}
\begin{axis}[
    axis lines=middle,
    xlabel=$x$ (m), ylabel=$y$ (m),
    xmin=0, xmax=44,
    ymin=0, ymax=12,
    xtick={0,10,20,30,40},
    ytick={0,5,10},
    grid=major,
    grid style={popline},
    width=0.9\linewidth,
    height=0.5\linewidth,
    clip=false
]
% Projectile: u=20, theta=45 deg, g=10
% x(t) = 10*sqrt(2)*t, y(t) = 10*sqrt(2)*t - 5 t^2 ; eliminate t: y = x - x^2/40 ; range 40 m
\addplot[popblue, thick, domain=0:40, samples=100] {x - x^2/40};
% Velocity vectors at t = 0.5, 1.0, 1.5 s
% position: x = 14.1421 t, y = 14.1421 t - 5 t^2
% velocity: vx = 14.1421, vy = 14.1421 - 10 t ; scale 0.3
\draw[->, poporange, thick] (7.07,5.82) -- ++(4.24,2.74);
\node[poporange, anchor=south east, font=\small] at (7.07,5.82) {$t=0.5$};
\draw[->, poporange, thick] (14.14,9.14) -- ++(4.24,1.24);
\node[poporange, anchor=south east, font=\small] at (14.14,9.14) {$t=1.0$};
\draw[->, poporange, thick] (21.21,9.96) -- ++(4.24,-0.26);
\node[poporange, anchor=south west, font=\small] at (21.21,9.96) {$t=1.5$};
\end{axis}
\end{tikzpicture}
\legende{Parabolic trajectory of a projectile ($u=20$ m/s, $\theta=45^\circ$, $g=10$ m/s$^2$) with velocity vectors at $t=0.5$, $1.0$, $1.5$ s. The velocity is always tangent to the path; its vertical component decreases steadily.}
\end{popfigure}

\begin{exercice}[Practice]
\begin{enumerate}
\item A particle moves with $\vec{r}(t) = (t^3 - 3t)\vec{i} + 2t^2\vec{j}$. Find $\vec{v}(t)$, $\vec{a}(t)$, the speed at $t=2$, and the angle of $\vec{v}$ with $Ox$ at $t=2$.
\item A particle has velocity $\vec{v}(t) = (4t-2)\vec{i} + 3t^2\vec{j}$. Given $\vec{r}(0) = \vec{i} + 2\vec{j}$, find $\vec{r}(t)$. Is the particle ever at rest?
\item For the motion $\vec{r}(t) = e^t\vec{i} + e^{-t}\vec{j}$, compute $\vec{v}\cdot\vec{a}$ as a function of $t$. Deduce the unique instant at which $\vec{v} \perp \vec{a}$, and describe how the speed behaves before and after that instant.
\item A projectile is fired from the top of a cliff 50 m high with initial velocity $30\vec{i} + 40\vec{j}$ (m/s), where $Oy$ is vertical upward from the base of the cliff. Take $g=10$ m/s$^2$. Find the time of flight, the horizontal distance travelled, and the velocity vector at impact.
\end{enumerate}
\end{exercice}

\begin{corrige}[Exercise 1]
$\vec{v} = (3t^2-3)\vec{i} + 4t\vec{j}$, \quad $\vec{a} = 6t\vec{i} + 4\vec{j}$.

At $t=2$: $\vec{v} = 9\vec{i} + 8\vec{j}$, speed $v = \sqrt{81+64} = \sqrt{145} \approx 12.04$ m/s.

$\tan\alpha = 8/9 \Rightarrow \alpha \approx 41.6^\circ$ above the positive $x$-axis.
\end{corrige}

\begin{corrige}[Exercise 2]
$\vec{r}(t) = \displaystyle\int \vec{v}\,dt = (2t^2-2t)\vec{i} + t^3\vec{j} + \vec{C}$.

$\vec{r}(0) = \vec{C} = \vec{i}+2\vec{j} \Rightarrow \vec{r}(t) = (2t^2-2t+1)\vec{i} + (t^3+2)\vec{j}$.

At rest we need $4t-2=0$ and $3t^2=0$, i.e. $t=1/2$ and $t=0$ simultaneously — impossible. The particle is never at rest.
\end{corrige}

\begin{corrige}[Exercise 3]
$\vec{v} = e^t\vec{i} - e^{-t}\vec{j}$, \quad $\vec{a} = e^t\vec{i} + e^{-t}\vec{j}$.
\[
\vec{v}\cdot\vec{a} = e^{2t} - e^{-2t}.
\]
This vanishes only when $e^{2t} = e^{-2t}$, i.e. $t = 0$. So $\vec{v} \perp \vec{a}$ only at $t=0$.

Since $\vec{v}\cdot\vec{a} = \dfrac{d}{dt}\!\left(\tfrac{1}{2}v^2\right)$: for $t<0$ the dot product is negative (speed decreasing), and for $t>0$ it is positive (speed increasing). The speed is minimum at $t=0$, where $v(0) = \sqrt{1+1} = \sqrt{2}$.
\end{corrige}

\begin{corrige}[Exercise 4]
With the origin at the base of the cliff: $\vec{r}(t) = 30t\,\vec{i} + (50 + 40t - 5t^2)\vec{j}$.

Impact when $y=0$: $50+40t-5t^2=0 \Rightarrow t^2-8t-10=0 \Rightarrow t = \dfrac{8+\sqrt{64+40}}{2} = 4+\sqrt{26} \approx 9.10$ s (positive root).

Horizontal distance: $x = 30(4+\sqrt{26}) \approx 273$ m.

Velocity at impact: $\vec{v} = 30\vec{i} + (40 - 10t)\vec{j} = 30\vec{i} - 10\sqrt{26}\,\vec{j} \approx 30\vec{i} - 51.0\vec{j}$ m/s.
\end{corrige}

\coursec{Polar coordinates and the rotating basis}

In the previous section we saw how to describe motion using Cartesian coordinates $(x,y)$ with the fixed basis vectors $\mathbf{i}$ and $\mathbf{j}$. This is perfect when a particle moves along straight lines or parabolic trajectories under uniform gravity. But many important motions — a satellite orbiting the Earth, a bead sliding on a rotating wire, a particle moving along a spiral — are \emph{naturally} described by a distance from a fixed centre and an angle. In such cases, insisting on Cartesian coordinates makes the algebra unnecessarily heavy. Polar coordinates $(r,\theta)$ and their associated \emph{rotating basis} $(\mathbf{e}_r,\mathbf{e}_\theta)$ are the tailor‑made tool for these situations.

\courssub{Polar coordinates and the rotating unit vectors}

Consider a particle $P$ moving in a plane. Choose a fixed origin $O$ and a fixed initial line (usually the positive $x$-axis). The position of $P$ is completely determined by:
\begin{itemize}
  \item the distance $r = OP \ge 0$ (the \emph{radial coordinate}),
  \item the angle $\theta$ measured anticlockwise from the initial line to $OP$ (the \emph{angular coordinate} or \emph{polar angle}).
\end{itemize}
The pair $(r,\theta)$ are the \textbf{polar coordinates} of $P$. The relations with Cartesian coordinates $(x,y)$ are the familiar ones:
\[
x = r\cos\theta,\qquad y = r\sin\theta,\qquad r = \sqrt{x^2+y^2},\qquad \tan\theta = \frac{y}{x}\quad (x\ne 0).
\]

\begin{popfigure}
\begin{tikzpicture}[scale=1.2, >=stealth]
  % Axes
  \draw[->] (-0.5,0) -- (4,0) node[right] {$x$};
  \draw[->] (0,-0.5) -- (0,3.5) node[above] {$y$};
  % Point P
  \coordinate (O) at (0,0);
  \coordinate (P) at (3.2,2.0);
  \draw[thick, popblue] (O) -- (P) node[midway, above left] {$r$};
  \fill[popblue] (P) circle (2pt) node[above right] {$P(r,\theta)$};
  % Angle theta
  \draw[poporange, thick] (0.8,0) arc (0:32:0.8);
  \node[poporange] at (1.1,0.3) {$\theta$};
  % Fixed basis
  \draw[->, thick, popdark] (O) -- (1.2,0) node[below right] {$\mathbf{i}$};
  \draw[->, thick, popdark] (O) -- (0,1.2) node[above left] {$\mathbf{j}$};
  % Rotating basis at P
  \draw[->, very thick, popgreen] (P) -- ++(0.6,0.375) node[above right] {$\mathbf{e}_r$};
  \draw[->, very thick, poppurple] (P) -- ++(-0.375,0.6) node[above left] {$\mathbf{e}_\theta$};
  % Dashed lines for projection
  \draw[dashed, thin] (P) -- (3.2,0);
  \draw[dashed, thin] (P) -- (0,2.0);
\end{tikzpicture}
\legende{Polar coordinates $(r,\theta)$ of a point $P$ and the rotating basis $(\mathbf{e}_r,\mathbf{e}_\theta)$ compared to the fixed Cartesian basis $(\mathbf{i},\mathbf{j})$.}
\end{popfigure}

At the point $P$ we define two unit vectors:
\begin{itemize}
  \item \textbf{Radial unit vector} $\mathbf{e}_r$: points from $O$ towards $P$, i.e. along $OP$ in the direction of increasing $r$.
  \item \textbf{Transverse unit vector} $\mathbf{e}_\theta$: perpendicular to $\mathbf{e}_r$, pointing in the direction of increasing $\theta$ (anticlockwise).
\end{itemize}

\begin{definition}[Polar unit vectors]
In terms of the fixed Cartesian basis $\{\mathbf{i},\mathbf{j}\}$,
\[
\boxed{\mathbf{e}_r = \cos\theta\,\mathbf{i} + \sin\theta\,\mathbf{j}},\qquad
\boxed{\mathbf{e}_\theta = -\sin\theta\,\mathbf{i} + \cos\theta\,\mathbf{j}}.
\]
These vectors are orthonormal: $\mathbf{e}_r\cdot\mathbf{e}_r = 1$, $\mathbf{e}_\theta\cdot\mathbf{e}_\theta = 1$, $\mathbf{e}_r\cdot\mathbf{e}_\theta = 0$, and $\mathbf{e}_r \times \mathbf{e}_\theta = \mathbf{k}$ (out of the page).
\end{definition}

Notice that $\mathbf{e}_r$ and $\mathbf{e}_\theta$ \emph{depend on $\theta$}. As the particle moves, $\theta$ changes, so the basis vectors rotate. This is the crucial difference from $\mathbf{i}$ and $\mathbf{j}$ which are constant in space.

The position vector of $P$ takes a beautifully simple form in this basis:
\[
\boxed{\mathbf{r} = r\,\mathbf{e}_r}.
\]
There is no $\mathbf{e}_\theta$ component because the position vector is purely radial.

\courssub{Differentiation of the rotating basis}

Since $\mathbf{e}_r$ and $\mathbf{e}_\theta$ vary with $\theta$, and $\theta$ varies with time $t$, we must differentiate them with respect to $t$. Using the chain rule:
\[
\frac{d\mathbf{e}_r}{dt} = \frac{d\mathbf{e}_r}{d\theta}\,\frac{d\theta}{dt},\qquad
\frac{d\mathbf{e}_\theta}{dt} = \frac{d\mathbf{e}_\theta}{d\theta}\,\frac{d\theta}{dt}.
\]
From the definitions:
\[
\frac{d\mathbf{e}_r}{d\theta} = -\sin\theta\,\mathbf{i} + \cos\theta\,\mathbf{j} = \mathbf{e}_\theta,\qquad
\frac{d\mathbf{e}_\theta}{d\theta} = -\cos\theta\,\mathbf{i} - \sin\theta\,\mathbf{j} = -\mathbf{e}_r.
\]
Denoting $\dot{\theta} = \dfrac{d\theta}{dt}$ (the \emph{angular velocity}), we obtain the fundamental formulas:

\begin{propriete}[Derivatives of the polar unit vectors]
\[
\boxed{\frac{d\mathbf{e}_r}{dt} = \dot{\theta}\,\mathbf{e}_\theta},\qquad
\boxed{\frac{d\mathbf{e}_\theta}{dt} = -\dot{\theta}\,\mathbf{e}_r}.
\]
These formulas are \textbf{examinable} and must be memorised. They express the fact that the basis $(\mathbf{e}_r,\mathbf{e}_\theta)$ rotates with angular velocity $\dot{\theta}$.
\end{propriete}

\begin{popfigure}
\begin{tikzpicture}[scale=1.3, >=stealth]
  % Two positions
  \coordinate (O) at (0,0);
  % Position 1
  \coordinate (P1) at (2.5,1.0);
  \draw[thick, popblue] (O) -- (P1);
  \fill[popblue] (P1) circle (2pt) node[below right] {$P_1$};
  \draw[->, very thick, popgreen] (P1) -- ++(0.923,0.385) node[above right] {$\mathbf{e}_r^{(1)}$};
  \draw[->, very thick, poppurple] (P1) -- ++(-0.385,0.923) node[above left] {$\mathbf{e}_\theta^{(1)}$};
  % Position 2
  \coordinate (P2) at (1.5,2.2);
  \draw[thick, popblue] (O) -- (P2);
  \fill[popblue] (P2) circle (2pt) node[above left] {$P_2$};
  \draw[->, very thick, popgreen] (P2) -- ++(0.557,0.830) node[above right] {$\mathbf{e}_r^{(2)}$};
  \draw[->, very thick, poppurple] (P2) -- ++(-0.830,0.557) node[above left] {$\mathbf{e}_\theta^{(2)}$};
  % Angle delta theta
  \draw[poporange, thick] (0.7,0) arc (0:55:0.7);
  \node[poporange] at (1.1,0.5) {$\Delta\theta$};
  % Arc showing rotation
  \draw[popmuted, dashed] (P1) arc (21.8:55.5:2.69);
  \node[popmuted] at (3.5,1.5) {rotation of basis};
  % Fixed axes
  \draw[->, thin] (-0.3,0) -- (3.5,0) node[right] {$x$};
  \draw[->, thin] (0,-0.3) -- (0,3) node[above] {$y$};
\end{tikzpicture}
\legende{As $\theta$ increases, the basis $(\mathbf{e}_r,\mathbf{e}_\theta)$ rotates. The change in $\mathbf{e}_r$ is in the direction of $\mathbf{e}_\theta$, and the change in $\mathbf{e}_\theta$ is opposite to $\mathbf{e}_r$.}
\end{popfigure}

\courssub{Physical interpretation: the rotating basis}

The formulas $\dot{\mathbf{e}}_r = \dot{\theta}\mathbf{e}_\theta$ and $\dot{\mathbf{e}}_\theta = -\dot{\theta}\mathbf{e}_r$ have a clear geometric meaning. Imagine a rigid rod $OP$ rotating about $O$ with angular velocity $\dot{\theta}$. The unit vector $\mathbf{e}_r$ is fixed to the rod. Its tip moves on a unit circle with speed $\dot{\theta}$ (since radius $=1$). The velocity of the tip is tangential to the circle, i.e. in the $\mathbf{e}_\theta$ direction, with magnitude $\dot{\theta}$. Hence $\dot{\mathbf{e}}_r = \dot{\theta}\mathbf{e}_\theta$. Similarly, $\mathbf{e}_\theta$ is also fixed to the rotating rod (it is just $\mathbf{e}_r$ rotated by $90^\circ$). Its tip moves on the same unit circle but its velocity is directed \emph{inwards} (towards $-\mathbf{e}_r$), giving $\dot{\mathbf{e}}_\theta = -\dot{\theta}\mathbf{e}_r$.

This interpretation is vital: \textbf{the polar basis is a rotating frame attached to the line $OP$}. Whenever you see $\dot{\mathbf{e}}_r$ or $\dot{\mathbf{e}}_\theta$, remember they are non‑zero precisely because the basis rotates.

\begin{savaistu}[Historical note]
The use of polar coordinates for orbital motion goes back to Newton's \emph{Principia} (1687). Newton derived the polar form of acceleration to prove that an inverse‑square central force produces elliptical orbits — a cornerstone of celestial mechanics. The rotating basis $(\mathbf{e}_r,\mathbf{e}_\theta)$ is essentially the \emph{Frenet frame} for planar curves when the curve is described relative to a fixed pole.
\end{savaistu}

\courssub{When are polar coordinates preferable?}

Polar coordinates simplify the mathematics whenever the motion has a natural centre of rotation or a radial symmetry. Typical situations include:
\begin{itemize}
  \item \textbf{Central force motion:} planets, satellites, charged particles in a Coulomb field (the force is along $\mathbf{e}_r$).
  \item \textbf{Motion on a rotating arm/wire:} a bead on a rotating spoke, a slider on a rotating rod.
  \item \textbf{Spiral trajectories:} $r = a\theta$, $r = ae^{b\theta}$ (e.g. a moth circling a light, the groove of a vinyl record).
  \item \textbf{Problems given in terms of $r$ and $\theta$:} e.g. ``a particle moves so that $r = 2\cos\theta$ and $\dot{\theta}=3$''.
\end{itemize}
In such cases, writing the equations of motion in Cartesian coordinates would introduce unnecessary trigonometric functions and coupled differential equations. The rotating basis \emph{diagonalises} the geometry: radial and transverse directions are the natural axes.

\courssub{Worked examples: converting between Cartesian and polar descriptions}

\begin{exoresolu}[Converting a Cartesian trajectory to polar form]
A particle moves along the curve $y = x^2$ in the Cartesian plane. Express its position vector in polar coordinates and find the polar unit vectors at the point where $x=1$.
\tcblower
\textbf{Solution.}
The Cartesian coordinates are $(x, x^2)$. The polar coordinates satisfy
\[
r = \sqrt{x^2 + x^4} = |x|\sqrt{1+x^2},\qquad \tan\theta = \frac{x^2}{x} = x \quad (x\ne 0).
\]
At $x=1$: $r = \sqrt{2}$, $\theta = \frac{\pi}{4}$.
The position vector is $\mathbf{r} = \sqrt{2}\,\mathbf{e}_r$.
The polar unit vectors at this point are
\[
\mathbf{e}_r = \cos\frac{\pi}{4}\,\mathbf{i} + \sin\frac{\pi}{4}\,\mathbf{j} = \frac{\sqrt{2}}{2}(\mathbf{i}+\mathbf{j}),\qquad
\mathbf{e}_\theta = -\sin\frac{\pi}{4}\,\mathbf{i} + \cos\frac{\pi}{4}\,\mathbf{j} = \frac{\sqrt{2}}{2}(-\mathbf{i}+\mathbf{j}).
\]
\textbf{Check:} $\mathbf{r} = \sqrt{2}\,\mathbf{e}_r = \mathbf{i}+\mathbf{j}$, which indeed corresponds to $(1,1)$ on $y=x^2$.
\end{exoresolu}

\begin{exoresolu}[Converting a polar description to Cartesian]
A particle moves such that its polar coordinates are $r = 2\cos\theta$ and $\theta = t$ (with $t$ in radians). Find its Cartesian coordinates $(x,y)$ as functions of $t$, and verify that the path is a circle.
\tcblower
\textbf{Solution.}
Given $r = 2\cos t$, $\theta = t$.
\[
x = r\cos\theta = (2\cos t)\cos t = 2\cos^2 t = 1 + \cos 2t,
\]
\[
y = r\sin\theta = (2\cos t)\sin t = \sin 2t.
\]
Eliminate $t$:
\[
(x-1)^2 + y^2 = \cos^2 2t + \sin^2 2t = 1.
\]
This is a circle of radius $1$ centred at $(1,0)$. The particle goes around this circle as $t$ increases.
\end{exoresolu}

\courssub{Method: Differentiating a vector expressed in the rotating basis}

Any vector $\mathbf{v}$ (velocity, acceleration, force, etc.) can be written in the polar basis as
\[
\mathbf{v} = v_r\,\mathbf{e}_r + v_\theta\,\mathbf{e}_\theta.
\]
To differentiate $\mathbf{v}$ with respect to time, we \textbf{must} use the product rule on both the components and the basis vectors:

\begin{methode}[Differentiation in the rotating basis]
To compute $\dfrac{d\mathbf{v}}{dt}$ for $\mathbf{v} = v_r\mathbf{e}_r + v_\theta\mathbf{e}_\theta$:
\begin{enumerate}
  \item Write $\mathbf{v} = v_r\mathbf{e}_r + v_\theta\mathbf{e}_\theta$.
  \item Differentiate using the product rule:
        \[
        \frac{d\mathbf{v}}{dt} = \dot{v}_r\mathbf{e}_r + v_r\dot{\mathbf{e}}_r + \dot{v}_\theta\mathbf{e}_\theta + v_\theta\dot{\mathbf{e}}_\theta.
        \]
  \item Substitute $\dot{\mathbf{e}}_r = \dot{\theta}\mathbf{e}_\theta$ and $\dot{\mathbf{e}}_\theta = -\dot{\theta}\mathbf{e}_r$:
        \[
        \frac{d\mathbf{v}}{dt} = (\dot{v}_r - v_\theta\dot{\theta})\,\mathbf{e}_r + (\dot{v}_\theta + v_r\dot{\theta})\,\mathbf{e}_\theta.
        \]
  \item Identify the radial and transverse components of the derivative.
\end{enumerate}
\end{methode}

This method will be used in the next section to derive the radial–transverse components of velocity and acceleration. For now, practise the differentiation of the basis vectors themselves.

\begin{exoresolu}[Differentiating a vector given in polar components]
A vector is given by $\mathbf{A} = r^2\mathbf{e}_r + r\dot{\theta}\mathbf{e}_\theta$, where $r$ and $\theta$ are functions of time. Find $\dfrac{d\mathbf{A}}{dt}$ in terms of $\mathbf{e}_r$ and $\mathbf{e}_\theta$.
\tcblower
\textbf{Solution.}
Let $v_r = r^2$, $v_\theta = r\dot{\theta}$.
Then $\dot{v}_r = 2r\dot{r}$, $\dot{v}_\theta = \dot{r}\dot{\theta} + r\ddot{\theta}$.
Using the method:
\[
\frac{d\mathbf{A}}{dt} = (\dot{v}_r - v_\theta\dot{\theta})\mathbf{e}_r + (\dot{v}_\theta + v_r\dot{\theta})\mathbf{e}_\theta
= (2r\dot{r} - r\dot{\theta}^2)\mathbf{e}_r + (\dot{r}\dot{\theta} + r\ddot{\theta} + r^2\dot{\theta})\mathbf{e}_\theta.
\]
\end{exoresolu}

\courssub{Common mistakes and how to avoid them}

\begin{attention}[The rotating basis trap]
\textbf{The most frequent error at GCE Advanced Level} is to treat $\mathbf{e}_r$ and $\mathbf{e}_\theta$ as constant vectors, i.e. to write
\[
\frac{d}{dt}(r\mathbf{e}_r) \stackrel{\text{WRONG}}{=} \dot{r}\mathbf{e}_r \quad\text{or}\quad \frac{d}{dt}(v_r\mathbf{e}_r + v_\theta\mathbf{e}_\theta) \stackrel{\text{WRONG}}{=} \dot{v}_r\mathbf{e}_r + \dot{v}_\theta\mathbf{e}_\theta.
\]
\textbf{This is false.} $\mathbf{e}_r$ and $\mathbf{e}_\theta$ rotate with the particle, so their time derivatives are $\dot{\theta}\mathbf{e}_\theta$ and $-\dot{\theta}\mathbf{e}_r$ respectively. \emph{Always} apply the product rule and substitute the basis derivatives. A missing $\dot{\theta}\mathbf{e}_\theta$ term will lose you the transverse component of velocity/acceleration and cost several marks.
\end{attention}

\begin{aretenir}[Key points of this section]
\begin{itemize}
  \item Polar coordinates: $x = r\cos\theta$, $y = r\sin\theta$; $r\ge 0$, $\theta$ measured anticlockwise from $Ox$.
  \item Rotating unit vectors: $\mathbf{e}_r = \cos\theta\,\mathbf{i}+\sin\theta\,\mathbf{j}$, $\mathbf{e}_\theta = -\sin\theta\,\mathbf{i}+\cos\theta\,\mathbf{j}$.
  \item Position vector: $\mathbf{r} = r\mathbf{e}_r$.
  \item Basis derivatives: $\dot{\mathbf{e}}_r = \dot{\theta}\mathbf{e}_\theta$, $\dot{\mathbf{e}}_\theta = -\dot{\theta}\mathbf{e}_r$ (memorise!).
  \item The basis rotates with angular velocity $\dot{\theta}$; this is why the derivatives are non‑zero.
  \item To differentiate any vector in this basis, use the product rule and substitute the basis derivatives.
  \item Polar coordinates are the natural choice for central forces, rotating constraints, and spiral paths.
\end{itemize}
\end{aretenir}

\begin{exercice}[Practice]
\begin{enumerate}
  \item A particle has polar coordinates $r = 3t$, $\theta = 2t^2$. Find $\mathbf{e}_r$ and $\mathbf{e}_\theta$ at $t=1$ in terms of $\mathbf{i}$ and $\mathbf{j}$.
  \item Express the Cartesian vector $\mathbf{v} = 3\mathbf{i} + 4\mathbf{j}$ in the polar basis at the point where $\theta = 30^\circ$.
  \item If $\mathbf{F} = (r\dot{\theta}^2)\mathbf{e}_r + (r\ddot{\theta}+2\dot{r}\dot{\theta})\mathbf{e}_\theta$, compute $\dfrac{d\mathbf{F}}{dt}$ in terms of $\mathbf{e}_r$, $\mathbf{e}_\theta$ and the time derivatives of $r$ and $\theta$.
  \item A point moves on the curve $r = \theta$ (an Archimedean spiral) with constant angular speed $\dot{\theta} = \omega$. Write its position vector and differentiate it once to find the velocity vector in polar components.
\end{enumerate}
\end{exercice}

\begin{corrige}[Exercises 1--4]
\begin{enumerate}
  \item At $t=1$: $r=3$, $\theta=2$ rad. $\mathbf{e}_r = \cos 2\,\mathbf{i} + \sin 2\,\mathbf{j}$, $\mathbf{e}_\theta = -\sin 2\,\mathbf{i} + \cos 2\,\mathbf{j}$.
  \item $\theta=30^\circ = \pi/6$. $\mathbf{e}_r = \frac{\sqrt{3}}{2}\mathbf{i}+\frac{1}{2}\mathbf{j}$, $\mathbf{e}_\theta = -\frac{1}{2}\mathbf{i}+\frac{\sqrt{3}}{2}\mathbf{j}$. Solve $v_r\mathbf{e}_r+v_\theta\mathbf{e}_\theta = 3\mathbf{i}+4\mathbf{j}$:
        $v_r = 3\cos\theta+4\sin\theta = \frac{3\sqrt{3}}{2}+2$, $v_\theta = -3\sin\theta+4\cos\theta = -\frac{3}{2}+2\sqrt{3}$.
  \item Apply the method: $\dot{F}_r = \dot{r}\dot{\theta}^2 + 2r\dot{\theta}\ddot{\theta}$, $\dot{F}_\theta = \dot{r}\ddot{\theta}+r\dddot{\theta}+2\ddot{r}\dot{\theta}+2\dot{r}\ddot{\theta}$.
        $\frac{d\mathbf{F}}{dt} = (\dot{F}_r - F_\theta\dot{\theta})\mathbf{e}_r + (\dot{F}_\theta + F_r\dot{\theta})\mathbf{e}_\theta$.
  \item $\mathbf{r} = \theta\mathbf{e}_r$. $\mathbf{v} = \dot{\mathbf{r}} = \dot{\theta}\mathbf{e}_r + \theta\dot{\mathbf{e}}_r = \omega\mathbf{e}_r + \theta\omega\mathbf{e}_\theta = \omega(\mathbf{e}_r + \theta\mathbf{e}_\theta)$.
\end{enumerate}
\end{corrige}

We are now fully equipped to derive the radial and transverse components of velocity and acceleration — the core of the next section. The rotating basis is not just a notation; it is a physical frame that moves with the particle's radial line, and its time derivatives encode the geometry of rotation.

\coursec{Derivation of radial-transverse components of velocity and acceleration}

In the previous section we introduced polar coordinates $(r,\theta)$ and the \cle{rotating basis} $(\vec{e}_r,\vec{e}_\theta)$, and we established the crucial differentiation rules
\[
\frac{d\vec{e}_r}{dt}=\dot{\theta}\,\vec{e}_\theta,
\qquad
\frac{d\vec{e}_\theta}{dt}=-\dot{\theta}\,\vec{e}_r .
\]
We are now ready to harvest the fruit of that work. Starting from the remarkably simple position vector $\vec{r}=r\,\vec{e}_r$, we shall \emph{derive} — not merely quote — the polar components of velocity and acceleration. This derivation is explicitly required by the syllabus and appears regularly at the GCE Advanced Level: examiners want to see the product rule applied to a moving basis, not a memorised formula.

\courssub{Setting up: the position vector in the rotating basis}

Consider a particle $P$ moving in a plane, with polar coordinates $(r(t),\theta(t))$ relative to a fixed pole $O$ and a fixed initial line. Its position vector is
\[
\vec{r}=\overrightarrow{OP}=r\,\vec{e}_r .
\]
This looks deceptively simple, but remember: \textbf{both} factors depend on time. The scalar $r(t)$ changes as the particle moves towards or away from $O$, and the unit vector $\vec{e}_r$ rotates as $\theta(t)$ changes. Differentiating $\vec{r}$ therefore requires the \cle{product rule} together with the basis differentiation rules recalled above.

\courssub{Velocity in polar coordinates}

\begin{propriete}[Velocity in polar components — proof examinable]
If a particle has polar coordinates $(r,\theta)$, then its velocity is
\[
\vec{v}=\dot{r}\,\vec{e}_r+r\dot{\theta}\,\vec{e}_\theta .
\]
The \cle{radial component} is $v_r=\dot{r}$ and the \cle{transverse component} is $v_\theta=r\dot{\theta}$.
\end{propriete}

\textbf{Derivation (to be reproduced in the examination).} Differentiate $\vec{r}=r\,\vec{e}_r$ with respect to time, applying the product rule:
\begin{align*}
\vec{v}=\frac{d\vec{r}}{dt}
&=\frac{d}{dt}\bigl(r\,\vec{e}_r\bigr)\\
&=\dot{r}\,\vec{e}_r+r\,\frac{d\vec{e}_r}{dt}\\
&=\dot{r}\,\vec{e}_r+r\bigl(\dot{\theta}\,\vec{e}_\theta\bigr)\\
&=\dot{r}\,\vec{e}_r+r\dot{\theta}\,\vec{e}_\theta . \qquad\blacksquare
\end{align*}

\textbf{Physical meaning of the two components.}
\begin{itemize}
\item The term $\dot{r}\,\vec{e}_r$ measures motion \emph{along the radius}: it is the velocity the particle would have if $\theta$ were frozen and only its distance from $O$ changed (like sliding along a rotating spoke).
\item The term $r\dot{\theta}\,\vec{e}_\theta$ measures motion \emph{around the pole}: it is the velocity due to the rotation of the radius vector. Note the factor $r$: the same angular speed $\dot{\theta}$ produces a larger transverse speed further from the pole — exactly why the tip of a ceiling fan blade moves faster than a point near the hub.
\end{itemize}

\begin{attention}[The transverse velocity is $r\dot{\theta}$, not $\dot{\theta}$]
$\dot{\theta}$ is an \emph{angular} speed (rad\,s$^{-1}$), not a linear speed. Multiplying by $r$ converts it to a linear speed. Writing $v_\theta=\dot{\theta}$ is a dimensional error and costs marks.
\end{attention}

Since $\vec{e}_r\perp\vec{e}_\theta$, the speed follows at once from Pythagoras:
\[
|\vec{v}|=\sqrt{\dot{r}^{2}+r^{2}\dot{\theta}^{2}} .
\]

\begin{popfigure}
\begin{center}
\begin{tikzpicture}[scale=1.1]
\draw[->,popmuted] (-0.5,0) -- (6.5,0) node[below left] {\small initial line};
\draw[->,popmuted] (0,-0.5) -- (0,4.2);
\fill (0,0) circle (1.5pt) node[below left] {$O$};
% spiral path r = 0.35*theta
\draw[popblue,thick,domain=0:5.6,samples=120,smooth] plot ({0.35*\x*cos(\x r)},{0.35*\x*sin(\x r)});
\node[popblue] at (4.6,2.6) {\small spiral path};
\pgfmathsetmacro{\th}{4.0}
\pgfmathsetmacro{\rr}{1.4}
\coordinate (P) at ({\rr*cos(\th r)},{\rr*sin(\th r)});
\fill[popdark] (P) circle (2pt) node[above right] {$P$};
% e_r and e_theta
\draw[->,popgreen,thick] (P) -- ++({1.3*cos(\th r)},{1.3*sin(\th r)}) node[right] {$\vec{e}_r$};
\draw[->,popgreen,thick] (P) -- ++({-1.3*sin(\th r)},{1.3*cos(\th r)}) node[above] {$\vec{e}_\theta$};
% velocity components
\draw[->,poporange,very thick] (P) -- ++({1.9*cos(\th r)},{1.9*sin(\th r)}) node[below right] {$\dot{r}\,\vec{e}_r$};
\draw[->,poppurple,very thick] (P) -- ++({-2.2*sin(\th r)},{2.2*cos(\th r)}) node[above left] {$r\dot{\theta}\,\vec{e}_\theta$};
% total velocity (sum)
\draw[->,popink,ultra thick] (P) -- ++({1.9*cos(\th r)-2.2*sin(\th r)},{1.9*sin(\th r)+2.2*cos(\th r)}) node[right] {$\vec{v}$};
% dashed parallelogram
\draw[dashed,popmuted] (P) ++({1.9*cos(\th r)},{1.9*sin(\th r)}) -- ++({-2.2*sin(\th r)},{2.2*cos(\th r)});
\draw[dashed,popmuted] (P) ++({-2.2*sin(\th r)},{2.2*cos(\th r)}) -- ++({1.9*cos(\th r)},{1.9*sin(\th r)});
% radius line
\draw[popmuted,dashed] (0,0) -- (P) node[midway,sloped,above] {\small $r$};
% theta arc
\draw[popgold,thick,->] (1.6,0) arc (0:{4.0*180/3.14159}:1.6);
\node[popgold] at (1.1,1.35) {$\theta$};
\end{tikzpicture}
\end{center}
\legende{Velocity at a point of a spiral path decomposed into its radial part $\dot{r}\,\vec{e}_r$ and transverse part $r\dot{\theta}\,\vec{e}_\theta$. The total velocity $\vec{v}$ is tangent to the path.}
\end{popfigure}

\courssub{Acceleration in polar coordinates}

\begin{propriete}[Acceleration in polar components — proof examinable]
If a particle has polar coordinates $(r,\theta)$, then its acceleration is
\[
\vec{a}=\bigl(\ddot{r}-r\dot{\theta}^{2}\bigr)\vec{e}_r+\bigl(r\ddot{\theta}+2\dot{r}\dot{\theta}\bigr)\vec{e}_\theta .
\]
The \cle{radial component} is $a_r=\ddot{r}-r\dot{\theta}^2$ and the \cle{transverse component} is $a_\theta=r\ddot{\theta}+2\dot{r}\dot{\theta}$.
\end{propriete}

\textbf{Derivation (to be reproduced in the examination).} Differentiate the velocity $\vec{v}=\dot{r}\,\vec{e}_r+r\dot{\theta}\,\vec{e}_\theta$ term by term, using the product rule on each product and the basis rules $\dfrac{d\vec{e}_r}{dt}=\dot{\theta}\vec{e}_\theta$, $\dfrac{d\vec{e}_\theta}{dt}=-\dot{\theta}\vec{e}_r$:
\begin{align*}
\vec{a}=\frac{d\vec{v}}{dt}
&=\frac{d}{dt}\bigl(\dot{r}\,\vec{e}_r\bigr)+\frac{d}{dt}\bigl(r\dot{\theta}\,\vec{e}_\theta\bigr)\\[2pt]
&=\Bigl(\ddot{r}\,\vec{e}_r+\dot{r}\,\frac{d\vec{e}_r}{dt}\Bigr)
+\Bigl(\dot{r}\dot{\theta}\,\vec{e}_\theta+r\ddot{\theta}\,\vec{e}_\theta+r\dot{\theta}\,\frac{d\vec{e}_\theta}{dt}\Bigr)\\[2pt]
&=\ddot{r}\,\vec{e}_r+\dot{r}\dot{\theta}\,\vec{e}_\theta
+\dot{r}\dot{\theta}\,\vec{e}_\theta+r\ddot{\theta}\,\vec{e}_\theta-r\dot{\theta}^{2}\,\vec{e}_r\\[2pt]
&=\bigl(\ddot{r}-r\dot{\theta}^{2}\bigr)\vec{e}_r+\bigl(r\ddot{\theta}+2\dot{r}\dot{\theta}\bigr)\vec{e}_\theta . \qquad\blacksquare
\end{align*}
Notice where each term comes from: the two copies of $\dot{r}\dot{\theta}$ arise from differentiating $\vec{e}_r$ in the first bracket and differentiating $r$ in the second bracket; the term $-r\dot{\theta}^2$ arises from differentiating $\vec{e}_\theta$.

\textbf{Identifying the physical terms.}
\begin{itemize}
\item $-r\dot{\theta}^{2}$ is the \cle{centripetal term}: even if $r$ is constant (circular motion), the particle accelerates towards the pole with magnitude $r\dot{\theta}^2=v_\theta^2/r$. The minus sign shows it points along $-\vec{e}_r$, i.e.\ inwards.
\item $2\dot{r}\dot{\theta}$ is the \cle{Coriolis-type term}: it appears only when the particle moves radially ($\dot{r}\neq0$) \emph{and} the radius rotates ($\dot{\theta}\neq0$). It is the extra sideways acceleration felt when sliding along a rotating arm — the same effect that deflects winds on the rotating Earth.
\end{itemize}

\begin{attention}[The radial acceleration is NOT simply $\ddot{r}$]
The single most common error at the GCE is writing $a_r=\ddot{r}$. The centripetal term $-r\dot{\theta}^2$ must \textbf{never} be forgotten: a particle on a circle of fixed radius has $\ddot{r}=0$ yet certainly accelerates (towards the centre). Likewise $a_\theta\neq r\ddot{\theta}$ in general — the Coriolis term $2\dot{r}\dot{\theta}$ must be included whenever $\dot r\neq0$.
\end{attention}

\begin{popfigure}
\begin{center}
\begin{tikzpicture}[scale=1.15]
\draw[->,popmuted] (-0.5,0) -- (6.5,0) node[below left] {\small initial line};
\draw[->,popmuted] (0,-0.5) -- (0,4.2);
\fill (0,0) circle (1.5pt) node[below left] {$O$};
% curved trajectory
\draw[popblue,thick,domain=0.4:5.4,samples=120,smooth] plot ({0.55*\x*cos(\x r)},{0.55*\x*sin(\x r)});
\node[popblue] at (4.9,1.6) {\small path};
\pgfmathsetmacro{\th}{3.2}
\pgfmathsetmacro{\rr}{1.76}
\coordinate (P) at ({\rr*cos(\th r)},{\rr*sin(\th r)});
\fill[popdark] (P) circle (2pt) node[below left] {$P$};
% basis
\draw[->,popgreen,thick] (P) -- ++({1.2*cos(\th r)},{1.2*sin(\th r)}) node[right] {$\vec{e}_r$};
\draw[->,popgreen,thick] (P) -- ++({-1.2*sin(\th r)},{1.2*cos(\th r)}) node[above] {$\vec{e}_\theta$};
% a_r (negative: inward)
\draw[->,poporange,very thick] (P) -- ++({-1.7*cos(\th r)},{-1.7*sin(\th r)}) node[below left] {$a_r\,\vec{e}_r$};
% a_theta
\draw[->,poppurple,very thick] (P) -- ++({-1.5*sin(\th r)},{1.5*cos(\th r)}) node[above left] {$a_\theta\,\vec{e}_\theta$};
% total a
\draw[->,popink,ultra thick] (P) -- ++({-1.7*cos(\th r)-1.5*sin(\th r)},{-1.7*sin(\th r)+1.5*cos(\th r)}) node[below] {$\vec{a}$};
\draw[dashed,popmuted] (P) ++({-1.7*cos(\th r)},{-1.7*sin(\th r)}) -- ++({-1.5*sin(\th r)},{1.5*cos(\th r)});
\draw[dashed,popmuted] (P) ++({-1.5*sin(\th r)},{1.5*cos(\th r)}) -- ++({-1.7*cos(\th r)},{-1.7*sin(\th r)});
\draw[popmuted,dashed] (0,0) -- (P);
\end{tikzpicture}
\end{center}
\legende{Acceleration decomposed into radial and transverse parts. Here $a_r=\ddot r-r\dot\theta^2<0$ (dominated by the centripetal term), so the radial part points inwards; $\vec{a}$ points to the concave side of the path, as it always must.}
\end{popfigure}

\courssub{An alternative form of the transverse acceleration}

\begin{propriete}[Transverse acceleration and the quantity $r^2\dot{\theta}$]
The transverse component of acceleration can be written as
\[
a_\theta=r\ddot{\theta}+2\dot{r}\dot{\theta}=\frac{1}{r}\frac{d}{dt}\bigl(r^{2}\dot{\theta}\bigr).
\]
Consequently, $a_\theta=0$ if and only if $r^{2}\dot{\theta}$ is \emph{constant}.
\end{propriete}

\textbf{Proof.} By the product rule,
\[
\frac{d}{dt}\bigl(r^{2}\dot{\theta}\bigr)=2r\dot{r}\dot{\theta}+r^{2}\ddot{\theta}
=r\bigl(2\dot{r}\dot{\theta}+r\ddot{\theta}\bigr)=r\,a_\theta,
\]
and dividing by $r$ gives the result. $\blacksquare$

This compact form is extremely useful: whenever a force is purely radial (a \emph{central force}, such as gravity acting on a satellite), the transverse equation of motion $m a_\theta=0$ immediately yields $r^2\dot{\theta}=\text{constant}$ — the conservation of angular momentum, which is the mathematical content of Kepler's second law.

\courssub{Method: derive, don't just quote}

\begin{methode}[Obtaining the polar formulae from $\vec{r}=r\vec{e}_r$]
The GCE regularly asks candidates to \emph{obtain} these formulae. The expected presentation is:
\begin{enumerate}
\item Write $\vec{r}=r\,\vec{e}_r$ and state the basis rules $\dot{\vec{e}}_r=\dot{\theta}\vec{e}_\theta$, $\dot{\vec{e}}_\theta=-\dot{\theta}\vec{e}_r$.
\item Differentiate once with the product rule to get $\vec{v}=\dot{r}\vec{e}_r+r\dot{\theta}\vec{e}_\theta$.
\item Differentiate again, applying the product rule to \emph{each} of the two terms and substituting the basis derivatives.
\item Collect the $\vec{e}_r$ terms and the $\vec{e}_\theta$ terms separately.
\item Box the final result and, if asked, identify the centripetal and Coriolis-type terms.
\end{enumerate}
Quoting the final formula without working, when the question says ``show that'' or ``derive'', earns little or no credit.
\end{methode}

\courssub{Consistency check with Cartesian coordinates}

A good mathematician always verifies a new tool against an old one. Let us check that the polar formulae reproduce the Cartesian results on a concrete motion.

\begin{exoresolu}[Spiral motion: $r=kt$, $\theta=\omega t$]
A particle moves so that its polar coordinates are $r=kt$ and $\theta=\omega t$, where $k$ and $\omega$ are positive constants (an Archimedean spiral). (a) Find the polar components of $\vec v$ and $\vec a$. (b) Verify the acceleration by computing it in Cartesian coordinates and resolving along $\vec e_r,\vec e_\theta$.
\tcblower
\textbf{(a) Polar route.} We have $\dot r=k$, $\ddot r=0$, $\dot\theta=\omega$, $\ddot\theta=0$. Hence
\[
\vec v=k\,\vec e_r+kt\,\omega\,\vec e_\theta=k\,\vec e_r+k\omega t\,\vec e_\theta,
\]
\[
a_r=\ddot r-r\dot\theta^2=0-kt\,\omega^2=-k\omega^2 t,
\qquad
a_\theta=r\ddot\theta+2\dot r\dot\theta=0+2k\omega=2k\omega .
\]
So $\vec a=-k\omega^2 t\,\vec e_r+2k\omega\,\vec e_\theta$.

\textbf{(b) Cartesian check.} The position is
\[
x=r\cos\theta=kt\cos\omega t,\qquad y=kt\sin\omega t .
\]
Differentiating twice with the product rule:
\begin{align*}
\dot x&=k\cos\omega t-k\omega t\sin\omega t, &
\dot y&=k\sin\omega t+k\omega t\cos\omega t,\\
\ddot x&=-2k\omega\sin\omega t-k\omega^2 t\cos\omega t, &
\ddot y&=2k\omega\cos\omega t-k\omega^2 t\sin\omega t .
\end{align*}
Now resolve $\vec a=\ddot x\,\vec i+\ddot y\,\vec j$ onto the rotating basis using $\vec e_r=\cos\theta\,\vec i+\sin\theta\,\vec j$ and $\vec e_\theta=-\sin\theta\,\vec i+\cos\theta\,\vec j$:
\begin{align*}
a_r&=\ddot x\cos\omega t+\ddot y\sin\omega t\\
&=\bigl(-2k\omega\sin\omega t-k\omega^2 t\cos\omega t\bigr)\cos\omega t
+\bigl(2k\omega\cos\omega t-k\omega^2 t\sin\omega t\bigr)\sin\omega t\\
&=-k\omega^2 t\bigl(\cos^2\omega t+\sin^2\omega t\bigr)=-k\omega^2 t,\\[3pt]
a_\theta&=-\ddot x\sin\omega t+\ddot y\cos\omega t\\
&=2k\omega\bigl(\sin^2\omega t+\cos^2\omega t\bigr)=2k\omega .
\end{align*}
Both components agree exactly with part (a). $\blacksquare$ Note how much \emph{shorter} the polar route is: this is precisely why the rotating basis exists.
\end{exoresolu}

\begin{exemplebox}[Speed on the spiral]
For the same spiral, the speed is
\[
|\vec v|=\sqrt{\dot r^{2}+r^{2}\dot\theta^{2}}=\sqrt{k^{2}+k^{2}\omega^{2}t^{2}}=k\sqrt{1+\omega^{2}t^{2}},
\]
which grows with time even though both $\dot r$ and $\dot\theta$ are constant — the transverse motion contributes more and more as $r$ increases.
\end{exemplebox}

\begin{exoresolu}[A numerical application]
A particle moves so that $r=(2t+1)\ \mathrm{m}$ and $\theta=3t\ \mathrm{rad}$, where $t$ is in seconds. Find, at the instant $t=1\ \mathrm{s}$: (a) the radial and transverse components of the velocity, and the speed; (b) the radial and transverse components of the acceleration, and its magnitude.
\tcblower
We have $\dot r=2$, $\ddot r=0$, $\dot\theta=3$, $\ddot\theta=0$, and at $t=1$: $r=3\ \mathrm{m}$.

\textbf{(a) Velocity.}
\[
v_r=\dot r=2\ \mathrm{m\,s^{-1}},\qquad
v_\theta=r\dot\theta=3\times3=9\ \mathrm{m\,s^{-1}} .
\]
\[
|\vec v|=\sqrt{v_r^{2}+v_\theta^{2}}=\sqrt{4+81}=\sqrt{85}\approx9.22\ \mathrm{m\,s^{-1}} .
\]

\textbf{(b) Acceleration.}
\[
a_r=\ddot r-r\dot\theta^{2}=0-3\times9=-27\ \mathrm{m\,s^{-2}},
\qquad
a_\theta=r\ddot\theta+2\dot r\dot\theta=0+2\times2\times3=12\ \mathrm{m\,s^{-2}} .
\]
\[
|\vec a|=\sqrt{a_r^{2}+a_\theta^{2}}=\sqrt{729+144}=\sqrt{873}\approx29.5\ \mathrm{m\,s^{-2}} .
\]
The negative sign of $a_r$ means the radial part of the acceleration points \emph{towards} the pole, dominated by the centripetal term $-r\dot\theta^2$.
\end{exoresolu}

\begin{exercice}[Circular motion as a special case]
A particle moves on a circle of radius $R$ with $\theta=\omega t$, $\omega$ constant. Starting from the polar formulae, show that $\vec v=R\omega\,\vec e_\theta$ and $\vec a=-R\omega^{2}\,\vec e_r$, and state the physical meaning of each.
\end{exercice}

\begin{corrige}[Exercise 1]
Here $r=R$ (constant), so $\dot r=\ddot r=0$, and $\dot\theta=\omega$, $\ddot\theta=0$. Then
\[
\vec v=\dot r\,\vec e_r+r\dot\theta\,\vec e_\theta=R\omega\,\vec e_\theta,
\qquad
\vec a=\bigl(0-R\omega^{2}\bigr)\vec e_r+\bigl(0+0\bigr)\vec e_\theta=-R\omega^{2}\vec e_r .
\]
The velocity is purely transverse (tangent to the circle) with constant magnitude $R\omega$; the acceleration is purely radial, directed \emph{towards the centre} with magnitude $R\omega^2=v^2/R$: this is the familiar centripetal acceleration of uniform circular motion, recovered as a special case of the general polar formulae.
\end{corrige}

\begin{exercice}[A rotating spoke]
An ant walks outwards along a straight spoke of a record turntable at a constant speed of $0.5\ \mathrm{m\,s^{-1}}$ relative to the spoke, starting from the centre at $t=0$. The turntable rotates at a constant angular speed of $2\ \mathrm{rad\,s^{-1}}$. Taking the centre as pole, write down $r(t)$ and $\theta(t)$, and find the polar components of the ant's velocity and acceleration at $t=2\ \mathrm{s}$.
\end{exercice}

\begin{corrige}[Exercise 2]
We have $r=0.5t$ and $\theta=2t$, so $\dot r=0.5$, $\ddot r=0$, $\dot\theta=2$, $\ddot\theta=0$. At $t=2\ \mathrm{s}$: $r=1\ \mathrm{m}$.
\[
v_r=\dot r=0.5\ \mathrm{m\,s^{-1}},\qquad v_\theta=r\dot\theta=1\times2=2\ \mathrm{m\,s^{-1}};
\]
\[
a_r=\ddot r-r\dot\theta^{2}=0-1\times4=-4\ \mathrm{m\,s^{-2}},\qquad
a_\theta=r\ddot\theta+2\dot r\dot\theta=0+2\times0.5\times2=2\ \mathrm{m\,s^{-2}} .
\]
Although the ant moves ``steadily'' in both $r$ and $\theta$, it has a genuine acceleration: an inward centripetal part of $4\ \mathrm{m\,s^{-2}}$ and a transverse Coriolis-type part of $2\ \mathrm{m\,s^{-2}}$.
\end{corrige}

\begin{aretenir}[Key results of this section]
\begin{itemize}
\item Velocity: $\vec v=\dot r\,\vec e_r+r\dot\theta\,\vec e_\theta$; speed $|\vec v|=\sqrt{\dot r^{2}+r^{2}\dot\theta^{2}}$.
\item Acceleration: $\vec a=\bigl(\ddot r-r\dot\theta^{2}\bigr)\vec e_r+\bigl(r\ddot\theta+2\dot r\dot\theta\bigr)\vec e_\theta$.
\item $-r\dot\theta^2$ is the centripetal term; $2\dot r\dot\theta$ is the Coriolis-type term.
\item $a_\theta=\dfrac{1}{r}\dfrac{d}{dt}\bigl(r^{2}\dot\theta\bigr)$, so $a_\theta=0\iff r^{2}\dot\theta=\text{constant}$.
\item Everything follows from $\vec r=r\vec e_r$ plus $\dot{\vec e}_r=\dot\theta\vec e_\theta$, $\dot{\vec e}_\theta=-\dot\theta\vec e_r$ — be ready to reproduce the derivation.
\end{itemize}
\end{aretenir}

\begin{savaistu}[Why the Coriolis term matters in Cameroon]
The term $2\dot r\dot\theta$ is the planar ancestor of the \emph{Coriolis acceleration} $2\vec\omega\wedge\vec v$ on the rotating Earth. It is why large-scale winds and ocean currents are deflected sideways, and why tropical storms spin. Near the equator — where Cameroon lies — the effect is weakest, which is one reason cyclones essentially never form within a few degrees of the equator: there is not enough Coriolis deflection to start the rotation.
\end{savaistu}

\coursec{Simple cases of radial-transverse motion}

In the previous section we derived the general expressions for velocity and acceleration in polar coordinates:
\[
\vec{v} = \dot{r}\,\vec{e}_r + r\dot{\theta}\,\vec{e}_\theta,
\qquad
\vec{a} = (\ddot{r} - r\dot{\theta}^2)\,\vec{e}_r + (r\ddot{\theta} + 2\dot{r}\dot{\theta})\,\vec{e}_\theta.
\]
These formulae are the master keys for any planar motion described by $r(t)$ and $\theta(t)$.  In this section we examine the most common special cases that appear in GCE Advanced Level examinations.  We will see how the general formulae simplify, what each term represents physically, and how to solve complete problems step by step.

\courssub{Uniform circular motion}

When a particle moves on a circle of fixed radius $r = R$ with constant angular speed $\dot{\theta} = \omega$, the motion is called \cle{uniform circular motion}.  Then $\dot{r}=0$, $\ddot{r}=0$, $\ddot{\theta}=0$ and the general formulae reduce to
\[
\vec{v} = R\omega\,\vec{e}_\theta,
\qquad
\vec{a} = -R\omega^2\,\vec{e}_r.
\]
The velocity is purely transverse, tangent to the circle, with constant magnitude $v = R\omega$.  The acceleration is purely radial, directed towards the centre of the circle, with magnitude $a = R\omega^2 = v^2/R$.  This is the \cle{centripetal acceleration}.

\begin{propriete}[Uniform circular motion]
In uniform circular motion the speed is constant but the velocity is not constant because its direction changes continuously.  The acceleration is purely centripetal: $\vec{a} = -\omega^2\vec{r}$ (or $a = v^2/R$ towards the centre).  This result is examinable and must be derived from the polar formulae by setting $r = \text{constant}$, $\dot{\theta} = \text{constant}$.
\end{propriete}

\begin{attention}[Constant speed $\neq$ zero acceleration]
A common mistake is to say ``speed is constant, therefore acceleration is zero''.  In uniform circular motion the speed is constant but the velocity vector rotates, so there is a non-zero acceleration directed radially inward.  Always check whether the direction of velocity changes.
\end{attention}

\begin{methode}[Circular motion from the general formulae]
To analyse any circular motion:
\begin{enumerate}
\item Set $r = R$ (constant) in the expressions for $\vec{v}$ and $\vec{a}$ \emph{before} differentiating.
\item Then $\dot{r}=0$, $\ddot{r}=0$ and the formulae become $\vec{v} = R\dot{\theta}\,\vec{e}_\theta$, $\vec{a} = -R\dot{\theta}^2\,\vec{e}_r + R\ddot{\theta}\,\vec{e}_\theta$.
\item Substitute the given $\dot{\theta}(t)$ and $\ddot{\theta}(t)$.
\end{enumerate}
\end{methode}

\begin{exemplebox}[Uniform circular motion on a roundabout]
A child sits on a roundabout of radius $R = 2.5\ \mathrm{m}$ rotating at a constant angular speed $\omega = 1.2\ \mathrm{rad\,s^{-1}}$.  Find the child's velocity and acceleration vectors, their magnitudes and directions.
\end{exemplebox}

\begin{exoresolu}[Uniform circular motion on a roundabout]
\textbf{Statement:} A child sits on a roundabout of radius $R = 2.5\ \mathrm{m}$ rotating at a constant angular speed $\omega = 1.2\ \mathrm{rad\,s^{-1}}$.  Find the child's velocity and acceleration vectors, their magnitudes and directions.
\tcblower
\textbf{Solution:}
\begin{itemize}
\item Given: $r = R = 2.5\ \mathrm{m}$, $\dot{\theta} = \omega = 1.2\ \mathrm{rad\,s^{-1}}$, $\ddot{\theta}=0$, $\dot{r}=0$, $\ddot{r}=0$.
\item Velocity: $\vec{v} = R\omega\,\vec{e}_\theta = 2.5 \times 1.2\,\vec{e}_\theta = 3.0\,\vec{e}_\theta\ \mathrm{m\,s^{-1}}$.  Magnitude $v = 3.0\ \mathrm{m\,s^{-1}}$, direction tangent to the circle (transverse).
\item Acceleration: $\vec{a} = -R\omega^2\,\vec{e}_r = -2.5 \times (1.2)^2\,\vec{e}_r = -3.6\,\vec{e}_r\ \mathrm{m\,s^{-2}}$.  Magnitude $a = 3.6\ \mathrm{m\,s^{-2}}$, direction radially inward (centripetal).
\end{itemize}
\end{exoresolu}

\begin{popfigure}
\begin{tikzpicture}[scale=1.2]
\draw[popblue, thick] (0,0) circle (2);
\fill[popdark] (0,0) circle (2pt) node[below left] {$O$};
\coordinate (P) at (60:2);
\fill[poporange] (P) circle (3pt) node[above right] {$P$};
\draw[->, popmuted] (0,0) -- (P) node[midway, above left] {$R$};
\draw[->, popgreen, thick] (P) -- ++({-sin(60)*1.5},{cos(60)*1.5}) node[midway, above right] {$\vec{v}=R\omega\,\vec{e}_\theta$};
\draw[->, poppink, thick] (P) -- ++({-cos(60)*1.2},{-sin(60)*1.2}) node[midway, below left] {$\vec{a}=-R\omega^2\,\vec{e}_r$};
\draw[->, popblue] (P) -- ++({cos(60)*0.6},{sin(60)*0.6}) node[above right] {$\vec{e}_r$};
\draw[->, popblue] (P) -- ++({-sin(60)*0.6},{cos(60)*0.6}) node[above left] {$\vec{e}_\theta$};
\draw[popmuted] (0.5,0) arc (0:60:0.5) node[midway, right] {$\theta$};
\end{tikzpicture}
\legende{Uniform circular motion: velocity tangent to the circle, centripetal acceleration pointing to the centre.}
\end{popfigure}

\courssub{Non-uniform circular motion}

If the radius remains constant ($r=R$) but the angular speed $\dot{\theta}$ varies with time, the motion is \cle{non-uniform circular motion}.  The general formulae with $\dot{r}=\ddot{r}=0$ give
\[
\vec{v} = R\dot{\theta}\,\vec{e}_\theta,
\qquad
\vec{a} = -R\dot{\theta}^2\,\vec{e}_r + R\ddot{\theta}\,\vec{e}_\theta.
\]
The radial component $-R\dot{\theta}^2$ is still the centripetal acceleration (now varying in magnitude because $\dot{\theta}$ varies).  The transverse component $R\ddot{\theta}$ is the \cle{tangential acceleration}, responsible for the change in speed.

\begin{propriete}[Non-uniform circular motion]
For a particle constrained to a circle of radius $R$:
\begin{itemize}
\item Radial acceleration: $a_r = -R\dot{\theta}^2$ (always directed towards the centre).
\item Transverse (tangential) acceleration: $a_\theta = R\ddot{\theta}$ (in the direction of increasing $\theta$ if $\ddot{\theta}>0$).
\item Speed: $v = R|\dot{\theta}|$; rate of change of speed: $\frac{dv}{dt} = R\ddot{\theta}$ (sign gives increase/decrease).
\end{itemize}
\end{propriete}

\begin{exemplebox}[Particle on a rotating turntable with angular acceleration]
A particle is fixed on a turntable of radius $R = 0.8\ \mathrm{m}$.  The turntable starts from rest and has a constant angular acceleration $\alpha = 0.5\ \mathrm{rad\,s^{-2}}$.  Find the velocity and acceleration vectors at $t = 4\ \mathrm{s}$.
\end{exemplebox}

\begin{exoresolu}[Particle on a rotating turntable with angular acceleration]
\textbf{Statement:} A particle is fixed on a turntable of radius $R = 0.8\ \mathrm{m}$.  The turntable starts from rest and has a constant angular acceleration $\alpha = 0.5\ \mathrm{rad\,s^{-2}}$.  Find the velocity and acceleration vectors at $t = 4\ \mathrm{s}$.
\tcblower
\textbf{Solution:}
\begin{itemize}
\item Given: $r = R = 0.8\ \mathrm{m}$, $\dot{r}=0$, $\ddot{r}=0$, $\ddot{\theta} = \alpha = 0.5\ \mathrm{rad\,s^{-2}}$.
\item Since the turntable starts from rest with constant angular acceleration, $\dot{\theta}(t) = \alpha t$.  At $t = 4\ \mathrm{s}$: $\dot{\theta} = 0.5 \times 4 = 2\ \mathrm{rad\,s^{-1}}$.
\item Velocity: $\vec{v} = R\dot{\theta}\,\vec{e}_\theta = 0.8 \times 2\,\vec{e}_\theta = 1.6\,\vec{e}_\theta\ \mathrm{m\,s^{-1}}$.  Magnitude $v = 1.6\ \mathrm{m\,s^{-1}}$, tangent to the circle.
\item Acceleration:
\[
\vec{a} = -R\dot{\theta}^2\,\vec{e}_r + R\ddot{\theta}\,\vec{e}_\theta
= -0.8 \times (2)^2\,\vec{e}_r + 0.8 \times 0.5\,\vec{e}_\theta
= -3.2\,\vec{e}_r + 0.4\,\vec{e}_\theta\ \mathrm{m\,s^{-2}}.
\]
\item Magnitude: $|\vec{a}| = \sqrt{(-3.2)^2 + (0.4)^2} = \sqrt{10.24 + 0.16} = \sqrt{10.4} \approx 3.22\ \mathrm{m\,s^{-2}}$.
\item The radial component $-3.2\ \mathrm{m\,s^{-2}}$ is the centripetal acceleration; the transverse component $0.4\ \mathrm{m\,s^{-2}}$ is the tangential acceleration that increases the speed.
\end{itemize}
\end{exoresolu}

\courssub{Motion with constant angular velocity but varying radius}

Consider a particle sliding along a rod that rotates with constant angular velocity $\dot{\theta} = \omega$ (e.g. a bead on a rotating wire, or a crane extending its boom while slewing at constant rate).  Then $\ddot{\theta}=0$, but $\dot{r}$ and $\ddot{r}$ are generally non-zero.  The acceleration becomes
\[
\vec{a} = (\ddot{r} - r\omega^2)\,\vec{e}_r + 2\dot{r}\omega\,\vec{e}_\theta.
\]
The transverse term $2\dot{r}\omega$ is the \cle{Coriolis acceleration} (in the rotating frame it appears as a fictitious force; in the inertial frame it is a real component of acceleration arising from the change of direction of the radial velocity).  It is perpendicular to the rod and proportional to the radial speed.

\begin{propriete}[Constant $\omega$, varying $r$]
If $\dot{\theta} = \omega = \text{constant}$ and $\ddot{\theta}=0$, then
\[
\vec{v} = \dot{r}\,\vec{e}_r + r\omega\,\vec{e}_\theta,
\qquad
\vec{a} = (\ddot{r} - r\omega^2)\,\vec{e}_r + 2\dot{r}\omega\,\vec{e}_\theta.
\]
The term $2\dot{r}\omega$ is the Coriolis acceleration.  It vanishes if the particle is momentarily at rest radially ($\dot{r}=0$).
\end{propriete}

\begin{attention}[Sign of the Coriolis term]
The Coriolis acceleration $2\dot{r}\omega\,\vec{e}_\theta$ is in the direction of $\vec{e}_\theta$ (i.e. the direction of rotation) if $\dot{r}>0$ (particle moving outward), and opposite to $\vec{e}_\theta$ if $\dot{r}<0$ (particle moving inward).  Do not confuse it with the centripetal term $-r\omega^2\vec{e}_r$ which is always radially inward.
\end{attention}

\begin{exemplebox}[Bead on a rotating horizontal rod]
A smooth horizontal rod rotates about a vertical axis through one end with constant angular velocity $\omega = 3\ \mathrm{rad\,s^{-1}}$.  A bead of mass $m$ is free to slide on the rod.  At a certain instant the bead is at distance $r = 0.4\ \mathrm{m}$ from the axis and has radial velocity $\dot{r} = 0.5\ \mathrm{m\,s^{-1}}$ outward and radial acceleration $\ddot{r} = 0.2\ \mathrm{m\,s^{-2}}$.  Find the velocity and acceleration of the bead at that instant.
\end{exemplebox}

\begin{exoresolu}[Bead on a rotating horizontal rod]
\textbf{Statement:} A smooth horizontal rod rotates about a vertical axis through one end with constant angular velocity $\omega = 3\ \mathrm{rad\,s^{-1}}$.  A bead of mass $m$ is free to slide on the rod.  At a certain instant the bead is at distance $r = 0.4\ \mathrm{m}$ from the axis and has radial velocity $\dot{r} = 0.5\ \mathrm{m\,s^{-1}}$ outward and radial acceleration $\ddot{r} = 0.2\ \mathrm{m\,s^{-2}}$.  Find the velocity and acceleration of the bead at that instant.
\tcblower
\textbf{Solution:}
\begin{itemize}
\item Given: $\omega = 3$, $\dot{r}=0.5$, $\ddot{r}=0.2$, $r=0.4$.
\item Velocity: $\vec{v} = \dot{r}\,\vec{e}_r + r\omega\,\vec{e}_\theta = 0.5\,\vec{e}_r + 0.4\times3\,\vec{e}_\theta = 0.5\,\vec{e}_r + 1.2\,\vec{e}_\theta\ \mathrm{m\,s^{-1}}$.
\item Acceleration:
\[
\vec{a} = (\ddot{r} - r\omega^2)\,\vec{e}_r + 2\dot{r}\omega\,\vec{e}_\theta
= (0.2 - 0.4\times9)\,\vec{e}_r + 2\times0.5\times3\,\vec{e}_\theta
= (0.2 - 3.6)\,\vec{e}_r + 3.0\,\vec{e}_\theta
= -3.4\,\vec{e}_r + 3.0\,\vec{e}_\theta\ \mathrm{m\,s^{-2}}.
\]
\item Magnitude of acceleration: $|\vec{a}| = \sqrt{(-3.4)^2 + 3.0^2} = \sqrt{11.56+9} = \sqrt{20.56} \approx 4.53\ \mathrm{m\,s^{-2}}$.
\item The radial component is inward (centripetal dominates), the transverse component is in the direction of rotation (Coriolis).
\end{itemize}
\end{exoresolu}

\begin{popfigure}
\begin{tikzpicture}[scale=1.3]
\draw[popdark, thick] (0,0) -- (3,0);
\fill[popdark] (0,0) circle (2pt) node[below left] {$O$};
\coordinate (B) at (1.5,0);
\fill[poporange] (B) circle (3pt) node[above] {$P$};
\draw[->, popblue, thick] (0,1) arc (90:180:1) node[midway, left] {$\omega$};
\draw[->, popgreen, thick] (B) -- ++(0.8,0) node[midway, above] {$\dot{r}\,\vec{e}_r$};
\draw[->, poppurple, thick] (B) -- ++(0,0.7) node[midway, right] {$r\omega\,\vec{e}_\theta$};
\draw[->, poppink, thick, dashed] (B) -- ++(0,0.9) node[midway, right] {$2\dot{r}\omega\,\vec{e}_\theta$};
\draw[->, poporange, thick, dashed] (B) -- ++(-1.2,0) node[midway, below] {$-r\omega^2\,\vec{e}_r$};
\draw[->, popgold, thick, dashed] (B) -- ++(0.5,0) node[midway, above] {$\ddot{r}\,\vec{e}_r$};
\draw[->, popblue] (B) -- ++(0.4,0) node[above right] {$\vec{e}_r$};
\draw[->, popblue] (B) -- ++(0,0.4) node[above left] {$\vec{e}_\theta$};
\end{tikzpicture}
\legende{Particle on a rotating rod: radial velocity $\dot{r}$, transverse velocity $r\omega$, Coriolis acceleration $2\dot{r}\omega$, centripetal acceleration $-r\omega^2$.}
\end{popfigure}

\courssub{General motion given by $r = f(t)$, $\theta = g(t)$ or $r = f(\theta)$}

In many problems the motion is defined explicitly by functions $r(t)$ and $\theta(t)$, or by a trajectory equation $r = f(\theta)$ together with a time law for $\theta$ (or for $r$).  The problem-solving strategy is always the same.

\begin{methode}[Problem-solving strategy for polar motion]
\begin{enumerate}
\item Identify the given functions: $r(t)$ and $\theta(t)$, or $r(\theta)$ and $\theta(t)$ (or $r(t)$ and $\theta(r)$).
\item Compute the time derivatives needed: $\dot{r}$, $\ddot{r}$, $\dot{\theta}$, $\ddot{\theta}$.  If $r = f(\theta)$, use the chain rule: $\dot{r} = f'(\theta)\dot{\theta}$, $\ddot{r} = f''(\theta)\dot{\theta}^2 + f'(\theta)\ddot{\theta}$.
\item Substitute into the general formulae:
\[
\vec{v} = \dot{r}\,\vec{e}_r + r\dot{\theta}\,\vec{e}_\theta,
\qquad
\vec{a} = (\ddot{r} - r\dot{\theta}^2)\,\vec{e}_r + (r\ddot{\theta} + 2\dot{r}\dot{\theta})\,\vec{e}_\theta.
\]
\item Evaluate at the required instant (given $t$ or $\theta$).
\item If asked for magnitude and direction of $\vec{v}$ or $\vec{a}$, compute
\[
|\vec{v}| = \sqrt{\dot{r}^2 + (r\dot{\theta})^2},\quad
\tan\alpha_v = \frac{r\dot{\theta}}{\dot{r}} \quad(\text{angle from }\vec{e}_r),
\]
and similarly for $\vec{a}$.
\item Always state the direction in words: ``radially outward/inward'', ``transverse in the direction of increasing $\theta$'', or give the angle relative to the radial direction.
\end{enumerate}
\end{methode}

\begin{exemplebox}[Spiral motion $r = a\theta$ with constant $\dot{\theta}$]
A particle moves along the spiral $r = 0.2\theta$ (with $r$ in metres, $\theta$ in radians) such that $\dot{\theta} = 2\ \mathrm{rad\,s^{-1}}$ constant.  Find the velocity and acceleration when $\theta = \pi/2$.
\end{exemplebox}

\begin{exoresolu}[Spiral motion $r = a\theta$ with constant $\dot{\theta}$]
\textbf{Statement:} A particle moves along the spiral $r = 0.2\theta$ (with $r$ in metres, $\theta$ in radians) such that $\dot{\theta} = 2\ \mathrm{rad\,s^{-1}}$ constant.  Find the velocity and acceleration when $\theta = \pi/2$.
\tcblower
\textbf{Solution:}
\begin{itemize}
\item Given: $r = 0.2\theta$, $\dot{\theta}=2$, $\ddot{\theta}=0$.
\item At $\theta = \pi/2$: $r = 0.2\times\pi/2 = 0.1\pi \approx 0.314\ \mathrm{m}$.
\item Derivatives: $\dot{r} = 0.2\dot{\theta} = 0.4\ \mathrm{m\,s^{-1}}$, $\ddot{r} = 0.2\ddot{\theta} = 0$.
\item Velocity:
\[
\vec{v} = \dot{r}\,\vec{e}_r + r\dot{\theta}\,\vec{e}_\theta
= 0.4\,\vec{e}_r + (0.1\pi\times2)\,\vec{e}_\theta
= 0.4\,\vec{e}_r + 0.2\pi\,\vec{e}_\theta \approx 0.4\,\vec{e}_r + 0.628\,\vec{e}_\theta\ \mathrm{m\,s^{-1}}.
\]
Magnitude $v = \sqrt{0.4^2 + (0.2\pi)^2} \approx \sqrt{0.16 + 0.394} = \sqrt{0.554} \approx 0.744\ \mathrm{m\,s^{-1}}$.
Direction: angle $\alpha$ from $\vec{e}_r$ with $\tan\alpha = \frac{0.2\pi}{0.4} = \frac{\pi}{2} \approx 1.571$, so $\alpha \approx 57.5^\circ$ towards $\vec{e}_\theta$.
\item Acceleration:
\[
\vec{a} = (\ddot{r} - r\dot{\theta}^2)\,\vec{e}_r + (r\ddot{\theta} + 2\dot{r}\dot{\theta})\,\vec{e}_\theta
= (0 - 0.1\pi\times4)\,\vec{e}_r + (0 + 2\times0.4\times2)\,\vec{e}_\theta
= -0.4\pi\,\vec{e}_r + 1.6\,\vec{e}_\theta \approx -1.257\,\vec{e}_r + 1.6\,\vec{e}_\theta\ \mathrm{m\,s^{-2}}.
\]
Magnitude $a = \sqrt{(0.4\pi)^2 + 1.6^2} \approx \sqrt{1.579 + 2.56} = \sqrt{4.139} \approx 2.034\ \mathrm{m\,s^{-2}}$.
Direction: angle $\beta$ from the inward radial direction ($-\vec{e}_r$) with $\tan\beta = \frac{1.6}{0.4\pi} = \frac{4}{\pi} \approx 1.273$, so $\beta \approx 51.9^\circ$ towards $\vec{e}_\theta$.
\end{itemize}
\end{exoresolu}

\begin{popfigure}
\begin{tikzpicture}[scale=1.2]
\begin{axis}[
    axis equal,
    axis lines=none,
    xmin=-2, xmax=2,
    ymin=-2, ymax=2,
    width=10cm,
    height=10cm
]
\addplot[popblue, thick, domain=0:12.566, samples=400, variable=\t]
    ({0.2*\t*cos(deg(\t))}, {0.2*\t*sin(deg(\t))});
\coordinate (P1) at ({0.2*1.571*cos(deg(1.571))}, {0.2*1.571*sin(deg(1.571))});
\fill[poporange] (P1) circle (3pt);
\draw[->, popgreen, thick] (P1) -- ++({0.4*cos(deg(1.571)) - 0.628*sin(deg(1.571))}, {0.4*sin(deg(1.571)) + 0.628*cos(deg(1.571))}) node[midway, above right] {$\vec{v}$};
\coordinate (P2) at ({0.2*3.142*cos(deg(3.142))}, {0.2*3.142*sin(deg(3.142))});
\fill[poporange] (P2) circle (3pt);
\draw[->, popgreen, thick] (P2) -- ++({0.4*cos(deg(3.142)) - 1.257*sin(deg(3.142))}, {0.4*sin(deg(3.142)) + 1.257*cos(deg(3.142))}) node[midway, below left] {$\vec{v}$};
\fill[popdark] (0,0) circle (2pt) node[below left] {$O$};
\end{axis}
\end{tikzpicture}
\legende{Spiral path $r = 0.2\theta$ with velocity vectors at $\theta=\pi/2$ and $\theta=\pi$.}
\end{popfigure}

\courssub{Physical interpretation in concrete situations}

The radial-transverse decomposition is not just a mathematical trick; each component has a clear physical meaning.

\begin{itemize}
\item \textbf{Radial component of velocity $\dot{r}$:} Rate of change of the distance from the origin.  Positive = moving outward, negative = moving inward.  Example: a satellite dish tracking a satellite --- $\dot{r}$ is the rate at which the distance to the satellite changes.
\item \textbf{Transverse component of velocity $r\dot{\theta}$:} Speed due to rotation about the origin.  Example: a point on a rotating crane boom --- $r\dot{\theta}$ is the speed of the tip due to slewing.
\item \textbf{Radial component of acceleration $\ddot{r} - r\dot{\theta}^2$:} Contains the radial acceleration $\ddot{r}$ (change of radial speed) minus the centripetal term $r\dot{\theta}^2$ (needed to keep the particle turning).  If the net radial acceleration is negative, the resultant force is inward.
\item \textbf{Transverse component of acceleration $r\ddot{\theta} + 2\dot{r}\dot{\theta}$:} Contains the tangential acceleration $r\ddot{\theta}$ (due to angular acceleration) and the Coriolis term $2\dot{r}\dot{\theta}$ (due to radial motion in a rotating system).  Example: a crane extending its boom while slewing --- the Coriolis term appears as a sideways force on the boom.
\end{itemize}

\begin{savaistu}[Coriolis effect in Cameroon]
The Coriolis acceleration $2\dot{r}\dot{\theta}$ is the same mathematical term that, on the rotating Earth, deflects moving air masses to the right in the northern hemisphere and to the left in the southern hemisphere.  Cameroon lies just north of the equator, so large-scale winds (like the Harmattan) are slightly deflected to the right.  In engineering, the Coriolis term must be accounted for in the design of rotating machinery (turbines, centrifuges) and in the guidance of rockets launched from a rotating platform.
\end{savaistu}

\courssub{Angular momentum and its conservation (exam-useful extra)}

For a particle of mass $m$ moving under a central force (force directed towards or away from a fixed point $O$), the transverse component of acceleration is zero: $a_\theta = 0$.  From the general formula $a_\theta = r\ddot{\theta} + 2\dot{r}\dot{\theta} = \frac{1}{r}\frac{d}{dt}(r^2\dot{\theta})$, we obtain
\[
\frac{d}{dt}(r^2\dot{\theta}) = 0 \quad\Longrightarrow\quad r^2\dot{\theta} = h = \text{constant}.
\]
The quantity $h = r^2\dot{\theta}$ is the \cle{specific angular momentum} (angular momentum per unit mass).  Its conservation is a powerful tool: it allows us to eliminate $\dot{\theta}$ in favour of $r$ and $h$, reducing the problem to a one-dimensional radial equation.

\begin{propriete}[Conservation of angular momentum]
If the transverse acceleration vanishes ($a_\theta = 0$), then $r^2\dot{\theta} = h$ is constant.  This holds for any central force (gravitational, electrostatic, tension in a string through a hole, etc.).  The areal velocity $\frac12 r^2\dot{\theta} = \frac12 h$ is also constant (Kepler's second law).
\end{propriete}

\begin{exemplebox}[Particle on a smooth horizontal table pulled by a string through a hole]
A particle of mass $m$ moves on a smooth horizontal table.  It is attached to a light inextensible string which passes through a small hole in the table.  The particle moves in a circle of radius $r_0$ with speed $v_0$.  The string is pulled downwards slowly so that the radius decreases to $r$.  Find the speed $v$ of the particle when the radius is $r$, and the tension in the string at that instant.
\end{exemplebox}

\begin{exoresolu}[Particle on a smooth horizontal table pulled by a string through a hole]
\textbf{Statement:} A particle of mass $m$ moves on a smooth horizontal table.  It is attached to a light inextensible string which passes through a small hole in the table.  The particle moves in a circle of radius $r_0$ with speed $v_0$.  The string is pulled downwards slowly so that the radius decreases to $r$.  Find the speed $v$ of the particle when the radius is $r$, and the tension in the string at that instant.
\tcblower
\textbf{Solution:}
\begin{itemize}
\item The force on the particle is the tension $T$ directed towards the hole (central force).  Hence transverse acceleration $a_\theta = 0$ and angular momentum is conserved.
\item Initial angular momentum per unit mass: $h = r_0 v_0$ (since initially $\dot{r}=0$, $v_0 = r_0\dot{\theta}_0$).
\item At radius $r$, the transverse velocity is $v_\theta = r\dot{\theta} = h/r = r_0 v_0 / r$.  The radial velocity $\dot{r}$ is not given, but the speed $v$ satisfies $v^2 = \dot{r}^2 + v_\theta^2$.  If the string is pulled ``slowly'', we may assume $\dot{r}$ is negligible compared to $v_\theta$, so $v \approx v_\theta = \frac{r_0 v_0}{r}$.
\item Tension provides the centripetal force: $T = m a_r = m(-r\dot{\theta}^2)$ (since $\ddot{r}\approx 0$ for slow pulling).  Using $\dot{\theta} = v_\theta/r = r_0 v_0 / r^2$:
\[
T = m r \left(\frac{r_0 v_0}{r^2}\right)^2 = \frac{m r_0^2 v_0^2}{r^3}.
\]
\end{itemize}
\end{exoresolu}

\courssub{Complete GCE-style worked example}

\begin{exemplebox}[GCE-style problem: motion defined by $r = 3t^2$, $\theta = 2t^3$]
A particle moves in a plane such that its polar coordinates at time $t$ are $r = 3t^2$ metres and $\theta = 2t^3$ radians ($t \ge 0$).  Find, at $t = 1\ \mathrm{s}$:
\begin{enumerate}
\item the velocity vector in terms of $\vec{e}_r$ and $\vec{e}_\theta$,
\item the acceleration vector in terms of $\vec{e}_r$ and $\vec{e}_\theta$,
\item the magnitude of the acceleration,
\item the angle between the acceleration vector and the radial direction.
\end{enumerate}
\end{exemplebox}

\begin{exoresolu}[GCE-style problem: motion defined by $r = 3t^2$, $\theta = 2t^3$]
\textbf{Statement:} A particle moves in a plane such that its polar coordinates at time $t$ are $r = 3t^2$ metres and $\theta = 2t^3$ radians ($t \ge 0$).  Find, at $t = 1\ \mathrm{s}$:
\begin{enumerate}
\item the velocity vector in terms of $\vec{e}_r$ and $\vec{e}_\theta$,
\item the acceleration vector in terms of $\vec{e}_r$ and $\vec{e}_\theta$,
\item the magnitude of the acceleration,
\item the angle between the acceleration vector and the radial direction.
\end{enumerate}
\tcblower
\textbf{Solution:}
\begin{enumerate}
\item Compute derivatives:
\[
r = 3t^2,\quad \dot{r} = 6t,\quad \ddot{r} = 6;
\]
\[
\theta = 2t^3,\quad \dot{\theta} = 6t^2,\quad \ddot{\theta} = 12t.
\]
At $t=1$: $r=3$, $\dot{r}=6$, $\ddot{r}=6$, $\dot{\theta}=6$, $\ddot{\theta}=12$.
\item Velocity:
\[
\vec{v} = \dot{r}\,\vec{e}_r + r\dot{\theta}\,\vec{e}_\theta = 6\,\vec{e}_r + 3\times6\,\vec{e}_\theta = 6\,\vec{e}_r + 18\,\vec{e}_\theta\ \mathrm{m\,s^{-1}}.
\]
\item Acceleration:
\[
\vec{a} = (\ddot{r} - r\dot{\theta}^2)\,\vec{e}_r + (r\ddot{\theta} + 2\dot{r}\dot{\theta})\,\vec{e}_\theta
= (6 - 3\times36)\,\vec{e}_r + (3\times12 + 2\times6\times6)\,\vec{e}_\theta
= (6 - 108)\,\vec{e}_r + (36 + 72)\,\vec{e}_\theta
= -102\,\vec{e}_r + 108\,\vec{e}_\theta\ \mathrm{m\,s^{-2}}.
\]
\item Magnitude of acceleration:
\[
|\vec{a}| = \sqrt{(-102)^2 + 108^2} = \sqrt{10404 + 11664} = \sqrt{22068} \approx 148.6\ \mathrm{m\,s^{-2}}.
\]
\item Angle between $\vec{a}$ and the radial direction ($\vec{e}_r$):
The radial component is negative (inward), the transverse component is positive (direction of increasing $\theta$).  The angle $\phi$ measured from the inward radial direction ($-\vec{e}_r$) towards $\vec{e}_\theta$ satisfies
\[
\tan\phi = \frac{108}{102} = \frac{18}{17} \approx 1.0588 \quad\Rightarrow\quad \phi \approx 46.6^\circ.
\]
Thus the acceleration makes an angle of about $46.6^\circ$ with the inward radial direction, towards the transverse direction.
\end{enumerate}
\end{exoresolu}

\courssub{Summary of key results}

\begin{aretenir}[Key formulae for special cases]
\begin{itemize}
\item \textbf{Uniform circular motion} ($r=R$, $\dot{\theta}=\omega$ constant):
$\vec{v}=R\omega\vec{e}_\theta$, $\vec{a}=-R\omega^2\vec{e}_r$, $a=v^2/R$.
\item \textbf{Non-uniform circular motion} ($r=R$):
$\vec{v}=R\dot{\theta}\vec{e}_\theta$, $\vec{a}=-R\dot{\theta}^2\vec{e}_r + R\ddot{\theta}\vec{e}_\theta$.
\item \textbf{Constant $\omega$, varying $r$} ($\dot{\theta}=\omega$, $\ddot{\theta}=0$):
$\vec{v}=\dot{r}\vec{e}_r + r\omega\vec{e}_\theta$, $\vec{a}=(\ddot{r}-r\omega^2)\vec{e}_r + 2\dot{r}\omega\vec{e}_\theta$.
\item \textbf{General motion} ($r(t)$, $\theta(t)$):
$\vec{v}=\dot{r}\vec{e}_r + r\dot{\theta}\vec{e}_\theta$,
$\vec{a}=(\ddot{r}-r\dot{\theta}^2)\vec{e}_r + (r\ddot{\theta}+2\dot{r}\dot{\theta})\vec{e}_\theta$.
\item \textbf{Angular momentum conservation}: if $a_\theta=0$ then $r^2\dot{\theta}=h$ constant.
\end{itemize}
\end{aretenir}

\begin{attention}[Exam tips]
\begin{itemize}
\item Always write the general formulae first, then substitute the given conditions.  This shows the examiner you know the derivation.
\item When asked for ``magnitude and direction'', give the magnitude as a number with units, and the direction in words (e.g. ``at $30^\circ$ to the inward radial direction towards the transverse direction'').
\item Do not confuse $\dot{\theta}$ (angular speed) with $\omega$ (often used for constant angular speed).  Use the notation given in the question.
\item In problems with $r = f(\theta)$, remember to use the chain rule: $\dot{r} = \frac{dr}{d\theta}\dot{\theta}$, $\ddot{r} = \frac{d^2r}{d\theta^2}\dot{\theta}^2 + \frac{dr}{d\theta}\ddot{\theta}$.
\item The Coriolis term $2\dot{r}\dot{\theta}$ appears whenever both $\dot{r}$ and $\dot{\theta}$ are non-zero.  It is a real acceleration in the inertial frame, not a fictitious force.
\end{itemize}
\end{attention}

\begin{exercice}[Practice problems]
\begin{enumerate}
\item A particle moves on the curve $r = 2\theta$ with $\dot{\theta} = 3\ \mathrm{rad\,s^{-1}}$ constant.  Find the velocity and acceleration when $\theta = \pi/3$.
\item A bead slides on a smooth wire shaped as the spiral $r = a\theta$ ($a>0$).  The wire rotates about the vertical axis through $O$ with constant angular velocity $\omega$.  If the bead is released from rest at $r = r_0$, find its radial velocity $\dot{r}$ as a function of $r$.  (Hint: use conservation of energy and angular momentum.)
\item A satellite is in an elliptical orbit described by $r = \frac{l}{1+e\cos\theta}$ with $\dot{\theta} = h/r^2$ (where $h$ is constant).  Show that the radial acceleration is $-\frac{h^2}{l r^2} = -\frac{h^2}{l^3}(1+e\cos\theta)^2$ and interpret the terms.
\end{enumerate}
\end{exercice}

\begin{corrige}[Exercise 1]
$r = 2\theta$, $\dot{\theta}=3$, $\ddot{\theta}=0$.  At $\theta=\pi/3$: $r=2\pi/3$, $\dot{r}=2\dot{\theta}=6$, $\ddot{r}=0$.
$\vec{v} = 6\vec{e}_r + (2\pi/3)\times3\vec{e}_\theta = 6\vec{e}_r + 2\pi\vec{e}_\theta$.
$\vec{a} = (0 - (2\pi/3)\times9)\vec{e}_r + (0 + 2\times6\times3)\vec{e}_\theta = -6\pi\vec{e}_r + 36\vec{e}_\theta$.
Magnitude $a = \sqrt{36\pi^2 + 1296} \approx 48.7\ \mathrm{m\,s^{-2}}$.
\end{corrige}

\begin{corrige}[Exercise 2]
Angular momentum conservation: $r^2\dot{\theta} = r_0^2\omega$ (since initially $\dot{r}=0$, $\dot{\theta}=\omega$).  But the wire forces $\dot{\theta}=\omega$, so angular momentum is not conserved because the wire exerts a transverse force.  Instead use energy: $\frac12 m(\dot{r}^2 + r^2\omega^2) = \frac12 m r_0^2\omega^2$ (initial kinetic energy).  Then $\dot{r}^2 = \omega^2(r_0^2 - r^2)$, so $\dot{r} = -\omega\sqrt{r_0^2 - r^2}$ (negative because $r$ decreases).
\end{corrige}

\begin{corrige}[Exercise 3]
Given $r = l/(1+e\cos\theta)$, $\dot{\theta}=h/r^2$.  Compute $\dot{r} = \frac{dr}{d\theta}\dot{\theta} = \frac{le\sin\theta}{(1+e\cos\theta)^2}\cdot\frac{h}{r^2} = \frac{he\sin\theta}{l}$.  Then $\ddot{r} = \frac{d}{dt}\dot{r} = \frac{he\cos\theta}{l}\dot{\theta} = \frac{he\cos\theta}{l}\cdot\frac{h}{r^2} = \frac{h^2e\cos\theta}{lr^2}$.
Radial acceleration: $a_r = \ddot{r} - r\dot{\theta}^2 = \frac{h^2e\cos\theta}{lr^2} - r\frac{h^2}{r^4} = \frac{h^2}{r^2}\left(\frac{e\cos\theta}{l} - \frac{1}{r}\right)$.
Substitute $1/r = (1+e\cos\theta)/l$: $a_r = \frac{h^2}{r^2}\left(\frac{e\cos\theta}{l} - \frac{1+e\cos\theta}{l}\right) = -\frac{h^2}{lr^2}$.
But $r = l/(1+e\cos\theta)$, so $a_r = -\frac{h^2}{l^3}(1+e\cos\theta)^2$.  This is the central acceleration for an inverse-square law force with $\mu = h^2/l$.  The term $-\frac{h^2}{l r^2}$ shows that the force is directed towards the centre and varies as the inverse square of the distance.
\end{corrige}

\newpage

\coursec{Integration activity}

\courssub{Aims of the activity}

This integration activity brings together all the skills developed in the chapter: describing a two-dimensional motion in \cle{Cartesian coordinates}, describing the same motion in \cle{polar coordinates}, and deriving and using the \cle{radial--transverse components} of velocity and acceleration. By the end of the activity, you should be able to:

\begin{itemize}
\item model a real rotating-and-sliding motion by polar equations $r(t)$ and $\theta(t)$;
\item compute the velocity $\vec{v} = \dot{r}\,\vec{e}_r + r\dot{\theta}\,\vec{e}_\theta$ and the acceleration $\vec{a} = \left(\ddot{r} - r\dot{\theta}^{2}\right)\vec{e}_r + \left(r\ddot{\theta} + 2\dot{r}\dot{\theta}\right)\vec{e}_\theta$ at a given instant;
\item recompute the same quantities in Cartesian coordinates and verify that the two descriptions agree;
\item interpret physically the radial and transverse components (centripetal and Coriolis-type effects) and draw an engineering conclusion;
\item communicate results clearly on a poster, including a diagram of the rotating basis.
\end{itemize}

\begin{experience}[Organisation]
Work in groups of three or four. Time allowed: one double period (about 90 minutes). Each group needs graph paper or a poster sheet, a ruler, a protractor, a calculator, and coloured markers. Roles: a \emph{modeller} (sets up the equations), a \emph{calculator} (performs the differentiations), a \emph{verifier} (checks with the Cartesian approach) and a \emph{reporter} (prepares the poster and presents).
\end{experience}

\courssub{Situation}

\begin{exemplebox}[The rotating crane of the Port of Douala]
The Port of Douala, on the Wouri estuary, is the main gateway for goods entering Cameroon. On the container terminal, a slewing crane unloads containers from lorries onto waiting wagons. During one operation, the crane operator performs two simultaneous motions:

\begin{itemize}
\item the \textbf{jib rotates} about a fixed vertical axis $O$ at a \textbf{constant angular velocity} $\omega = 0.2\ \mathrm{rad\,s^{-1}}$;
\item the \textbf{trolley} carrying a container of mass $m$ slides \textbf{outward along the jib} at a \textbf{constant speed} $u = 0.5\ \mathrm{m\,s^{-1}}$, measured relative to the jib.
\end{itemize}

At time $t = 0$, the trolley is at a distance $r_0 = 4\ \mathrm{m}$ from the axis $O$, and the jib points along a fixed reference direction $Ox$ (so $\theta = 0$). The motion takes place in a horizontal plane, and friction in the trolley mechanism is negligible.

The safety engineer of the terminal wants to know, at the instant $t = 10\ \mathrm{s}$:

\begin{enumerate}
\item the exact velocity and acceleration of the container;
\item whether the transverse guide cable of the trolley must mainly resist a \emph{radial} effect (pulling the container towards or away from the axis) or a \emph{transverse} effect (pushing it sideways along the direction of rotation).
\end{enumerate}

Your group is the team of junior engineers asked to produce the report.
\end{exemplebox}

\textbf{Data summary.}

\begin{center}
\begin{tabularx}{\linewidth}{|l|X|}\hline
\rowcolor{popblueL} \textbf{Quantity} & \textbf{Value} \\ \hline
Angular velocity of the jib & $\omega = \dot{\theta} = 0.2\ \mathrm{rad\,s^{-1}}$ (constant) \\ \hline
Radial speed of the trolley & $u = \dot{r} = 0.5\ \mathrm{m\,s^{-1}}$ (constant) \\ \hline
Initial radius & $r_0 = 4\ \mathrm{m}$ \\ \hline
Instant of interest & $t = 10\ \mathrm{s}$ \\ \hline
\end{tabularx}
\end{center}

\courssub{Tasks}

\textbf{Part A — Modelling the motion.}

\begin{enumerate}
\item Explain why the polar equations of motion of the container are
\[ r(t) = 4 + 0.5t \qquad \text{and} \qquad \theta(t) = 0.2t . \]
\item Compute the position of the container at $t = 10\ \mathrm{s}$: give $r$, $\theta$ (in radians and in degrees), and then the Cartesian coordinates $(x, y) = (r\cos\theta, r\sin\theta)$.
\item On your poster, draw the jib at $t = 0$ and at $t = 10\ \mathrm{s}$, mark the position of the container, and draw the rotating basis vectors $\vec{e}_r$ and $\vec{e}_\theta$ at that position.
\end{enumerate}

\textbf{Part B — Velocity and acceleration in polar (radial--transverse) components.}

\begin{enumerate}
\setcounter{enumi}{3}
\item Compute $\dot{r}$, $\ddot{r}$, $\dot{\theta}$, $\ddot{\theta}$. Which of these are zero, and why?
\item Write the velocity in the form $\vec{v} = \dot{r}\,\vec{e}_r + r\dot{\theta}\,\vec{e}_\theta$ and evaluate its components at $t = 10\ \mathrm{s}$. Deduce the speed $|\vec{v}|$.
\item Write the acceleration in the form
\[ \vec{a} = \left(\ddot{r} - r\dot{\theta}^{2}\right)\vec{e}_r + \left(r\ddot{\theta} + 2\dot{r}\dot{\theta}\right)\vec{e}_\theta \]
and evaluate its components at $t = 10\ \mathrm{s}$. Deduce the magnitude $|\vec{a}|$.
\item Identify, in the acceleration, the term due to the rotation alone (centripetal-type term) and the term due to the combination of rotation and outward sliding (Coriolis-type term $2\dot{r}\dot{\theta}$).
\end{enumerate}

\textbf{Part C — Verification in Cartesian coordinates.}

\begin{enumerate}
\setcounter{enumi}{7}
\item Write $x(t) = r(t)\cos\theta(t)$ and $y(t) = r(t)\sin\theta(t)$ explicitly, then differentiate twice to obtain $\dot{x}, \dot{y}, \ddot{x}, \ddot{y}$ at $t = 10\ \mathrm{s}$.
\item Check that $\sqrt{\dot{x}^{2} + \dot{y}^{2}}$ and $\sqrt{\ddot{x}^{2} + \ddot{y}^{2}}$ give the same values as in Part B. Why must this be the case?
\end{enumerate}

\textbf{Part D — Engineering conclusion.}

\begin{enumerate}
\setcounter{enumi}{9}
\item Compare the magnitudes of the radial and transverse components of $\vec{a}$. Which effect dominates? Advise the safety engineer on which effect the guide cable must mainly resist.
\item Discuss with the whole class: for this problem, what are the advantages of the polar description over the Cartesian one? When would Cartesian coordinates be preferable?
\end{enumerate}

\courssub{Model answer}

\textbf{Part A.}

\textbf{1.} The trolley moves outward at constant speed $u = 0.5\ \mathrm{m\,s^{-1}}$ starting from $r_0 = 4\ \mathrm{m}$, so $r(t) = r_0 + ut = 4 + 0.5t$. The jib rotates at constant angular velocity $\omega = 0.2\ \mathrm{rad\,s^{-1}}$ starting from $\theta = 0$, so $\theta(t) = \omega t = 0.2t$.

\textbf{2.} At $t = 10\ \mathrm{s}$:
\[ r = 4 + 0.5(10) = 9\ \mathrm{m}, \qquad \theta = 0.2(10) = 2\ \mathrm{rad} \approx 114.6^{\circ}. \]
Hence
\[ x = 9\cos 2 \approx 9(-0.4161) \approx -3.75\ \mathrm{m}, \qquad y = 9\sin 2 \approx 9(0.9093) \approx 8.18\ \mathrm{m}. \]

\textbf{3.} The diagram shows the jib at both instants and the rotating basis at the position of the container.

\begin{popfigure}
\begin{tikzpicture}[scale=0.42]
\draw[->, gray] (-5,0) -- (6,0) node[below left] {$Ox$};
\draw[->, gray] (0,-5) -- (0,6) node[below left] {$Oy$};
\fill (0,0) circle (3pt) node[below right] {$O$};
\draw[popmuted, thick] (0,0) -- (4,0) node[midway, below] {jib at $t=0$};
\fill[popmuted] (4,0) circle (3pt) node[below] {$t=0$};
\draw[popblue, very thick] (0,0) -- ({9*cos(114.592)}, {9*sin(114.592)}) node[midway, above left] {jib at $t=10$ s};
\fill[poporange] ({9*cos(114.592)}, {9*sin(114.592)}) circle (4pt);
\draw[->, popdark, thick] ({9*cos(114.592)}, {9*sin(114.592)}) -- ++({2.2*cos(114.592)}, {2.2*sin(114.592)}) node[above left] {$\vec{e}_r$};
\draw[->, popgreen, thick] ({9*cos(114.592)}, {9*sin(114.592)}) -- ++({-2.2*sin(114.592)}, {2.2*cos(114.592)}) node[below left] {$\vec{e}_\theta$};
\draw[->, popgold, thick] (2.2,0) arc (0:114.592:2.2) node[midway, right] {$\theta$};
\end{tikzpicture}
\legende{The crane jib at $t = 0$ and at $t = 10\ \mathrm{s}$, with the rotating basis $(\vec{e}_r, \vec{e}_\theta)$ at the container.}
\end{popfigure}

\textbf{Part B.}

\textbf{4.} Since $r(t) = 4 + 0.5t$ and $\theta(t) = 0.2t$:
\[ \dot{r} = 0.5\ \mathrm{m\,s^{-1}}, \quad \ddot{r} = 0, \qquad \dot{\theta} = 0.2\ \mathrm{rad\,s^{-1}}, \quad \ddot{\theta} = 0. \]
The second derivatives vanish because both the radial sliding and the rotation are at \emph{constant} rates.

\textbf{5.} Velocity:
\[ \vec{v} = \dot{r}\,\vec{e}_r + r\dot{\theta}\,\vec{e}_\theta = 0.5\,\vec{e}_r + (9)(0.2)\,\vec{e}_\theta = 0.5\,\vec{e}_r + 1.8\,\vec{e}_\theta \quad (\mathrm{m\,s^{-1}}). \]
Since $\vec{e}_r \perp \vec{e}_\theta$,
\[ |\vec{v}| = \sqrt{0.5^{2} + 1.8^{2}} = \sqrt{0.25 + 3.24} = \sqrt{3.49} \approx 1.87\ \mathrm{m\,s^{-1}}. \]

\textbf{6.} Acceleration:
\begin{align*}
a_r &= \ddot{r} - r\dot{\theta}^{2} = 0 - 9(0.2)^{2} = -0.36\ \mathrm{m\,s^{-2}}, \\
a_\theta &= r\ddot{\theta} + 2\dot{r}\dot{\theta} = 0 + 2(0.5)(0.2) = 0.20\ \mathrm{m\,s^{-2}}.
\end{align*}
So $\vec{a} = -0.36\,\vec{e}_r + 0.20\,\vec{e}_\theta$ and
\[ |\vec{a}| = \sqrt{0.36^{2} + 0.20^{2}} = \sqrt{0.1296 + 0.04} = \sqrt{0.1696} \approx 0.41\ \mathrm{m\,s^{-2}}. \]

\textbf{7.} The radial term $-r\dot{\theta}^{2} = -0.36\ \mathrm{m\,s^{-2}}$ is the \cle{centripetal-type} term: even a trolley fixed on the rotating jib would have this inward acceleration. The transverse term $2\dot{r}\dot{\theta} = 0.20\ \mathrm{m\,s^{-2}}$ is the \cle{Coriolis-type} term: it exists \emph{only} because the trolley slides outward while the jib rotates.

\textbf{Part C.}

\textbf{8.} With $x(t) = (4 + 0.5t)\cos(0.2t)$ and $y(t) = (4 + 0.5t)\sin(0.2t)$:
\begin{align*}
\dot{x} &= 0.5\cos(0.2t) - 0.2(4 + 0.5t)\sin(0.2t), \\
\dot{y} &= 0.5\sin(0.2t) + 0.2(4 + 0.5t)\cos(0.2t), \\
\ddot{x} &= -0.2\cos(0.2t) - 0.04(4+0.5t)\cos(0.2t) - 0.1\sin(0.2t)\ \text{(collected below)}, \\
\ddot{y} &= \text{(similarly)}.
\end{align*}
At $t = 10$ ($r = 9$, $\theta = 2$ rad, $\cos 2 \approx -0.4161$, $\sin 2 \approx 0.9093$):
\begin{align*}
\dot{x} &= 0.5(-0.4161) - 1.8(0.9093) \approx -0.208 - 1.637 = -1.845\ \mathrm{m\,s^{-1}}, \\
\dot{y} &= 0.5(0.9093) + 1.8(-0.4161) \approx 0.455 - 0.749 = -0.294\ \mathrm{m\,s^{-1}}.
\end{align*}
Differentiating again and simplifying:
\begin{align*}
\ddot{x} &= -0.2\sin(0.2t) - 0.2\sin(0.2t)\cdot\tfrac{1}{1} - 0.04(4+0.5t)\cos(0.2t) \\
&= -0.2\sin(0.2t) - \big[0.1\sin(0.2t) + 0.04(4+0.5t)\cos(0.2t)\big],
\end{align*}
which at $t = 10$ gives
\[ \ddot{x} = -0.2(0.9093) - 0.36(-0.4161) \approx -0.182 + 0.150 = -0.032\ \mathrm{m\,s^{-2}}, \]
\[ \ddot{y} = 0.2(-0.4161) - 0.36(0.9093) \approx -0.083 - 0.327 = -0.410\ \mathrm{m\,s^{-2}}. \]

\textbf{9.} Check of magnitudes:
\[ \sqrt{\dot{x}^{2} + \dot{y}^{2}} = \sqrt{1.845^{2} + 0.294^{2}} \approx \sqrt{3.404 + 0.086} \approx 1.87\ \mathrm{m\,s^{-1}}, \]
\[ \sqrt{\ddot{x}^{2} + \ddot{y}^{2}} = \sqrt{0.032^{2} + 0.410^{2}} \approx 0.41\ \mathrm{m\,s^{-2}}. \]
Both agree with Part B. This \emph{must} be so: the velocity and acceleration are \cle{vectors}, intrinsic objects; their \emph{components} depend on the basis chosen, but their \emph{magnitudes} do not.

\begin{attention}[A frequent error in Part C]
When differentiating $x(t) = (4 + 0.5t)\cos(0.2t)$, students often forget that \emph{both} factors depend on $t$: the product rule is compulsory, and the chain rule brings the factor $0.2$. Forgetting the product rule is equivalent to forgetting the Coriolis term $2\dot{r}\dot{\theta}$ in polar form.
\end{attention}

\textbf{Part D.}

\textbf{10.} At $t = 10\ \mathrm{s}$: $|a_r| = 0.36\ \mathrm{m\,s^{-2}}$ and $|a_\theta| = 0.20\ \mathrm{m\,s^{-2}}$. The \textbf{radial (centripetal-type) effect dominates}, pulling the container \emph{inward} towards the axis; the cable and trolley restraint must therefore mainly resist this inward radial pull, while still catering for the non-negligible sideways Coriolis-type push (about $56\%$ of the radial value). Note also that $a_r = -r\omega^{2}$ grows as the trolley moves outward, so the restraint requirement increases during the operation.

\textbf{11.} Plenary discussion. The polar description is superior here because the motion is \emph{organised around a fixed centre}: the data ($\omega$, $u$) are given directly as $\dot{\theta}$ and $\dot{r}$, the formulas for $\vec{v}$ and $\vec{a}$ are short, and each component has a direct physical meaning (radial pull, sideways push). The Cartesian description requires the product and chain rules at every differentiation and mixes the two effects together. Cartesian coordinates would be preferable for motion along straight guides or when forces are given in fixed horizontal/vertical directions.

\begin{aretenir}[What this activity consolidates]
\begin{itemize}
\item Constant $\dot{r}$ and $\dot{\theta}$ do \emph{not} imply zero acceleration: rotation alone produces $-r\dot{\theta}^{2}\,\vec{e}_r$, and sliding on a rotating arm produces $2\dot{r}\dot{\theta}\,\vec{e}_\theta$.
\item Polar and Cartesian descriptions are two views of the same vector quantities; magnitudes are invariant.
\item The choice of coordinate system is a modelling decision: choose the system that matches the geometry of the problem.
\end{itemize}
\end{aretenir}

\newpage

\coursec{Methods and worked examples}

\courssub{Methods and Worked Examples}

\begin{methode}[Finding Cartesian Components of Velocity and Acceleration]
\begin{enumerate}
\item Write the position vector $\vec{r}(t) = x(t)\,\vec{i} + y(t)\,\vec{j}$.
\item Differentiate once with respect to time to obtain the velocity vector: $\vec{v}(t) = \dot{x}(t)\,\vec{i} + \dot{y}(t)\,\vec{j}$.
\item Differentiate again to obtain the acceleration vector: $\vec{a}(t) = \ddot{x}(t)\,\vec{i} + \ddot{y}(t)\,\vec{j}$.
\item The speed is the magnitude of the velocity vector: $v = \sqrt{\dot{x}^2 + \dot{y}^2}$.
\item To find the direction of motion, compute the angle $\phi$ that $\vec{v}$ makes with the positive $x$-axis: $\tan\phi = \frac{\dot{y}}{\dot{x}}$ (taking quadrant into account).
\end{enumerate}
\end{methode}

\begin{exoresolu}[Cartesian Components – Velocity, Acceleration and Speed]
A particle moves in the $xy$-plane such that its position vector at time $t$ seconds is given by
\[
\vec{r}(t) = (t^3 - 3t)\,\vec{i} + (2t^2 - 4)\,\vec{j} \quad (\text{metres}).
\]
\begin{enumerate}
\item Find the velocity and acceleration vectors at time $t$.
\item Determine the speed of the particle when $t = 2$ s.
\item Find the angle between the velocity and acceleration vectors at $t = 2$ s.
\end{enumerate}
\tcblower
\textbf{Solution.}
\begin{enumerate}
\item Differentiating $\vec{r}(t)$:
\[
\vec{v}(t) = \frac{d\vec{r}}{dt} = (3t^2 - 3)\,\vec{i} + (4t)\,\vec{j} \ \mathrm{m\,s^{-1}}.
\]
\[
\vec{a}(t) = \frac{d\vec{v}}{dt} = (6t)\,\vec{i} + 4\,\vec{j} \ \mathrm{m\,s^{-2}}.
\]
\item At $t = 2$ s:
\[
\vec{v}(2) = (3\cdot4 - 3)\,\vec{i} + 8\,\vec{j} = 9\,\vec{i} + 8\,\vec{j} \ \mathrm{m\,s^{-1}}.
\]
Speed $v = |\vec{v}(2)| = \sqrt{9^2 + 8^2} = \sqrt{81+64} = \sqrt{145} \approx 12.04 \ \mathrm{m\,s^{-1}}$.
\item At $t = 2$ s:
\[
\vec{a}(2) = 12\,\vec{i} + 4\,\vec{j} \ \mathrm{m\,s^{-2}}.
\]
The angle $\theta$ between $\vec{v}(2)$ and $\vec{a}(2)$ is given by
\[
\cos\theta = \frac{\vec{v}(2)\cdot\vec{a}(2)}{|\vec{v}(2)|\,|\vec{a}(2)|}
= \frac{9\cdot12 + 8\cdot4}{\sqrt{145}\cdot\sqrt{12^2+4^2}}
= \frac{108+32}{\sqrt{145}\cdot\sqrt{160}}
= \frac{140}{\sqrt{145}\cdot4\sqrt{10}} = \frac{35}{\sqrt{1450}}.
\]
Thus $\theta = \cos^{-1}\left(\frac{35}{\sqrt{1450}}\right) \approx 22.6^\circ$.
\end{enumerate}
\end{exoresolu}

\begin{methode}[Expressing Velocity and Acceleration in Polar Coordinates]
\begin{enumerate}
\item Let the position of a particle be given by polar coordinates $(r,\theta)$ with unit vectors $\vec{e}_r$ (radial) and $\vec{e}_\theta$ (transverse).
\item The position vector is $\vec{r} = r\,\vec{e}_r$.
\item The velocity vector is $\vec{v} = \dot{r}\,\vec{e}_r + r\dot{\theta}\,\vec{e}_\theta$.
\item The acceleration vector is $\vec{a} = (\ddot{r} - r\dot{\theta}^2)\,\vec{e}_r + (r\ddot{\theta} + 2\dot{r}\dot{\theta})\,\vec{e}_\theta$.
\item If the path is given as $r = f(\theta)$ and $\dot{\theta}$ is known, compute $\dot{r} = f'(\theta)\dot{\theta}$ and $\ddot{r} = f''(\theta)\dot{\theta}^2 + f'(\theta)\ddot{\theta}$.
\end{enumerate}
\end{methode}

\begin{exoresolu}[Polar Components – Given Path and Angular Speed]
A particle moves along the curve $r = 2\theta$ (with $r$ in metres, $\theta$ in radians). At the instant when $\theta = \frac{\pi}{4}$, the angular velocity is $\dot{\theta} = 3\ \mathrm{rad\,s^{-1}}$ and the angular acceleration is $\ddot{\theta} = 0$.
\begin{enumerate}
\item Find the radial and transverse components of velocity at this instant.
\item Find the radial and transverse components of acceleration at this instant.
\item Express the velocity and acceleration vectors in terms of $\vec{e}_r$ and $\vec{e}_\theta$.
\end{enumerate}
\tcblower
\textbf{Solution.}
Given $r = 2\theta$, we have $\dot{r} = 2\dot{\theta}$ and $\ddot{r} = 2\ddot{\theta}$.
At $\theta = \frac{\pi}{4}$:
$r = 2\cdot\frac{\pi}{4} = \frac{\pi}{2}$ m.
$\dot{\theta} = 3$, $\ddot{\theta} = 0$.
Thus $\dot{r} = 2\times3 = 6\ \mathrm{m\,s^{-1}}$, $\ddot{r} = 0$.
\begin{enumerate}
\item Radial velocity: $v_r = \dot{r} = 6\ \mathrm{m\,s^{-1}}$.
Transverse velocity: $v_\theta = r\dot{\theta} = \frac{\pi}{2}\times3 = \frac{3\pi}{2}\ \mathrm{m\,s^{-1}}$.
\item Radial acceleration: $a_r = \ddot{r} - r\dot{\theta}^2 = 0 - \frac{\pi}{2}\times 9 = -\frac{9\pi}{2}\ \mathrm{m\,s^{-2}}$.
Transverse acceleration: $a_\theta = r\ddot{\theta} + 2\dot{r}\dot{\theta} = \frac{\pi}{2}\times0 + 2\times6\times3 = 36\ \mathrm{m\,s^{-2}}$.
\item $\vec{v} = 6\,\vec{e}_r + \frac{3\pi}{2}\,\vec{e}_\theta \ \mathrm{m\,s^{-1}}$.
$\vec{a} = -\frac{9\pi}{2}\,\vec{e}_r + 36\,\vec{e}_\theta \ \mathrm{m\,s^{-2}}$.
\end{enumerate}
\end{exoresolu}

\begin{methode}[Deriving Radial and Transverse Components of Velocity and Acceleration]
\begin{enumerate}
\item Start with the position vector in polar coordinates: $\vec{r} = r\,\vec{e}_r$.
\item Recall that the unit vectors $\vec{e}_r$ and $\vec{e}_\theta$ rotate with the particle: $\frac{d\vec{e}_r}{dt} = \dot{\theta}\,\vec{e}_\theta$, $\frac{d\vec{e}_\theta}{dt} = -\dot{\theta}\,\vec{e}_r$.
\item Differentiate $\vec{r}$ once: $\vec{v} = \dot{r}\,\vec{e}_r + r\frac{d\vec{e}_r}{dt} = \dot{r}\,\vec{e}_r + r\dot{\theta}\,\vec{e}_\theta$.
\item Differentiate $\vec{v}$: $\vec{a} = \ddot{r}\,\vec{e}_r + \dot{r}\frac{d\vec{e}_r}{dt} + \dot{r}\dot{\theta}\,\vec{e}_\theta + r\ddot{\theta}\,\vec{e}_\theta + r\dot{\theta}\frac{d\vec{e}_\theta}{dt}$.
\item Substitute the derivatives of the unit vectors and collect terms in $\vec{e}_r$ and $\vec{e}_\theta$:
$\vec{a} = (\ddot{r} - r\dot{\theta}^2)\,\vec{e}_r + (r\ddot{\theta} + 2\dot{r}\dot{\theta})\,\vec{e}_\theta$.
\end{enumerate}
\end{methode}

\begin{exoresolu}[Derivation of Radial-Transverse Acceleration Components]
Derive the expressions for the radial and transverse components of acceleration for a particle moving in a plane, using polar coordinates $(r,\theta)$ and the rotating unit vectors $\vec{e}_r$, $\vec{e}_\theta$.
\tcblower
\textbf{Proof.}
Let the position vector be $\vec{r} = r\,\vec{e}_r$, where $\vec{e}_r = \cos\theta\,\vec{i} + \sin\theta\,\vec{j}$ and $\vec{e}_\theta = -\sin\theta\,\vec{i} + \cos\theta\,\vec{j}$.
Differentiating $\vec{e}_r$ and $\vec{e}_\theta$ with respect to time:
\[
\frac{d\vec{e}_r}{dt} = (-\sin\theta\,\vec{i} + \cos\theta\,\vec{j})\dot{\theta} = \dot{\theta}\,\vec{e}_\theta,
\]
\[
\frac{d\vec{e}_\theta}{dt} = (-\cos\theta\,\vec{i} - \sin\theta\,\vec{j})\dot{\theta} = -\dot{\theta}\,\vec{e}_r.
\]
Now differentiate $\vec{r}$:
\[
\vec{v} = \frac{d\vec{r}}{dt} = \dot{r}\,\vec{e}_r + r\frac{d\vec{e}_r}{dt} = \dot{r}\,\vec{e}_r + r\dot{\theta}\,\vec{e}_\theta.
\]
Differentiate $\vec{v}$:
\[
\vec{a} = \frac{d\vec{v}}{dt} = \ddot{r}\,\vec{e}_r + \dot{r}\frac{d\vec{e}_r}{dt} + \dot{r}\dot{\theta}\,\vec{e}_\theta + r\ddot{\theta}\,\vec{e}_\theta + r\dot{\theta}\frac{d\vec{e}_\theta}{dt}.
\]
Substitute $\frac{d\vec{e}_r}{dt} = \dot{\theta}\,\vec{e}_\theta$ and $\frac{d\vec{e}_\theta}{dt} = -\dot{\theta}\,\vec{e}_r$:
\[
\vec{a} = \ddot{r}\,\vec{e}_r + \dot{r}\dot{\theta}\,\vec{e}_\theta + \dot{r}\dot{\theta}\,\vec{e}_\theta + r\ddot{\theta}\,\vec{e}_\theta - r\dot{\theta}^2\,\vec{e}_r.
\]
Group the $\vec{e}_r$ and $\vec{e}_\theta$ terms:
\[
\vec{a} = (\ddot{r} - r\dot{\theta}^2)\,\vec{e}_r + (r\ddot{\theta} + 2\dot{r}\dot{\theta})\,\vec{e}_\theta.
\]
Hence the radial component of acceleration is $a_r = \ddot{r} - r\dot{\theta}^2$ and the transverse component is $a_\theta = r\ddot{\theta} + 2\dot{r}\dot{\theta}$.

\begin{popfigure}
\begin{tikzpicture}[scale=1.2]
\draw[->] (-2,0) -- (2,0) node[right] {$x$};
\draw[->] (0,-2) -- (0,2) node[above] {$y$};
\draw[thick,->] (0,0) -- (1.5,1) node[above right] {$\vec{r}$};
\draw[dashed] (0,0) -- (1.5,0);
\draw[dashed] (1.5,0) -- (1.5,1);
\node at (0.8,0.2) {$\theta$};
\draw[->,popblue] (1.5,1) -- (2.2,1.5) node[above right] {$\vec{e}_r$};
\draw[->,poporange] (1.5,1) -- (1,2) node[above left] {$\vec{e}_\theta$};
\end{tikzpicture}
\legende{Polar coordinates and rotating unit vectors.}
\end{popfigure}
\end{exoresolu}

\begin{methode}[Analysing Simple Radial-Transverse Motion]
\begin{enumerate}
\item Identify the given components of acceleration ($a_r$, $a_\theta$) or constraints (e.g., $a_\theta = 0$, $r = \text{constant}$, $\dot{\theta} = \text{constant}$).
\item Write the equations of motion:
$a_r = \ddot{r} - r\dot{\theta}^2$,
$a_\theta = r\ddot{\theta} + 2\dot{r}\dot{\theta}$.
\item Use the given conditions to simplify. Common cases:
\begin{itemize}
\item Circular motion: $r = \text{constant} \Rightarrow \dot{r}=0,\ \ddot{r}=0$. Then $a_r = -r\dot{\theta}^2$, $a_\theta = r\ddot{\theta}$.
\item Central force / zero transverse acceleration: $a_\theta = 0 \Rightarrow r\ddot{\theta} + 2\dot{r}\dot{\theta} = 0 \Rightarrow \frac{d}{dt}(r^2\dot{\theta}) = 0 \Rightarrow r^2\dot{\theta} = h$ (constant).
\item Constant radial acceleration: $a_r = \text{constant}$.
\end{itemize}
\item Solve the resulting differential equations using initial conditions.
\end{enumerate}
\end{methode}

\begin{exoresolu}[Circular Motion with Constant Angular Speed]
A particle moves in a circle of radius $a = 3$ m with constant angular speed $\omega = 4\ \mathrm{rad\,s^{-1}}$.
\begin{enumerate}
\item Write down the radial and transverse components of velocity and acceleration.
\item Find the magnitude of the acceleration.
\item If the particle has mass $m = 2$ kg, find the magnitude of the net force acting on it.
\end{enumerate}
\tcblower
\textbf{Solution.}
For circular motion, $r = a = 3$ m, $\dot{r} = 0$, $\ddot{r} = 0$. $\dot{\theta} = \omega = 4$, $\ddot{\theta} = 0$.
\begin{enumerate}
\item Radial velocity: $v_r = \dot{r} = 0$.
Transverse velocity: $v_\theta = r\dot{\theta} = 3 \times 4 = 12\ \mathrm{m\,s^{-1}}$.
Radial acceleration: $a_r = \ddot{r} - r\dot{\theta}^2 = 0 - 3 \times 16 = -48\ \mathrm{m\,s^{-2}}$.
Transverse acceleration: $a_\theta = r\ddot{\theta} + 2\dot{r}\dot{\theta} = 0$.
\item Magnitude of acceleration: $a = \sqrt{a_r^2 + a_\theta^2} = 48\ \mathrm{m\,s^{-2}}$ (directed radially inward).
\item Net force magnitude: $F = m a = 2 \times 48 = 96\ \mathrm{N}$ (directed toward the centre).
\end{enumerate}
\end{exoresolu}

\begin{methode}[Using Conservation of Angular Momentum (Zero Transverse Acceleration)]
\begin{enumerate}
\item When the transverse acceleration is zero ($a_\theta = 0$), the angular momentum per unit mass $h = r^2\dot{\theta}$ is constant.
\item This occurs for any central force (force directed toward or away from a fixed point).
\item The radial equation becomes $\ddot{r} - r\dot{\theta}^2 = a_r$. Substitute $\dot{\theta} = h/r^2$ to get $\ddot{r} - \frac{h^2}{r^3} = a_r$.
\item If $a_r$ is given as a function of $r$, this is a differential equation for $r(t)$. Alternatively, use the energy integral if the force is conservative.
\end{enumerate}
\end{methode}

\begin{exoresolu}[Central Force Motion – Inverse Square Law]
A particle of unit mass moves under a central attractive force of magnitude $\frac{\mu}{r^2}$ (where $\mu > 0$ is constant). Initially, the particle is at distance $r = r_0$ from the centre, with radial velocity $\dot{r} = 0$ and transverse velocity $v_\theta = v_0$.
\begin{enumerate}
\item Show that the angular momentum per unit mass $h = r^2\dot{\theta}$ is constant and find its value in terms of $r_0$ and $v_0$.
\item Derive the radial equation of motion: $\ddot{r} - \frac{h^2}{r^3} = -\frac{\mu}{r^2}$.
\item Using the substitution $u = 1/r$, show that the orbit equation is $\frac{d^2u}{d\theta^2} + u = \frac{\mu}{h^2}$.
\end{enumerate}
\tcblower
\textbf{Solution.}
\begin{enumerate}
\item Since the force is central, the transverse acceleration is zero: $a_\theta = 0$. Hence $\frac{d}{dt}(r^2\dot{\theta}) = 0$, so $h = r^2\dot{\theta}$ is constant. At $t=0$, $r = r_0$, $\dot{\theta} = v_\theta/r_0 = v_0/r_0$. Thus $h = r_0^2 \cdot \frac{v_0}{r_0} = r_0 v_0$.
\item The radial acceleration is $a_r = \ddot{r} - r\dot{\theta}^2$. The force per unit mass is $-\mu/r^2$ (negative because attractive). So $\ddot{r} - r\dot{\theta}^2 = -\frac{\mu}{r^2}$. Substitute $\dot{\theta} = h/r^2$: $\ddot{r} - r\frac{h^2}{r^4} = -\frac{\mu}{r^2}$, i.e., $\ddot{r} - \frac{h^2}{r^3} = -\frac{\mu}{r^2}$.
\item Let $u = 1/r$. Then $r = 1/u$, $\dot{r} = -\frac{1}{u^2}\dot{u} = -\frac{1}{u^2}\frac{du}{d\theta}\dot{\theta} = -h\frac{du}{d\theta}$ (since $\dot{\theta} = h u^2$). Then $\ddot{r} = \frac{d}{dt}(-h\frac{du}{d\theta}) = -h\frac{d^2u}{d\theta^2}\dot{\theta} = -h^2 u^2 \frac{d^2u}{d\theta^2}$.
Substitute into the radial equation:
$-h^2 u^2 \frac{d^2u}{d\theta^2} - h^2 u^3 = -\mu u^2$.
Divide by $-h^2 u^2$ (assuming $u \neq 0$):
$\frac{d^2u}{d\theta^2} + u = \frac{\mu}{h^2}$.
This is the orbit equation.
\end{enumerate}
\end{exoresolu}

\begin{methode}[Solving Motion with Given Radial and Transverse Accelerations]
\begin{enumerate}
\item Write the two equations:
$a_r = \ddot{r} - r\dot{\theta}^2$,
$a_\theta = r\ddot{\theta} + 2\dot{r}\dot{\theta}$.
\item If $a_\theta = 0$, use $r^2\dot{\theta} = h$ (constant) to eliminate $\dot{\theta}$.
\item If $a_r$ is a function of $r$ only, the radial equation becomes $\ddot{r} - \frac{h^2}{r^3} = a_r(r)$.
\item Multiply by $\dot{r}$ and integrate to obtain the energy equation: $\frac{1}{2}\dot{r}^2 + \frac{h^2}{2r^2} + V(r) = E$, where $V(r)$ is the potential such that $a_r = -dV/dr$.
\item Alternatively, if the path $r(\theta)$ is required, use the substitution $u = 1/r$ and the orbit equation $\frac{d^2u}{d\theta^2} + u = -\frac{a_r}{h^2 u^2}$.
\end{enumerate}
\end{methode}

\begin{exoresolu}[Radial Oscillation with Zero Transverse Acceleration]
A particle moves in a plane such that its transverse acceleration is zero and its radial acceleration is $a_r = -\omega^2 r$, where $\omega > 0$ is constant. Initially, at $t=0$, $r = a$, $\dot{r} = 0$, and $\dot{\theta} = \Omega$.
\begin{enumerate}
\item Show that $r^2\dot{\theta} = h$ is constant and find $h$ in terms of $a$ and $\Omega$.
\item Show that the radial motion satisfies $\ddot{r} + \omega^2 r = \frac{h^2}{r^3}$.
\item Given that $\omega^2 = \frac{h^2}{a^4}$, solve for $r(t)$ and describe the motion.
\end{enumerate}
\tcblower
\textbf{Solution.}
\begin{enumerate}
\item $a_\theta = 0 \Rightarrow \frac{d}{dt}(r^2\dot{\theta}) = 0 \Rightarrow r^2\dot{\theta} = h$ (constant). At $t=0$, $r=a$, $\dot{\theta}=\Omega$, so $h = a^2\Omega$.
\item The radial equation is $\ddot{r} - r\dot{\theta}^2 = a_r = -\omega^2 r$. Substitute $\dot{\theta} = h/r^2$:
$\ddot{r} - r\frac{h^2}{r^4} = -\omega^2 r \Rightarrow \ddot{r} - \frac{h^2}{r^3} = -\omega^2 r$.
Rearrange: $\ddot{r} + \omega^2 r = \frac{h^2}{r^3}$.
\item Given $\omega^2 = h^2/a^4$. At $t=0$, $r=a$, $\dot{r}=0$. Check if $r=a$ is an equilibrium: $\ddot{r} = 0$ when $r=a$? LHS: $\omega^2 a = (h^2/a^4)a = h^2/a^3$. RHS: $h^2/a^3$. So $r=a$ satisfies the equation with $\ddot{r}=0$. But we need to see if the particle stays at $r=a$. Since $\dot{r}=0$ initially and $\ddot{r}=0$, the particle remains at $r=a$ for all $t$. Then $\dot{\theta} = h/a^2 = \Omega$ constant. The motion is uniform circular motion of radius $a$ with angular speed $\Omega$. (If the initial conditions were slightly different, the radial motion would be oscillatory about $r=a$.)
\end{enumerate}
\end{exoresolu}

\newpage

\coursec{Exercises}

\courssub{I check my knowledge}

\begin{exercice}[MCQ: Cartesian acceleration]
A particle moves in the plane with position vector $\vec{r}(t) = (3t^2)\,\vec{i} + (4t)\,\vec{j}$ (SI units). What is its acceleration vector at $t = 2\,\mathrm{s}$?
\begin{enumerate}
\item $6\,\vec{i} + 4\,\vec{j}$
\item $6\,\vec{i}$
\item $12\,\vec{i} + 8\,\vec{j}$
\item $12\,\vec{i}$
\end{enumerate}
\end{exercice}

\begin{exercice}[True/False: Radial acceleration in polar coordinates]
\textbf{Statement:} In polar coordinates, the radial component of acceleration is always $\ddot{r}$.
\begin{itemize}
\item State whether the statement is true or false.
\item Justify your answer by writing the correct expression for the radial component of acceleration.
\end{itemize}
\end{exercice}

\begin{exercice}[Definition: Polar unit vectors]
Define the radial unit vector $\vec{e}_r$ and the transverse unit vector $\vec{e}_\theta$ in terms of the Cartesian unit vectors $\vec{i}$ and $\vec{j}$ and the angle $\theta$. State how they vary with $\theta$.
\end{exercice}

\begin{exercice}[Direct calculation: Velocity, acceleration, speed and angle]
A particle moves with position vector $\vec{r}(t) = (t^2 + 1)\,\vec{i} + (2t^3 - t)\,\vec{j}$ (SI units). At $t = 1\,\mathrm{s}$, find:
\begin{enumerate}
\item the velocity vector $\vec{v}$,
\item the acceleration vector $\vec{a}$,
\item the speed $v$,
\item the angle that $\vec{v}$ makes with the $Ox$ axis.
\end{enumerate}
\end{exercice}

\courssub{I practise}

\begin{exercice}[Uniform circular motion in Cartesian coordinates]
A particle moves such that its coordinates are $x = 3\cos(2t)$ and $y = 3\sin(2t)$ (distances in metres, time in seconds).
\begin{enumerate}
\item Show that the trajectory is a circle centred at the origin.
\item Prove that the motion is uniform circular motion.
\item Find the speed of the particle.
\item Show that the acceleration is directed towards the origin and has magnitude $12\,\mathrm{m\,s^{-2}}$.
\end{enumerate}
\end{exercice}

\begin{exercice}[Derivation of polar components of velocity and acceleration]
Starting from the position vector $\vec{r} = r\,\vec{e}_r$, where $\vec{e}_r$ and $\vec{e}_\theta$ are the polar unit vectors, derive the expressions for:
\begin{enumerate}
\item the velocity vector $\vec{v}$ in terms of $\dot{r}$, $r$, $\dot{\theta}$, $\vec{e}_r$, $\vec{e}_\theta$,
\item the acceleration vector $\vec{a}$ in terms of $\ddot{r}$, $\dot{r}$, $r$, $\ddot{\theta}$, $\dot{\theta}$, $\vec{e}_r$, $\vec{e}_\theta$.
\end{enumerate}
Clearly show the differentiation steps and the use of $\frac{d\vec{e}_r}{dt} = \dot{\theta}\,\vec{e}_\theta$ and $\frac{d\vec{e}_\theta}{dt} = -\dot{\theta}\,\vec{e}_r$.
\end{exercice}

\begin{exercice}[Spiral motion in polar coordinates]
A particle describes the spiral given by $r = 2t$ and $\theta = 3t$ (SI units).
\begin{enumerate}
\item Find the radial and transverse components of velocity, $v_r$ and $v_\theta$, at $t = 2\,\mathrm{s}$.
\item Find the radial and transverse components of acceleration, $a_r$ and $a_\theta$, at $t = 2\,\mathrm{s}$.
\item Deduce the magnitudes of the velocity and acceleration vectors at that instant.
\end{enumerate}
\end{exercice}

\begin{exercice}[Car on a roundabout]
A car travels round a roundabout of radius $20\,\mathrm{m}$ at a constant speed of $8\,\mathrm{m\,s^{-1}}$.
\begin{enumerate}
\item Find the angular velocity $\dot{\theta}$ of the car.
\item Find the centripetal acceleration of the car.
\item Comment on the direction of the acceleration vector relative to the velocity vector.
\end{enumerate}
\end{exercice}

\courssub{I prepare for the GCE}

\begin{exercice}[Bead on a rotating rod]
A bead slides outward on a smooth rod which rotates in a horizontal plane with constant angular velocity $\omega$. The distance of the bead from the pivot is given by $r = r_0 e^{kt}$, where $r_0$ and $k$ are positive constants.
\begin{enumerate}
\item Show that the transverse component of the bead's acceleration is $2k\omega r$.
\item Find the radial component of the acceleration.
\item Hence write down the magnitude of the total acceleration.
\end{enumerate}
\end{exercice}

\begin{exercice}[Particle on a circle with angular acceleration]
A particle moves on a circle of radius $5\,\mathrm{m}$. Its angular acceleration is constant at $\ddot{\theta} = 0.4\,\mathrm{rad\,s^{-2}}$. At the instant when $\dot{\theta} = 2\,\mathrm{rad\,s^{-1}}$, find:
\begin{enumerate}
\item the radial component of acceleration,
\item the transverse component of acceleration,
\item the magnitude of the total acceleration.
\end{enumerate}
\end{exercice}

\begin{exercice}[Satellite dish tracking]
A satellite dish tracks a target such that the distance from the pivot to the target is constant at $r = 50\,\mathrm{m}$ and the dish rotates with constant angular velocity $\dot{\theta} = 0.1\,\mathrm{rad\,s^{-1}}$. A small part slides along the dish arm with radial velocity $\dot{r} = 0.3\,\mathrm{m\,s^{-1}}$ (assume $\ddot{r} = 0$).
\begin{enumerate}
\item Find the radial and transverse components of the acceleration of the part.
\item Find the magnitude of the total acceleration.
\item Explain the physical origin of each term in the acceleration components (e.g., centripetal, Coriolis, etc.).
\end{enumerate}
\end{exercice}

\newpage

\coursec{Solutions to the exercises}

\begin{corrige}[Exercise 1 — MCQ: Cartesian acceleration]
We differentiate the position vector twice with respect to time: position $\to$ velocity $\to$ acceleration.

Given $\vec{r}(t) = (3t^2)\,\vec{i} + (4t)\,\vec{j}$:
\begin{align*}
\vec{v}(t) &= \frac{d\vec{r}}{dt} = (6t)\,\vec{i} + 4\,\vec{j},\\[2pt]
\vec{a}(t) &= \frac{d\vec{v}}{dt} = 6\,\vec{i} + 0\,\vec{j} = 6\,\vec{i}.
\end{align*}
The acceleration is \emph{constant} (independent of $t$), so at $t = 2\,\mathrm{s}$:
\[ \vec{a}(2) = 6\,\vec{i}\ \mathrm{m\,s^{-2}}. \]

\textbf{Correct answer: (b)} $6\,\vec{i}$.
\end{corrige}

\begin{attention}[Trap in Exercise 1]
The trap is to differentiate only once, or to evaluate the wrong vector at $t = 2$. Option (c), $12\vec{i}+8\vec{j}$, corresponds to evaluating the \emph{velocity} $\vec{v}(t) = 6t\,\vec{i} + 4\,\vec{j}$ at $t = 2$, not the acceleration. Always check how many differentiations are required: position $\to$ velocity $\to$ acceleration, and check \emph{which} vector the question asks for before substituting the value of $t$.
\end{attention}

\begin{corrige}[Exercise 2 — True/False: Radial acceleration in polar coordinates]
\textbf{The statement is FALSE.}

When we differentiate $\vec{v} = \dot{r}\,\vec{e}_r + r\dot{\theta}\,\vec{e}_\theta$, the rotation of the basis vectors produces extra terms. The correct expression is
\[ a_r = \ddot{r} - r\dot{\theta}^2. \]

The term $\ddot{r}$ alone gives the radial acceleration only when the transverse direction does not rotate, i.e. when $\dot{\theta} = 0$ (purely radial motion along a fixed line). As soon as the particle has angular motion, the \emph{centripetal} term $-r\dot{\theta}^2$ appears, directed towards the pole.

\textbf{Counter-example:} for uniform circular motion, $r = R$ constant, so $\dot{r} = \ddot{r} = 0$, yet the radial acceleration is $a_r = -R\dot{\theta}^2 \neq 0$. This proves that $a_r \neq \ddot{r}$ in general.
\end{corrige}

\begin{corrige}[Exercise 3 — Definition: Polar unit vectors]
Let $Oxy$ be a Cartesian frame with unit vectors $\vec{i}, \vec{j}$, and let $M$ be a point with polar coordinates $(r,\theta)$, where $\theta = (\vec{i}, \overrightarrow{OM})$.

The \cle{radial unit vector} $\vec{e}_r$ is the unit vector along $\overrightarrow{OM}$, and the \cle{transverse unit vector} $\vec{e}_\theta$ is obtained from $\vec{e}_r$ by a rotation of $+\dfrac{\pi}{2}$ (direct sense):
\[ \vec{e}_r = \cos\theta\,\vec{i} + \sin\theta\,\vec{j}, \qquad \vec{e}_\theta = -\sin\theta\,\vec{i} + \cos\theta\,\vec{j}. \]

\textbf{Variation with $\theta$:} unlike $\vec{i}$ and $\vec{j}$, which are fixed, the polar vectors \emph{depend on $\theta$} and therefore rotate as the particle moves. Differentiating with respect to $\theta$:
\[ \frac{d\vec{e}_r}{d\theta} = -\sin\theta\,\vec{i} + \cos\theta\,\vec{j} = \vec{e}_\theta, \qquad \frac{d\vec{e}_\theta}{d\theta} = -\cos\theta\,\vec{i} - \sin\theta\,\vec{j} = -\vec{e}_r. \]
Hence, by the chain rule, $\dfrac{d\vec{e}_r}{dt} = \dot{\theta}\,\vec{e}_\theta$ and $\dfrac{d\vec{e}_\theta}{dt} = -\dot{\theta}\,\vec{e}_r$. Note also that $(\vec{e}_r, \vec{e}_\theta)$ is always an orthonormal direct basis: $\vec{e}_r \cdot \vec{e}_\theta = 0$ and $\|\vec{e}_r\| = \|\vec{e}_\theta\| = 1$.
\end{corrige}

\begin{popfigure}
\begin{center}
\begin{tikzpicture}[scale=2.2, >=stealth]
\draw[->, gray] (-0.3,0) -- (2.3,0) node[below left] {$x$};
\draw[->, gray] (0,-0.3) -- (0,1.8) node[below left] {$y$};
\draw[->, thick, popblue] (0,0) -- (0.9,0) node[below] {$\vec{i}$};
\draw[->, thick, popblue] (0,0) -- (0,0.9) node[left] {$\vec{j}$};
\coordinate (M) at (1.5,1.2);
\draw[dashed, popmuted] (0,0) -- (M);
\fill[popdark] (M) circle (0.02) node[above right] {$M(r,\theta)$};
\draw[->, thick, poporange] (M) -- ($(M)+(0.39,0.312)$) node[right] {$\vec{e}_r$};
\draw[->, thick, popgreen] (M) -- ($(M)+(-0.312,0.39)$) node[above] {$\vec{e}_\theta$};
\draw[->, poppurple] (0.55,0) arc (0:38.7:0.55) node[midway, right] {$\theta$};
\end{tikzpicture}
\end{center}
\legende{The rotating polar basis $(\vec{e}_r, \vec{e}_\theta)$ attached to the moving point $M$, compared with the fixed Cartesian basis $(\vec{i}, \vec{j})$.}
\end{popfigure}

\begin{corrige}[Exercise 4 — Direct calculation: Velocity, acceleration, speed and angle]
Given $\vec{r}(t) = (t^2+1)\,\vec{i} + (2t^3 - t)\,\vec{j}$.

\textbf{1. Velocity vector.} Differentiate each component:
\[ \vec{v}(t) = \frac{d\vec{r}}{dt} = (2t)\,\vec{i} + (6t^2 - 1)\,\vec{j}. \]
At $t = 1$: $\vec{v}(1) = 2\,\vec{i} + 5\,\vec{j}\ \mathrm{m\,s^{-1}}$.

\textbf{2. Acceleration vector.} Differentiate again:
\[ \vec{a}(t) = \frac{d\vec{v}}{dt} = 2\,\vec{i} + (12t)\,\vec{j}. \]
At $t = 1$: $\vec{a}(1) = 2\,\vec{i} + 12\,\vec{j}\ \mathrm{m\,s^{-2}}$.

\textbf{3. Speed.} The speed is the magnitude of the velocity:
\[ v = \|\vec{v}(1)\| = \sqrt{2^2 + 5^2} = \sqrt{4 + 25} = \sqrt{29} \approx 5.39\ \mathrm{m\,s^{-1}}. \]

\textbf{4. Angle with the $Ox$ axis.} If $\alpha$ is the angle between $\vec{v}$ and $Ox$, then
\[ \tan\alpha = \frac{v_y}{v_x} = \frac{5}{2} = 2.5 \quad\Longrightarrow\quad \alpha = \arctan(2.5) \approx 68.2^\circ. \]
Since both components are positive, the vector lies in the first quadrant and this value of $\alpha$ is the correct one (no ambiguity of $\arctan$).
\end{corrige}

\begin{corrige}[Exercise 5 — Uniform circular motion in Cartesian coordinates]
Given $x = 3\cos(2t)$, $y = 3\sin(2t)$.

\textbf{1. Trajectory is a circle centred at $O$.}
\[ x^2 + y^2 = 9\cos^2(2t) + 9\sin^2(2t) = 9\bigl(\cos^2(2t)+\sin^2(2t)\bigr) = 9. \]
So every point of the trajectory satisfies $x^2+y^2 = 3^2$: the particle moves on the \cle{circle} of centre $O$ and radius $R = 3\,\mathrm{m}$. $\blacksquare$

\textbf{2. The motion is uniform circular.} We must show that the angular velocity is constant. The polar angle satisfies
\[ \tan\theta = \frac{y}{x} = \frac{3\sin(2t)}{3\cos(2t)} = \tan(2t) \quad\Longrightarrow\quad \theta = 2t \ (\text{mod } 2\pi). \]
Hence $\dot{\theta} = 2\,\mathrm{rad\,s^{-1}}$, a \emph{constant}. Motion on a circle with constant angular velocity is, by definition, \cle{uniform circular motion}. $\blacksquare$

\textbf{3. Speed.} Differentiate the coordinates:
\[ \dot{x} = -6\sin(2t), \qquad \dot{y} = 6\cos(2t). \]
\[ v = \sqrt{\dot{x}^2 + \dot{y}^2} = \sqrt{36\sin^2(2t) + 36\cos^2(2t)} = \sqrt{36} = 6\ \mathrm{m\,s^{-1}}. \]
(Check: $v = R\dot{\theta} = 3 \times 2 = 6\,\mathrm{m\,s^{-1}}$. The speed is constant, confirming ``uniform''.)

\textbf{4. Acceleration.} Differentiate once more:
\[ \ddot{x} = -12\cos(2t), \qquad \ddot{y} = -12\sin(2t). \]
Therefore
\[ \vec{a} = -12\cos(2t)\,\vec{i} - 12\sin(2t)\,\vec{j} = -4\bigl(3\cos(2t)\,\vec{i} + 3\sin(2t)\,\vec{j}\bigr) = -4\,\vec{r}. \]
Since $\vec{a} = -4\,\overrightarrow{OM}$, the acceleration is collinear with $\overrightarrow{OM}$ and directed \emph{towards the origin} (centripetal). Its magnitude is
\[ \|\vec{a}\| = 4\|\vec{r}\| = 4 \times 3 = 12\ \mathrm{m\,s^{-2}}. \]
(Check: $a = R\dot{\theta}^2 = 3 \times 4 = 12\,\mathrm{m\,s^{-2}} = \dfrac{v^2}{R} = \dfrac{36}{3}$.) $\blacksquare$
\end{corrige}

\begin{corrige}[Exercise 6 — Derivation of polar components of velocity and acceleration]
We start from $\vec{r} = r\,\vec{e}_r$ and use
\[ \frac{d\vec{e}_r}{dt} = \dot{\theta}\,\vec{e}_\theta, \qquad \frac{d\vec{e}_\theta}{dt} = -\dot{\theta}\,\vec{e}_r. \]

\textbf{1. Velocity.} Differentiate the product $r\,\vec{e}_r$ (both factors depend on $t$):
\begin{align*}
\vec{v} = \frac{d\vec{r}}{dt} &= \frac{d}{dt}\bigl(r\,\vec{e}_r\bigr) = \dot{r}\,\vec{e}_r + r\,\frac{d\vec{e}_r}{dt} \\
&= \dot{r}\,\vec{e}_r + r\bigl(\dot{\theta}\,\vec{e}_\theta\bigr).
\end{align*}
\[ \vec{v} = \dot{r}\,\vec{e}_r + r\dot{\theta}\,\vec{e}_\theta \qquad (v_r = \dot{r},\ v_\theta = r\dot{\theta}). \]

\textbf{2. Acceleration.} Differentiate $\vec{v}$ term by term, applying the product rule to each:
\begin{align*}
\vec{a} = \frac{d\vec{v}}{dt} &= \frac{d}{dt}\bigl(\dot{r}\,\vec{e}_r\bigr) + \frac{d}{dt}\bigl(r\dot{\theta}\,\vec{e}_\theta\bigr) \\[2pt]
&= \Bigl(\ddot{r}\,\vec{e}_r + \dot{r}\,\frac{d\vec{e}_r}{dt}\Bigr) + \Bigl(\dot{r}\dot{\theta}\,\vec{e}_\theta + r\ddot{\theta}\,\vec{e}_\theta + r\dot{\theta}\,\frac{d\vec{e}_\theta}{dt}\Bigr) \\[2pt]
&= \ddot{r}\,\vec{e}_r + \dot{r}\bigl(\dot{\theta}\vec{e}_\theta\bigr) + \dot{r}\dot{\theta}\,\vec{e}_\theta + r\ddot{\theta}\,\vec{e}_\theta + r\dot{\theta}\bigl(-\dot{\theta}\vec{e}_r\bigr) \\[2pt]
&= \ddot{r}\,\vec{e}_r - r\dot{\theta}^2\,\vec{e}_r + \dot{r}\dot{\theta}\,\vec{e}_\theta + \dot{r}\dot{\theta}\,\vec{e}_\theta + r\ddot{\theta}\,\vec{e}_\theta.
\end{align*}
Grouping the radial terms (in $\vec{e}_r$) and the transverse terms (in $\vec{e}_\theta$):
\[ \vec{a} = \bigl(\ddot{r} - r\dot{\theta}^2\bigr)\,\vec{e}_r + \bigl(r\ddot{\theta} + 2\dot{r}\dot{\theta}\bigr)\,\vec{e}_\theta, \]
so that
\[ a_r = \ddot{r} - r\dot{\theta}^2, \qquad a_\theta = r\ddot{\theta} + 2\dot{r}\dot{\theta} = \frac{1}{r}\frac{d}{dt}\bigl(r^2\dot{\theta}\bigr). \]
(The last equality follows from $\dfrac{d}{dt}\bigl(r^2\dot{\theta}\bigr) = 2r\dot{r}\dot{\theta} + r^2\ddot{\theta}$; dividing by $r$ gives $2\dot{r}\dot{\theta} + r\ddot{\theta}$.)
\end{corrige}

\begin{aretenir}[Names of the polar acceleration terms]
The term $-r\dot{\theta}^2$ is the \cle{centripetal acceleration}; the term $2\dot{r}\dot{\theta}$ is the \cle{Coriolis (compound) acceleration}, present whenever the particle moves radially \emph{and} the radius vector rotates simultaneously. The term $r\ddot{\theta}$ is the tangential acceleration due to angular speeding-up, and $\ddot{r}$ is the radial acceleration due to change of radial speed.
\end{aretenir}

\begin{corrige}[Exercise 7 — Spiral motion in polar coordinates]
Given $r = 2t$ and $\theta = 3t$, we have
\[ \dot{r} = 2, \quad \ddot{r} = 0, \quad \dot{\theta} = 3, \quad \ddot{\theta} = 0. \]
At $t = 2\,\mathrm{s}$: $r = 4\,\mathrm{m}$.

\textbf{1. Velocity components.}
\[ v_r = \dot{r} = 2\ \mathrm{m\,s^{-1}}, \qquad v_\theta = r\dot{\theta} = 4 \times 3 = 12\ \mathrm{m\,s^{-1}}. \]

\textbf{2. Acceleration components.}
\begin{align*}
a_r &= \ddot{r} - r\dot{\theta}^2 = 0 - 4 \times 3^2 = -36\ \mathrm{m\,s^{-2}},\\
a_\theta &= r\ddot{\theta} + 2\dot{r}\dot{\theta} = 0 + 2 \times 2 \times 3 = 12\ \mathrm{m\,s^{-2}}.
\end{align*}

\textbf{3. Magnitudes.} Since $(\vec{e}_r, \vec{e}_\theta)$ is orthonormal:
\[ v = \sqrt{v_r^2 + v_\theta^2} = \sqrt{2^2 + 12^2} = \sqrt{148} = 2\sqrt{37} \approx 12.2\ \mathrm{m\,s^{-1}}. \]
\[ a = \sqrt{a_r^2 + a_\theta^2} = \sqrt{(-36)^2 + 12^2} = \sqrt{1296 + 144} = \sqrt{1440} = 12\sqrt{10} \approx 37.9\ \mathrm{m\,s^{-2}}. \]

\textbf{Remark.} Even though $\ddot r = 0$ and $\ddot\theta = 0$ (both $r$ and $\theta$ vary ``uniformly''), the acceleration is \emph{not} zero: the rotation of the polar basis generates the centripetal term $-r\dot\theta^2 = -36\,\mathrm{m\,s^{-2}}$ and the Coriolis term $2\dot r\dot\theta = 12\,\mathrm{m\,s^{-2}}$.
\end{corrige}

\begin{corrige}[Exercise 8 — Car on a roundabout]
The car moves on a circle of radius $r = R = 20\,\mathrm{m}$ (constant) at constant speed $v = 8\,\mathrm{m\,s^{-1}}$.

\textbf{1. Angular velocity.} For circular motion, $v = R\dot{\theta}$, so
\[ \dot{\theta} = \frac{v}{R} = \frac{8}{20} = 0.4\ \mathrm{rad\,s^{-1}}. \]

\textbf{2. Centripetal acceleration.} Since $r$ is constant, $\dot{r} = \ddot{r} = 0$; since the speed is constant, $\ddot{\theta} = 0$. The acceleration reduces to
\[ a_r = -R\dot{\theta}^2 = -20 \times (0.4)^2 = -20 \times 0.16 = -3.2\ \mathrm{m\,s^{-2}}, \qquad a_\theta = 0. \]
The magnitude of the centripetal acceleration is
\[ a = R\dot{\theta}^2 = \frac{v^2}{R} = \frac{64}{20} = 3.2\ \mathrm{m\,s^{-2}}. \]

\textbf{3. Direction relative to the velocity.} The velocity is purely transverse: $\vec{v} = R\dot{\theta}\,\vec{e}_\theta$ (tangent to the circle). The acceleration is purely radial: $\vec{a} = -R\dot{\theta}^2\,\vec{e}_r$ (pointing towards the centre of the roundabout). Since $\vec{e}_r \perp \vec{e}_\theta$, the acceleration is \cle{perpendicular to the velocity} at every instant. This is why the speed stays constant: a force (here the friction of the tyres on the road) perpendicular to the velocity changes only the \emph{direction} of $\vec{v}$, never its magnitude.
\end{corrige}

\begin{corrige}[Exercise 9 — Bead on a rotating rod (GCE-type)]
Given: $\dot{\theta} = \omega$ (constant), so $\ddot{\theta} = 0$; and $r = r_0 e^{kt}$.

\textbf{Preliminary derivatives:}
\[ \dot{r} = k r_0 e^{kt} = kr, \qquad \ddot{r} = k^2 r_0 e^{kt} = k^2 r. \]

\textbf{1. Transverse acceleration.} Using $a_\theta = r\ddot{\theta} + 2\dot{r}\dot{\theta}$:
\[ a_\theta = r \times 0 + 2(kr)(\omega) = 2k\omega r. \qquad \blacksquare \]

\textbf{2. Radial acceleration.}
\[ a_r = \ddot{r} - r\dot{\theta}^2 = k^2 r - r\omega^2 = (k^2 - \omega^2)\,r. \]

\textbf{3. Magnitude of the total acceleration.}
\[ a = \sqrt{a_r^2 + a_\theta^2} = \sqrt{(k^2-\omega^2)^2 r^2 + 4k^2\omega^2 r^2} = r\sqrt{(k^2-\omega^2)^2 + 4k^2\omega^2}. \]
Expanding: $(k^2-\omega^2)^2 + 4k^2\omega^2 = k^4 - 2k^2\omega^2 + \omega^4 + 4k^2\omega^2 = k^4 + 2k^2\omega^2 + \omega^4 = (k^2+\omega^2)^2$. Hence the elegant result
\[ a = (k^2 + \omega^2)\,r = (k^2+\omega^2)\,r_0 e^{kt}. \]

\textbf{Indicative mark scheme (8 marks):}
\begin{center}
\begin{tabularx}{\linewidth}{|l|X|}
\hline
\rowcolor{popblueL} \textbf{Marks} & \textbf{Expected work} \\
\hline
M1 & Correct differentiation: $\dot r = kr$, $\ddot r = k^2 r$ \\
\hline
M1 A1 & Use of $a_\theta = r\ddot\theta + 2\dot r\dot\theta$ with $\ddot\theta = 0$; result $2k\omega r$ \\
\hline
M1 A1 & Use of $a_r = \ddot r - r\dot\theta^2$; result $(k^2-\omega^2)r$ \\
\hline
M1 & $a = \sqrt{a_r^2 + a_\theta^2}$ attempted \\
\hline
A1 & Simplification to $(k^2+\omega^2)r$ (recognising the perfect square) \\
\hline
A1 & Final answer fully simplified, in terms of the given quantities \\
\hline
\end{tabularx}
\end{center}

\textbf{Examiner's expectations:} quote the general polar formulas \emph{before} substituting; state explicitly that $\ddot\theta = 0$ because $\omega$ is constant; do not forget that the Coriolis term $2\dot r\dot\theta$ survives even though $\ddot\theta = 0$ — this is the most frequently lost mark.
\end{corrige}

\begin{corrige}[Exercise 10 — Particle on a circle with angular acceleration (GCE-type)]
Given: $r = R = 5\,\mathrm{m}$ constant, so $\dot{r} = \ddot{r} = 0$; $\ddot{\theta} = 0.4\,\mathrm{rad\,s^{-2}}$; at the instant considered, $\dot{\theta} = 2\,\mathrm{rad\,s^{-1}}$.

\textbf{1. Radial component.}
\[ a_r = \ddot{r} - r\dot{\theta}^2 = 0 - 5 \times (2)^2 = -20\ \mathrm{m\,s^{-2}}. \]
The negative sign confirms the inward (centripetal) direction; magnitude $20\,\mathrm{m\,s^{-2}}$.

\textbf{2. Transverse component.}
\[ a_\theta = r\ddot{\theta} + 2\dot{r}\dot{\theta} = 5 \times 0.4 + 0 = 2\ \mathrm{m\,s^{-2}}. \]

\textbf{3. Magnitude of total acceleration.}
\[ a = \sqrt{a_r^2 + a_\theta^2} = \sqrt{(-20)^2 + 2^2} = \sqrt{404} = 2\sqrt{101} \approx 20.1\ \mathrm{m\,s^{-2}}. \]

\textbf{Indicative mark scheme (6 marks):}
\begin{center}
\begin{tabularx}{\linewidth}{|l|X|}
\hline
\rowcolor{popblueL} \textbf{Marks} & \textbf{Expected work} \\
\hline
B1 & Stating $\dot r = \ddot r = 0$ (motion on a fixed circle) \\
\hline
M1 A1 & $a_r = -r\dot\theta^2$ correctly computed: $-20\ \mathrm{m\,s^{-2}}$ \\
\hline
M1 A1 & $a_\theta = r\ddot\theta$ correctly computed: $2\ \mathrm{m\,s^{-2}}$ \\
\hline
A1 & $a = \sqrt{404} \approx 20.1\ \mathrm{m\,s^{-2}}$ with correct units \\
\hline
\end{tabularx}
\end{center}

\textbf{Examiner's expectations:} candidates must realise that with constant angular acceleration the acceleration is \emph{not} purely centripetal: the transverse part $r\ddot\theta$ (tangential acceleration) is responsible for the increase of the speed. A common error is to answer $a = r\dot\theta^2 = 20\,\mathrm{m\,s^{-2}}$ for the total acceleration; this scores only part of the marks.
\end{corrige}

\begin{corrige}[Exercise 11 — Satellite dish tracking (GCE-type)]
Given: $\dot{\theta} = 0.1\,\mathrm{rad\,s^{-1}}$ constant, so $\ddot{\theta} = 0$; the part slides along the arm with $\dot{r} = 0.3\,\mathrm{m\,s^{-1}}$ constant, so $\ddot{r} = 0$; at the instant considered $r = 50\,\mathrm{m}$.

\textbf{1. Acceleration components.}
\begin{align*}
a_r &= \ddot{r} - r\dot{\theta}^2 = 0 - 50 \times (0.1)^2 = -0.5\ \mathrm{m\,s^{-2}},\\[2pt]
a_\theta &= r\ddot{\theta} + 2\dot{r}\dot{\theta} = 0 + 2 \times 0.3 \times 0.1 = 0.06\ \mathrm{m\,s^{-2}}.
\end{align*}

\textbf{2. Magnitude of the total acceleration.}
\[ a = \sqrt{a_r^2 + a_\theta^2} = \sqrt{(0.5)^2 + (0.06)^2} = \sqrt{0.25 + 0.0036} = \sqrt{0.2536} \approx 0.504\ \mathrm{m\,s^{-2}}. \]

\textbf{3. Physical origin of each term.}
\begin{itemize}
\item $-r\dot{\theta}^2 = -0.5\,\mathrm{m\,s^{-2}}$: the \cle{centripetal acceleration}. Even at constant $r$, the rotation of the dish permanently bends the velocity vector of the part towards the pivot; this term is provided by the mechanical constraint of the arm.
\item $2\dot{r}\dot{\theta} = 0.06\,\mathrm{m\,s^{-2}}$: the \cle{Coriolis (compound) acceleration}. It appears because the part slides radially \emph{while} the arm rotates: as $r$ increases, the transverse speed $r\dot\theta$ must increase, which requires a transverse acceleration. It is present only when $\dot r \neq 0$ \emph{and} $\dot\theta \neq 0$ simultaneously.
\item $\ddot{r} = 0$: no radial ``push'' acceleration, since the sliding speed is constant.
\item $r\ddot{\theta} = 0$: no tangential acceleration due to angular speeding-up, since the dish rotates uniformly.
\end{itemize}

\textbf{Indicative mark scheme (8 marks):}
\begin{center}
\begin{tabularx}{\linewidth}{|l|X|}
\hline
\rowcolor{popblueL} \textbf{Marks} & \textbf{Expected work} \\
\hline
B1 & Identifying $\ddot r = 0$ and $\ddot\theta = 0$ from the data \\
\hline
M1 A1 & $a_r = -r\dot\theta^2 = -0.5\ \mathrm{m\,s^{-2}}$ \\
\hline
M1 A1 & $a_\theta = 2\dot r\dot\theta = 0.06\ \mathrm{m\,s^{-2}}$ \\
\hline
M1 A1 & $a = \sqrt{0.2536} \approx 0.504\ \mathrm{m\,s^{-2}}$ \\
\hline
A1 & Correct physical interpretation: centripetal origin of $-r\dot\theta^2$ and Coriolis origin of $2\dot r\dot\theta$ \\
\hline
\end{tabularx}
\end{center}

\textbf{Examiner's expectations:} the explanation mark requires the candidate to name the two effects correctly and to link each to a mathematical term; vague statements such as ``the acceleration is due to the rotation'' do not earn the mark. Note the trap: although the target distance is described as constant, the \emph{part} slides along the arm, so $\dot r \neq 0$ for the part — read the question carefully to see which quantity each datum refers to.
\end{corrige}

\newpage

\coursec{Summary sheet}

\begin{popfigure}
\begin{tikzpicture}[mindmap, grow cyclic, every node/.style={concept, execute at begin node=\hskip0pt}, concept color=popblue!80, level 1/.style={level distance=4.5cm, sibling angle=360/7}]
\node[concept] {Motion in 2D}
child[concept color=popblue] {node[concept] {Cartesian}}
child[concept color=poporange] {node[concept] {Polar}}
child[concept color=popgreen] {node[concept] {Rotating Basis}}
child[concept color=poppurple] {node[concept] {Radial Velocity}}
child[concept color=poppink] {node[concept] {Radial Accel.}}
child[concept color=popgold] {node[concept] {Simple Cases}}
child[concept color=popmuted] {node[concept] {GCE Focus}};
\end{tikzpicture}
\legende{Mind map of Chapter FMA11: Motion of Particles in Two Dimensions}
\end{popfigure}

\begin{bilan}
\textbf{Key Definitions}
\begin{itemize}
\item Position vector in Cartesian coordinates: $\vec{r}=x\vec{i}+y\vec{j}$.
\item Velocity vector: $\vec{v}=\dot{x}\vec{i}+\dot{y}\vec{j}$; acceleration: $\vec{a}=\ddot{x}\vec{i}+\ddot{y}\vec{j}$.
\item Polar coordinates $(r,\theta)$ with unit vectors $\hat{r}$ (radial) and $\hat{\theta}$ (transverse), where $\hat{r}=\cos\theta\vec{i}+\sin\theta\vec{j}$, $\hat{\theta}=-\sin\theta\vec{i}+\cos\theta\vec{j}$.
\item Velocity in polar basis: $\vec{v}=\dot{r}\hat{r}+r\dot{\theta}\hat{\theta}$.
\item Acceleration in polar basis: $\vec{a}=(\ddot{r}-r\dot{\theta}^2)\hat{r}+(r\ddot{\theta}+2\dot{r}\dot{\theta})\hat{\theta}$.
\end{itemize}

\textbf{Fundamental Formulas and Theorems}
\begin{itemize}
\item Derivatives of rotating unit vectors: $\dot{\hat{r}}=\dot{\theta}\hat{\theta}$, $\dot{\hat{\theta}}=-\dot{\theta}\hat{r}$ (examinable proof).
\item Central force motion: transverse acceleration zero $\Rightarrow r^2\dot{\theta}=h$ (constant angular momentum per unit mass).
\item Circular motion ($r=R$ constant): $\vec{v}=R\dot{\theta}\hat{\theta}$, $\vec{a}=-R\dot{\theta}^2\hat{r}+R\ddot{\theta}\hat{\theta}$.
\item Motion with constant angular velocity $\dot{\theta}=\omega$: $\vec{a}=(\ddot{r}-r\omega^2)\hat{r}+2\dot{r}\omega\hat{\theta}$.
\item Radial motion only ($\dot{\theta}=0$): reduces to 1D motion along a line through origin.
\end{itemize}

\textbf{Standard Solution Methods}
\begin{itemize}
\item Convert given Cartesian equations to polar using $x=r\cos\theta$, $y=r\sin\theta$ and differentiate.
\item Express forces in radial/transverse components; apply Newton's second law component-wise.
\item For central forces, use $r^2\dot{\theta}=h$ to eliminate $\dot{\theta}$ and obtain a differential equation in $r$ (or $u=1/r$).
\item When path is given as $r=f(\theta)$, compute $\dot{r}=f'(\theta)\dot{\theta}$ and $\ddot{r}=f''(\theta)\dot{\theta}^2+f'(\theta)\ddot{\theta}$.
\item Always check dimensions and limiting cases (e.g., circular motion, straight line through origin).
\end{itemize}

\textbf{Common Pitfalls}
\begin{itemize}
\item Forgetting that $\hat{r}$ and $\hat{\theta}$ depend on $\theta$; their time derivatives are not zero.
\item Sign error in transverse acceleration: the Coriolis term $2\dot{r}\dot{\theta}$ is positive in the $\hat{\theta}$ direction.
\item Confusing $\ddot{r}$ (second derivative of radial distance) with the total radial acceleration component $\ddot{r}-r\dot{\theta}^2$.
\item Misapplying the chain rule when differentiating $r(\theta)$: $\dot{r}=r'\dot{\theta}$, $\ddot{r}=r''\dot{\theta}^2+r'\ddot{\theta}$.
\item Omitting the derivation of $\dot{\hat{r}}$ and $\dot{\hat{\theta}}$ in proof questions; GCE awards marks for each step.
\end{itemize}

\textbf{GCE Examination Tips}
\begin{itemize}
\item The derivation of radial-transverse components from Cartesian is a frequent 5-7 mark question; write it clearly with labelled steps.
\item Typical problem: "A particle moves along the curve $r=a(1+\cos\theta)$ with constant angular velocity $\omega$. Find the acceleration." Use the formulas directly.
\item For central force problems, remember $h=r^2\dot{\theta}$ is constant; substitute to get Binet's equation if needed.
\item Sketch the polar curve and indicate $\hat{r}$, $\hat{\theta}$ directions; this helps visualise component signs.
\item Always define your notation at the start: $\vec{r}$, $\hat{r}$, $\hat{\theta}$, $\dot{\theta}$, etc.
\item Check your final answer for correct units: acceleration in $\mathrm{m\,s^{-2}}$, velocity in $\mathrm{m\,s^{-1}}$.
\end{itemize}
\end{bilan}

\findecours

\end{document}
